\documentclass[11pt]{article}
\usepackage[margin=1in]{geometry}
\usepackage{amsmath,amssymb,amsthm,mathtools}
\usepackage{enumitem}
\usepackage{hyperref}
\usepackage{xcolor}
\usepackage{microtype}
\usepackage{booktabs}
\usepackage{array}

\hypersetup{colorlinks=true, linkcolor=blue!60!black, citecolor=blue!60!black, urlcolor=blue!60!black}

\newtheorem{theorem}{Theorem}[section]
\newtheorem{conjecture}[theorem]{Conjecture}
\newtheorem{lemma}[theorem]{Lemma}
\newtheorem{proposition}[theorem]{Proposition}
\newtheorem{corollary}[theorem]{Corollary}
\newtheorem{definition}[theorem]{Definition}
\newtheorem{observation}[theorem]{Observation}
\theoremstyle{remark}
\newtheorem{remark}[theorem]{Remark}
\newtheorem{convention}[theorem]{Convention}
\theoremstyle{plain}
\newtheorem{fact}[theorem]{Fact}
\theoremstyle{plain}

\newcommand{\one}{\mathbf 1}
\newcommand{\ip}[2]{\left\langle #1,#2\right\rangle}
\DeclareMathOperator{\Tr}{Tr}
\newcommand{\norm}[1]{\left\|#1\right\|}
\newcommand{\eps}{\varepsilon}
\newcommand{\pos}[1]{{(#1)_{+}}}
\newcommand{\wtP}{\widetilde{\mathcal P}}
\newcommand{\Pcal}{\mathcal P}
% macros used by the imported Zone-B section
\newcommand{\Kfrak}{\mathfrak K}
\newcommand{\Kstar}{K^{*}}
\newcommand{\kxi}{\kappa_\xi}
\newcommand{\kbar}{\bar\kappa}
\newcommand{\lbar}{\bar\ell}
\newcommand{\epsbar}{\bar\eps}
\newcommand{\RII}{[R2]}
% macros used by the imported Zone-C section
\newcommand{\eu}{\mathrm{e}}
\newcommand{\Lhat}{\widehat\Lambda}
\newcommand{\D}{\mathcal D}

\title{Towards a cleaner proof of the odd-cycle bound:\\
one identity, intrinsic thresholds,\\ and analytic replacements for the computer-assisted steps}
\author{Working note by Claude, for T.K., H.Y., and collaborators}
\date{July 17, 2026}

\begin{document}
\maketitle

\begin{abstract}
The odd-cycle bound $t(C_m,W)\ge p^m-p(1-p)^{m-1}$ (all graphons $W$ of edge
density $p$, all odd $m\ge3$) is proved in the consolidated notes
[R1]~=~\texttt{paper\_new.tex} (region $p\ge\tfrac23$, and the per-cycle
theorems $C_3$--$C_{13}$) and [R2]~=~\texttt{paper\_region2\_v2.tex} (region
$\tfrac12<p<\tfrac23$, odd $m\ge15$).  The combined proof has two
unsatisfying features: a many-layered case split ($p\ge\tfrac23$ versus
$\tfrac12<p<\tfrac23$; within the second region, $m\le13$ versus $m\ge15$;
within the first, residual pairs $m\le61$ versus $m\ge63$), and five
computer-assisted certificate steps.  This note does three things.
(1)~It shows that the density split at $p=\tfrac23$ is \emph{intrinsic}
rather than ad hoc: $q=1-p=\tfrac13$ is exactly the density below which the
near-extremal multipartite family carries several spectral modes above the
frontier line $\lambda=q$, and exactly the density at which the safe-mass
separation $L<q$ fails; in the other direction, the pointwise mechanism of
the high-density proof provably cannot cross $q\approx0.4118$.  Any proof
along the present lines must switch mechanisms inside $(\tfrac13,0.4118)$,
and $\tfrac13$ is the canonical switch point.
(2)~It reorganises the proof around a single spectral-shift identity with
one budget and one frontier dichotomy, so the split becomes a principled
branching.
(3)~It removes case-split and computer-assisted content from the chain:
the Region-II scalar route is extended from $m\ge15$ to \emph{all} odd
$m\ge3$ (\S\ref{sec:smallm}), which deletes the per-cycle theorems
$C_5$--$C_{13}$ from the main chain; the Zone-B $23$-box certificate
(\S\ref{sec:zoneB}) and the Zone-C $3001$-box certificate
(\S\ref{sec:zoneC}) are replaced by fully analytic proofs; the
Appendix-B constant inequalities of [R1] are proved analytically for all
real parameters (\S\ref{sec:appconst}), eliminating the scripted sweep;
and the last and largest certificate, $\operatorname{StripCert}(61)$, is
eliminated by a unified strip theorem (\S\ref{sec:strip}) covering
$r\ge1$ and all odd $m$ at one stroke --- subsuming also the separate
$r=1$ theorem, the $m\ge63$ analytic-strip section, and the Appendix-B
constants of [R1].
Every new proof was validated numerically (dense grids plus exact rational
arithmetic at extremes) and audited by an independent verifier and an
adversarial reviewer; the per-section provenance notes record this.
\end{abstract}

\tableofcontents

\section{Introduction}\label{sec:intro}

Let $W:[0,1]^2\to[0,1]$ be a graphon with edge density $p=t(K_2,W)$ and let
$m\ge3$ be odd.  The statement under discussion is
\begin{equation}\label{eq:conj}
        t(C_m,W)\ \ge\ g_m(p):=p^m-p(1-p)^{m-1},
\end{equation}
sharp at every balanced complete multipartite graphon (equality at
$p=1-\tfrac1k$, Observation~\ref{obs:equality}).  Write $q=1-p$ and
$U=1-W$.  The bound is trivial for $p\le\tfrac12$.  The two consolidated
notes prove it everywhere else:
\begin{itemize}[leftmargin=2em]
\item \textbf{Region I} ($p\ge\tfrac23$, i.e.\ $q\le\tfrac13$): [R1], via
the two-sided spectral-shift identity, one edge deletion, the exact kernel
expansion, and diagonal positivity of the kernels $\wtP_{m,r}(q,\ell)$.
\item \textbf{Region II} ($\tfrac12<p<\tfrac23$, i.e.\
$\tfrac13<q<\tfrac12$): [R2], via the one-sided identity, the one-frontier
reduction, the exact Huber shape elimination, and the scalar inequality
$R_m\le C_m\psi(\xi,\rho)$, proved for $m\ge15$; odd $m\le13$ is quoted
from the per-cycle theorems $C_5$--$C_{13}$ of [R1].
\end{itemize}
Throughout, [R1] and [R2] refer to \texttt{paper\_new.tex} and
\texttt{paper\_region2\_v2.tex} (both in this directory), and results of
theirs are cited by their \LaTeX{} labels, set in italics, e.g.\
Lemma~\emph{threegeo} of [R2].

Two features of the combined proof feel unsatisfactory; this note
addresses both.

\paragraph{The case split.}  The density split at $p=\tfrac23$, the length
split at $m=15$ inside Region II, and the residual-strip split at $m=63$
inside Region I all look like artifacts.  Section~\ref{sec:dichotomy}
shows that the first split is \emph{not} an artifact: it sits at an
intrinsic structural threshold of the problem, and each of the two
mechanisms genuinely dies on the far side of its region
(Propositions~\ref{prop:frontier-count}, \ref{prop:safe-mass} and
Remark~\ref{fact:ceiling}).  What \emph{can} be removed is the other two
splits: the length split at $m=15$ disappears in \S\ref{sec:smallm}, and
the residual-strip split is addressed in \S\ref{sec:strip}.  What remains
is one principled branching: a single identity and budget
(Section~\ref{sec:architecture}), a dichotomy on the frontier spectrum,
and two payment schemes whose validity regions overlap on
$(\tfrac13,0.41)$ and together tile $(0,\tfrac12)$.

\paragraph{The computer-assisted steps.}  After the July-12 state, the
computer-assisted content of the full proof was exactly:
\begin{enumerate}[label=\textup{(C\arabic*)}]
\item\label{c:strip} $\operatorname{StripCert}(61)$: exact rational
interval verification of the two-sided surplus criterion for all residual
diagonal pairs ($r\ge2$, $n=m-2r>2r$) with odd $m\le61$, over
$(q,\ell)\in[0,\tfrac13]\times[0,\tfrac12]$ ([R1], Prop.~\emph{M61}).
\item\label{c:appendix} The Appendix-B constant sweep of [R1]: scripted
verification of $\frac{99}{100m}P(\theta)B_{0,1}(\theta)^m\ge1$ for the
integer pairs with odd $63\le m\le499$, plus twelve scripted integer
comparisons.
\item\label{c:zoneB} \texttt{zoneB\_certifier.py}: $23$ exact rational
boxes for the Zone-B battle ([R2], Lemma~\emph{zoneB-battle}).
\item\label{c:zoneC} \texttt{zoneC\_certifier.py}: $3001$ exact rational
boxes (the prose count ``2997'' in [R2] is stale), subdivision depth
$28$, cycle lengths up to $m=1716$, for the moderate-$e$ part of Zone C
([R2], Lemma~\emph{zoneC-mod}).
\item\label{c:c513} Indirectly: the per-cycle theorems $C_5$--$C_{13}$
(needed only for odd $m\le13$ inside Region II) contain their own
computer-assisted steps (the $C_9$ Gram-matrix minors, the $C_{13}$
frontier split).
\end{enumerate}
The state after this note: \ref{c:strip}, \ref{c:appendix},
\ref{c:zoneB} and \ref{c:zoneC} are \emph{eliminated} (fully analytic
replacements: \S\ref{sec:strip}, \S\ref{sec:appconst}, \S\ref{sec:zoneB},
\S\ref{sec:zoneC}; the strip theorem's only finite residue is a $21$-row
table of one-line exact rational inequalities in the sanctioned displayed
style), and \ref{c:c513} is removed from the main chain by the small-$m$
extension (\S\ref{sec:smallm}), whose own residual certificate content is
far smaller and structurally uniform.  Section~\ref{sec:map} tabulates
exactly what remains.

\paragraph{Numerical ground truth.}  Two calibration computations inform
everything below.  First, the Region-II scalar inequality
$R_m\le C_m\psi(\xi,\rho)$ holds on its admissible domain for \emph{all}
odd $3\le m\le13$, with healthy margins: the minimum relative margin is
$0.100$ for $m=3$ and $\tfrac16$ for $m=5$ (both attained in the limit at
the Tur\'an corner, with closed-form corner ratios
$c_m'=3m(m-2)/((2\sqrt{2/3}-1)(m2^{m-1}-m+6))<1$), and $0.37/0.29/0.24/0.20$
for $m=7/9/11/13$ (attained at the pinch $\kappa=\kappa_{\max}$,
$e=\Theta(1/m)$).  (A first float-grid scan reported a $2\cdot10^{-4}$
margin near the corner; this was a \texttt{float64} cancellation artifact
in $R_m\sim\delta^2$, resolved by high-precision and exact rational
evaluation.)  So the restriction $m\ge15$ in [R2] is an artifact of the
zone constants, not of the reduction --- and \S\ref{sec:smallm} removes it.
Second, for $q\le\tfrac13$ the frontier route breaks structurally
(Section~\ref{sec:dichotomy}): no literal unification of the two regions
is available along present lines, and the correct target is a principled
branching, not a single mechanism.

\section{One identity, one budget, one dichotomy}\label{sec:architecture}

Everything in both regions rests on the same three objects.  Decompose
$L^2[0,1]=\mathbb R\one\oplus\one^\perp$ and write, for the complement
kernel $U=1-W$,
\[
        T_U=\begin{pmatrix}q&g^*\\ g&A\end{pmatrix},
        \qquad g=T_U\one-q\one\in\one^\perp,
        \qquad A=P_{\one^\perp}T_UP_{\one^\perp}.
\]

\begin{lemma}[Spectral-shift identity; {[R2]}, Prop.~\emph{shift}]\label{lem:shift}
For every graphon $W$ and every odd $m\ge3$,
\[
        t(C_m,W)=p^m-\Tr(A^m)+m[z^m]\mathcal L_W(z),
        \qquad
        \mathcal L_W(z)=-\log\Bigl(1-\frac{z^2\ip g{(I+zA)^{-1}g}}{1-pz}\Bigr),
\]
and the shift series $\mathcal L_W$ has nonnegative coefficients.  Adding
the same identity for $U$ gives the two-sided form
$t(C_m,W)+t(C_m,U)=p^m+q^m+S_m$ of [R1], Thm.~\emph{two-sided}.
\end{lemma}

\begin{lemma}[Hilbert--Schmidt budget; {[R2]}, eq.~\emph{HS-budget}]\label{lem:budget}
$\Tr(A^2)+2\norm g_2^2\le pq$, and $\norm A\le\tfrac12$.
\end{lemma}

\begin{lemma}[No-frontier case, all densities; {[R2]}, Lem.~\emph{no-frontier}]\label{lem:no-frontier}
If $\lambda_{\max}(A)\le q$, then \eqref{eq:conj} holds --- for every
$q\in(0,\tfrac12]$ and every odd $m\ge3$.
\end{lemma}

\begin{proof}[Proof sketch]
Negative eigenvalues contribute $\nu^m<0\le q^{m-2}\nu^2$ to $\Tr(A^m)$;
eigenvalues in $[0,q]$ contribute $\nu^m\le q^{m-2}\nu^2$; summing against
the budget gives $\Tr(A^m)\le pq^{m-1}$, and the shift term is
nonnegative.
\end{proof}

Call an eigenvalue of $A$ \emph{frontier} if it exceeds $q$.  By
Lemma~\ref{lem:no-frontier} the whole conjecture reduces to graphons whose
complement compression has at least one frontier eigenvalue, and the
entire remaining content of both regional proofs is a \emph{payment
scheme}: the shift term (and, in Region I, the deletion gap) must pay for
the frontier excess $\Tr(A^m)-pq^{m-1}$.  The two regions differ only in
which payment scheme is available, and the next section shows that this
difference is forced by the geometry of the near-extremal family.

\section{The intrinsic thresholds}\label{sec:dichotomy}

\begin{observation}[Equality family]\label{obs:equality}
For $p=1-\tfrac1k$ ($k\ge2$) and $W$ the balanced complete $k$-partite
graphon, $T_W$ has eigenvalue $p$ on $\one$ and $-\tfrac1k$ with
multiplicity $k-1$ on $\one^\perp$, so for odd $m$
\[
        t(C_m,W)=p^m-(k-1)k^{-m}=p^m-p(1-p)^{m-1}:
\]
equality in \eqref{eq:conj}.  Dually, $U$ is the disjoint union of $k$
blocks of measure $\tfrac1k$; here $g=0$ and $A$ has the eigenvalue
$q=\tfrac1k$ with multiplicity $k-1$ --- $k-1$ modes sitting exactly
\emph{on} the frontier line $\lambda=q$.
\end{observation}

The next proposition is elementary but organises everything: the number
of frontier modes is controlled by $p/q$, and this control degenerates
exactly at $q=\tfrac13$.

\begin{proposition}[Frontier count]\label{prop:frontier-count}
For every graphon, the number $N_f$ of eigenvalues of $A$ exceeding $q$
satisfies
\[
        N_f\ <\ \frac{\Tr(A^2)}{q^2}\ \le\ \frac{pq}{q^2}=\frac pq .
\]
In particular $N_f\le1$ for $q>\tfrac13$ (the one-frontier lemma of
[R2]).  Conversely, for every $q<\tfrac13$ there are graphons with
$N_f\ge2$ lying on the near-extremal family: if $U$ is the disjoint union
of three blocks of measure $\beta$, then $q=3\beta^2$ and $A$ has the
eigenvalue $\beta$ with multiplicity two, which is a frontier eigenvalue
iff $\beta<\tfrac13$, i.e.\ iff $q<\tfrac13$.  More generally $k$ blocks
of measure $\beta<\tfrac1k$ give $N_f=k-1$ at $q=k\beta^2<\tfrac1k$.
\end{proposition}

\begin{proof}
Each frontier eigenvalue contributes more than $q^2$ to $\Tr(A^2)$, and
the budget (Lemma~\ref{lem:budget}) caps $\Tr(A^2)$ by $pq$.  ($N_f=2$
requires $2q^2<pq$, i.e.\ $q<\tfrac13$; this is the contrapositive of the
one-frontier lemma.)  For the block family: on $\one^\perp$, differences
of block indicators are eigenfunctions of $T_U$ with eigenvalue $\beta$,
orthogonal to $\one$; there are $k-1$ independent ones, and they survive
the compression unchanged.  $q=t(K_2,U)=k\beta^2$, and $\beta>k\beta^2$
iff $\beta<\tfrac1k$.
\end{proof}

\begin{proposition}[Safe-mass separation]\label{prop:safe-mass}
Suppose $A$ has a frontier eigenvalue $\alpha>q$, and let
$L^2=pq-\alpha^2$ bound the square mass of the remaining spectrum
(Lemma~\ref{lem:budget}).  Then
\[
        q>\tfrac13\ \Longrightarrow\ L<q,
        \qquad\text{while}\qquad
        q<\tfrac13\ \Longrightarrow\ L>q
        \ \text{ for all }\alpha\ \text{close enough to }q .
\]
\end{proposition}

\begin{proof}
$L<q\iff q(1-2q)<\alpha^2$.  If $q>\tfrac13$ then
$\alpha^2>q^2\ge q(1-2q)$.  If $q<\tfrac13$ then at $\alpha=q$ one has
$q^2<q(1-2q)$, so by continuity $L>q$ for $\alpha$ near $q$.
\end{proof}

The master-defect ledger of [R2] needs both facts: one frontier mode (so
a single scalar $\alpha$ carries all the excess) and $L<q<\alpha$ (so the
safe spectrum is strictly separated below the frontier and $f=\alpha-L>0$
prices the safe channel).  By Propositions~\ref{prop:frontier-count}
and~\ref{prop:safe-mass} both fail for $q<\tfrac13$, and they fail
\emph{on the near-extremal multipartite family itself}, not on exotic
configurations: as $\beta\uparrow\tfrac1k$ the block family converges to
the equality case of Observation~\ref{obs:equality} while carrying $k-1$
frontier modes.  A one-frontier payment scheme below $q=\tfrac13$ would
have to pay $k-1$ excesses with slack that vanishes on the equality
family --- there is no room.  This is the structural reason Region I must
use an \emph{aggregate} mechanism: the kernel expansion of [R1]
integrates over the whole spectral measure at once.

In the other direction, the Region-I mechanism has its own hard ceiling.

\begin{remark}[Ceiling of the aggregate route; {[R1]}, Rmk.~\emph{thresholds}]\label{fact:ceiling}
The aggregate route bounds $t(C_m,U)$ by the path density
$t(P_{m-1},U)$ (one edge deletion) and then requires pointwise positivity
of the diagonal kernels $\wtP_{m,r}(q,\ell)$.  The exact positivity
threshold satisfies $\inf_{m,r}q^*(m,r)=0.411756\ldots$ (attained as
$r/m\to0.17656$), so the route proves densities $q\le0.41$ at best and
\emph{cannot} reach $q=\tfrac12$: beyond the threshold the kernels
themselves go negative --- the deletion discards too much.  (For
$q\le\tfrac13$, by contrast, every threshold is at distance $\ge0.078$:
the fat-margin regime that the analytic replacements of
\S\ref{sec:strip} exploit.)
\end{remark}

\begin{remark}[Why $q=\tfrac13$ is the canonical switch point]\label{rmk:canonical}
The two mechanisms are simultaneously valid on a genuine overlap window:
numerically $(\tfrac13,0.4117)$ for the aggregate route (proved to
$\tfrac13$) and $(\tfrac13,\tfrac12)$ for the frontier route.  The split
could be placed anywhere in the overlap, but $\tfrac13$ is distinguished
three times over: it is where the one-frontier structure theorem begins
(Proposition~\ref{prop:frontier-count}), where the safe-mass separation
begins (Proposition~\ref{prop:safe-mass}), and where the multipartite
equality family ends ($p=1-\tfrac1k\ge\tfrac23$ for $k\ge3$; inside
Region II the bound is strict, with equality approached only at the two
corners $q\uparrow\tfrac12$ and $q\downarrow\tfrac13$ --- precisely the
pinch and Tur\'an corners where the Region-II proof is tight).  The case
split is thus a feature of the extremal landscape: Region I contains
infinitely many equality structures and needs a mechanism that is exact
on all of them (the kernel expansion is: at every $q=\tfrac1k$ the
deletion gap and the shift vanish together on the equality family);
Region II contains none, but its two corners pin the inequality from both
ends.  A proof with no split at all would need a payment scheme that is
simultaneously exact on infinitely many multipartite equality structures
and strictly positive on the one-frontier configurations --- nothing of
the sort is currently in sight, and the numerics above say the two
existing schemes cannot individually be stretched to cover $(0,\tfrac12)$.
\end{remark}

\section{Region II for all odd $m\ge3$: removing the per-cycle theorems}\label{sec:smallm}

\subsection*{Provenance and verification status}
Proved by Claude on 2026-07-17 (prover--verifier--skeptic pipeline; one
repair round fixing five display-level constant errors flagged by the
adversarial reviewer, whose final verdict was \emph{sound}; independent
verifier: \emph{pass}, $124$ checks, $0$ failures, including independent
re-derivations of the two certificate steps).  The result
extends the Region-II scalar inequality to all odd $3\le m\le13$ through
the single relaxed $\lambda=1$ certificate $(T_m)$; together with the
$m$-uniformity of the [R2] reduction (Remark~\emph{muniform} below) this
deletes the per-cycle theorems $C_5$--$C_{13}$ from the main chain.  Two
steps remain certificate-based (exact rational, uniform protocol): the
two-variable Bernstein tables of Prop.~\emph{band} and the per-$m$
subdivision of Prop.~\emph{moderate}; everything else is prose.  The
content of this section is the note
\texttt{claude\_jul17/smallm\_regionII\_proof\_v2.tex}, inlined verbatim
(sectioning demoted, labels prefixed where needed).

\subsection{Statement, notation and result status}\label{sec:smallm-intro}

This note extends the Region-II scalar inequality of
\texttt{paper\_region2\_v2.tex} (cited below as [R2]) from odd $m\ge15$ to all
odd $3\le m\le 13$, so that the per-cycle theorems for $C_5,\dots,C_{13}$ can
be removed from the main chain.  All notation is from [R2]:
$q\in(\tfrac13,\tfrac12)$, $p=1-q$, $\alpha\in(q,r(q)]$ the frontier
eigenvalue, $L=\sqrt{pq-\alpha^2}$, $d=\alpha-q$, $e=1-2\alpha$,
$f=\alpha-L$, $\kappa=d/e$, $x=\alpha/p$, $\tau=q/\alpha$, $y=L/p$,
$\ell=L/\alpha$, $u=1-x$, $\delta=\alpha-\tfrac13$,
\[
 k_m(\lambda)=\frac{p^{m-1}-\lambda^{m-1}}{p+\lambda},\quad
 A_m=2L^{m-2}+mk_m(\alpha),\quad
 B_m=2L^{m-2}+mk_m(L),\quad
 R_m=\alpha^m+L^m-pq^{m-1},
\]
and $C_m$, $\xi$, $\rho$, $\psi(\xi,\rho)$ as in [R2, Def.~huber].  We use the
$(e,\kappa)$ chart of [R2, Lem.~chart]: the admissible $(q,\alpha)$ are
exactly $0<e<\tfrac13$, $0<\kappa\le\kappa_{\max}(e)=\frac{1-e}{1+e}$,
$\kappa<\kappa_q(e)=\frac{1-3e}{6e}$, and on it
$L^2=\alpha e-d(d+e)$, $\xi=(1-e)^2\kappa/e$, $\delta=\frac{1-3e}{6}$.

\begin{theorem}\label{thm:smallm-main}
For every odd $m$ with $3\le m\le13$ and every admissible $(q,\alpha)$,
\begin{equation}\label{eq:smallm-target}
 R_m\ \le\ C_m\,\psi(\xi,\rho).
\end{equation}
Consequently (Remark~\ref{rem:muniform}), the odd-cycle bound
$t(C_m,W)\ge p^m-pq^{m-1}$ holds in Region II for these $m$.
\end{theorem}

\begin{remark}[$m$-uniformity of the reduction in {[R2]}]\label{rem:muniform}
The statement of [R2, Thm.~huber] carries the standing hypothesis
``odd $m\ge15$'' of that paper's domain $D$, but its proof --- and the proofs
of the whole reduction chain it rests on, namely
[R2, Prop.~shift, Lem.~no-frontier, Lem.~one-frontier,
Lem.~forced-variance, Prop.~master-defect, Prop.~psi-forms] --- is
$m$-pointwise: $m$ enters only through the quantities $A_m$, $B_m$, $R_m$
defined above, and the only structural inputs are
$0<A_m<B_m$ [R2, Rem.~knegative] and the frontier geometry $L<q<\alpha<p$,
which hold for every odd $m\ge3$ on the admissible domain.
($A_m<B_m$ because $k_m$ is strictly decreasing on $[0,p]$ and $L<\alpha$;
$A_m>0$ because $\alpha<p$ gives $k_m(\alpha)>0$.)
Hence the implication
``\eqref{eq:smallm-target} on the admissible domain $\Rightarrow$
$t(C_m,W)\ge p^m-pq^{m-1}$ in Region II'' is valid verbatim for every odd
$m\ge3$, and Theorem~\ref{thm:smallm-main} may legitimately invoke it for
$3\le m\le13$.
\end{remark}

\emph{Status of the proof.}  Steps proved analytically below: the certificate
relaxation (Lemma~\ref{lem:relax}), the positivity gate
(Lemma~\ref{lem:gate}), the small-$e$ range $e\le1/60$
(Proposition~\ref{prop:smalle}), the tail cut on the Tur\'an band
(Lemma~\ref{lem:tail}: analytic reduction to six one-line exact rational
inequalities, displayed in Table~\ref{tab:tailints}), and the reduction of
the Tur\'an band $\delta\le\Delta_m$ to two explicit polynomial
nonnegativity claims (Proposition~\ref{prop:band}), which are then certified
by exact Bernstein coefficient tables.  The remaining moderate range
($1/60\le e\le\tfrac13-2\Delta_m$) is covered by an exact-rational certified
subdivision (Proposition~\ref{prop:moderate}), in the same certificate
protocol as [R2]; replacing these two certificate steps by prose arguments
is left open.

Throughout, a terminating decimal denotes that exact rational number, and
every decimal comparison is rounded in the safe direction with the
underlying exact comparison stated; all displayed inequalities are
additionally machine-verified by \texttt{validate\_smallm\_v2.py}.

\subsection{The relaxed certificate and the positivity gate}

\begin{lemma}[Relaxed $\lambda=1$ certificate]\label{lem:relax}
On the admissible domain, with $w:=\sqrt{2\alpha}$,
\begin{equation}\label{eq:rho-id}
 \rho=\frac{A_m w}{4B_m f},
\end{equation}
and consequently
\begin{equation}\label{eq:Tm}
 C_m\psi(\xi,\rho)\ \ge\ \sqrt{2\alpha}\,B_mf\Bigl(d-\frac{e^2}{16\alpha^2(1+\rho)}\Bigr)
 \ \ge\ \sqrt{2\alpha}\,B_mfd-\frac{e^2B_m^2f^2}{4\alpha^2A_m}.
\end{equation}
Hence \eqref{eq:smallm-target} follows from
\begin{equation}\label{eq:Tm2}
 (T_m):\qquad 4\alpha^2A_m R_m\ \le\ 4\alpha^2 A_m\sqrt{2\alpha}\,B_mfd-e^2B_m^2f^2 .
\end{equation}
\end{lemma}

\begin{proof}
By [R2, Def.~huber], $\rho=\frac{A_m}{B_m}\frac{\sqrt\alpha}{2\sqrt2 f}$, and
$\frac{\sqrt\alpha}{2\sqrt2}=\frac{\sqrt{2\alpha}}4$, giving \eqref{eq:rho-id}.
The first inequality in \eqref{eq:Tm} is the dual certificate $\lambda=1$
[R2, eq.~cert-II].  For the second, $1+\rho\ge\rho$ and \eqref{eq:rho-id} give
\[
 \sqrt{2\alpha}\,B_mf\cdot\frac{e^2}{16\alpha^2(1+\rho)}
 \le \sqrt{2\alpha}\,B_mf\cdot\frac{e^2}{16\alpha^2}\cdot\frac{4B_mf}{A_m\sqrt{2\alpha}}
 =\frac{e^2B_m^2f^2}{4\alpha^2A_m}. \qedhere
\]
\end{proof}

\begin{lemma}[Positivity gate; all odd $m\ge3$]\label{lem:gate}
If $R_m>0$ then
\begin{equation}\label{eq:gate}
 (m-1)\kappa\ >\ x\bigl(1+\kappa-\ell^{m-2}\bigr).
\end{equation}
\end{lemma}

\begin{proof}
By [R2, eq.~defect-form], $R_m=\alpha^{m-1}(\alpha-p\tau^{m-1})+L^m$, so
$R_m>0$ means $\tau^{m-1}<x(1+\ell^m)$, since
$L^m/(p\alpha^{m-1})=x\ell^m$.  Bernoulli's inequality gives
$\tau^{m-1}=(1-d/\alpha)^{m-1}\ge1-(m-1)d/\alpha$, whence
\[
 (m-1)\frac d\alpha\ >\ 1-x-x\ell^m .
\]
Multiply by $\alpha/e$ and use $1-x=u=\frac{d+e}p=\frac{x(d+e)}\alpha$, so
that $(1-x)\alpha/e=x(1+\kappa)$; finally
$\ell^m\alpha/e\le\ell^{m-2}$ because $\ell^2\le e/\alpha$
[R2, Lem.~elementary].
\end{proof}

\begin{corollary}[Gate on the small-$e$ range]\label{cor:gate-smalle}
Let $3\le m\le13$ be odd, $e\le1/60$, and $R_m>0$.  Then, with the standing
bounds $x\ge\frac{1-e}{1+3e}\ge\frac{59}{63}>0.9365$ and
$\ell\le\sqrt{2e/(1-e)}\le1.4262\sqrt e\le0.1842$,
\begin{equation}\label{eq:gate-smalle}
 \kappa\ >\ \frac{x(1-\ell)}{m-1-x}\ \ge\ \frac{0.763}{m-1}\ \ge\ 0.0635,
 \qquad
 \xi=\frac{(1-e)^2\kappa}e\ \ge\ 3.68,
\end{equation}
and for $m=3$ moreover $\kappa>0.7183$.  In particular the region
$\{\,e\le1/60,\ \xi\le1,\ R_m>0\,\}$ is empty for $3\le m\le13$, and
\begin{equation}\label{eq:eps-smalle}
 \eps:=\frac{e^2}{16\alpha^2(1+\rho)d}=\frac{e}{4(1-e)^2(1+\rho)\kappa}
 \ \le\ \min\Bigl(0.068,\ \frac{e}{4(1-e)^2\kappa}\Bigr).
\end{equation}
\end{corollary}

\begin{proof}
Since $m\ge3$ and $\ell<1$, $\ell^{m-2}\le\ell$, so \eqref{eq:gate} gives
$\kappa(m-1-x)>x(1-\ell)$, i.e.\ the first bound; then
$x(1-\ell)\ge0.9365\cdot0.8158>0.763$ and $m-1-x<m-1$.  For $m=3$,
$t\mapsto t/(2-t)$ is increasing, so
$\kappa>0.9365\cdot0.8158/(2-0.9365)>0.7183$.  Next,
$\xi\ge(59/60)^2\cdot\frac{0.763}{12}\cdot60>3.68$ using $m\le13$; in
particular $\xi>1$.  Finally $\eps=1/(4\xi(1+\rho))\le1/(4\cdot3.68)<0.068$
and $\rho\ge0$; the equality in \eqref{eq:eps-smalle} is the chart identity
$\xi=(1-e)^2\kappa/e$.
\end{proof}
(The displayed decimal replacements: $0.9365\cdot0.8158=0.76399\ldots>0.763$,
$0.9365\cdot0.8158/1.0635=0.71837\ldots>0.7183$,
$(59/60)^2\cdot0.763\cdot5=3.68889\ldots>3.68$, and
$1/14.72=0.06793\ldots<0.068$.)

\subsection{The small-$e$ range $e\le1/60$}

\begin{proposition}\label{prop:smalle}
Let $3\le m\le13$ be odd, $e\le1/60$, $R_m>0$.  Then
$R_m\le\sqrt{2\alpha}B_mf\bigl(d-\frac{e^2}{16\alpha^2(1+\rho)}\bigr)\le C_m\psi$.
\end{proposition}

\begin{lemma}[Reduction; m-adapted {[R2, Lem.~reduceA]}]\label{lem:reduce}
Under the hypotheses of Proposition~\ref{prop:smalle}, the first inequality
follows from
\begin{equation}\label{eq:red}
 \sqrt{1-e}\,\frac{(1-\ell)(1-y^{m-1})(1-\eps)}{1+y}
 \ \ge\
 x^{m-2}\Bigl[\frac{m-1}{mx}-\frac{1+1/\kappa}{m}\Bigr]+\Lambda,
 \qquad
 \Lambda:=\frac{x^{m-2}(e/\alpha)^{m/2-1}}{m\kappa}.
\end{equation}
\end{lemma}

\begin{proof}
Exactly as in [R2, Lem.~reduceA]: $B_m\ge mk_m(L)=\frac{mp^{m-2}(1-y^{m-1})}{1+y}$,
$\sqrt{2\alpha}=\sqrt{1-e}$, $f=\alpha(1-\ell)$ give: payment
$\ge md\alpha p^{m-2}\cdot\mathrm{LHS}\eqref{eq:red}$; Bernoulli
$\tau^{m-1}\ge1-(m-1)d/\alpha$ and $L^m\le(\alpha e)^{m/2}$ give
$R_m\le md\alpha p^{m-2}\cdot\mathrm{RHS}\eqref{eq:red}$.
(No $m\ge15$ constant enters.)
\end{proof}

\begin{proof}[Proof of Proposition~\ref{prop:smalle}]
Standing bounds for $e\le1/60$, each stated with its exact justification:
$x\ge\frac{1-e}{1+3e}\ge\frac{59}{63}>0.9365$ (from $d\le e$);
$\ell^2\le e/\alpha=\frac{2e}{1-e}\le\frac{120}{59}e\le2.04e$, hence
$\ell\le1.4262\sqrt e\le0.1842$
($1.4262^2=2.03404\ldots>120/59=2.03389\ldots$;
$1.4262/\sqrt{60}<0.1842$); $y\le\ell$ (as $p\ge\alpha$);
$\sqrt{1-e}\ge1-0.51e$ (as $(1-0.51e)^2-(1-e)=-0.02e+0.2601e^2\le0$ for
$e\le0.02/0.2601$); $1/x=1+(1+\kappa)e/\alpha\le1+2.04(1+\kappa)e$
(as $1/\alpha\le120/59\le2.04$);
$u=1-x=\frac{(1+\kappa)e}{p}\ge1.9047(1+\kappa)e$
(as $p=\frac{1+e}2+d\le\frac{1+3e}2\le\frac{21}{40}$ and
$40/21=1.90476\ldots\ge1.9047$); and Corollary~\ref{cor:gate-smalle}.

\emph{Case $m=3$.}  Using
$\frac{(1-\ell)(1-y^2)(1-\eps)}{1+y}\ge(1-\ell)(1-y)(1-y^2)(1-\eps)$ and
$(1-a_1)\cdots(1-a_k)\ge1-\sum a_i$ on $[0,1]$,
\[
 \mathrm{LHS}\eqref{eq:red}\ \ge\ 1-2\ell-y^2-\eps-0.51e
 \ \ge\ 1-0.3684-0.034-0.068-0.0085\ >\ 0.52,
\]
where $y^2\le\ell^2\le2.04e\le2.04/60=0.034$.  On the other side, for $m=3$
the bracket in \eqref{eq:red} is $\frac{2}{3x}-\frac{1+1/\kappa}3$, so
$\mathrm{RHS}\eqref{eq:red}=\frac23-\frac{x(1+1/\kappa)}3+\Lambda$ with
$\Lambda=\frac{x\sqrt{e/\alpha}}{3\kappa}$.  Using $\kappa>0.7183$,
$1+1/\kappa>2$, $x\ge0.9365$, $x\le1$, $\sqrt{e/\alpha}\le1.4262\sqrt e$
and $1.4262/(3\cdot0.7183)=0.66183\ldots\le0.662$:
\[
 \mathrm{RHS}\eqref{eq:red}\le\frac23-\frac{0.9365\cdot2}{3}+0.662\sqrt e
 \le\frac23-0.6243+0.662\cdot0.1291<0.128<0.52,
\]
($\sqrt{1/60}<0.1291$ as $0.1291^2=0.01666\,681>1/60$).  So \eqref{eq:red}
holds for $m=3$.

\emph{Case $5\le m\le13$.}  Put $z:=(m-2)u$, $c_m:=\frac{m-3}{2(m-2)}$,
$\beta_m:=\frac{1.0711}{m-2}$ (so that $2.04(1+\kappa)e\le\beta_m z$, since
$\beta_mz\ge1.0711\cdot1.9047(1+\kappa)e=2.04012\ldots(1+\kappa)e$), and
\[
 \Gamma_0(e):=2.8524\sqrt e+0.51e+(2.04e)^2,\qquad
 \sigma_m(e):=1.21\,(2.04e)^{(m-2)/2}.
\]
Then $2\ell+0.51e+y^{m-1}\le\Gamma_0$ (using
$y^{m-1}\le\ell^4\le(2.04e)^2$ for $m\ge5$), and
$\Lambda\le\sigma_mx^{m-2}$: indeed
$(e/\alpha)^{(m-2)/2}\le(2.04e)^{(m-2)/2}$, and by
Corollary~\ref{cor:gate-smalle} and $m\le13$,
\[
 \frac1{m\kappa}\ <\ \frac{m-1}{0.763\,m}\ \le\ \frac{12}{0.763\cdot13}
 \ =\ 1.20979\ldots\ <\ 1.21 .
\]
Since $x^{m-2}\le(1+z+c_mz^2)^{-1}$ [R2, eq.~E-a] and
$1/x\le1+\beta_mz$, $\frac1{mx}\ge\frac1m$, \eqref{eq:red} follows from
\begin{equation}\label{eq:F}
 F:=\bigl(1-\Gamma_0-\eps\bigr)(1+z+c_mz^2)-\frac1x+\frac{2+1/\kappa}m-\sigma_m\ \ge\ 0 .
\end{equation}
We record two endpoint evaluations, rounded up
($2.8524/\sqrt{60}<2.8524/7.7459<0.36825$, since $7.7459^2=59.99896681<60$;
$0.51/60=0.0085$; $(2.04/60)^2=0.001156$; $1.293/60=0.02155$):
\begin{equation}\label{eq:G-ends}
 \Gamma_0(e)\le\Gamma_0(1/60)<0.378,\qquad
 \Gamma_1(e):=\Gamma_0(e)+1.293e\le\Gamma_1(1/60)<0.3995,
\end{equation}
both $\Gamma_0,\Gamma_1$ being increasing in $e$; and
$\sigma_m(e)\le\sigma_5(1/60)=1.21\cdot(0.034)^{3/2}<0.0076$ for all
$5\le m\le13$ (as $2.04e\le0.034<1$).

\emph{Sub-case $\kappa<\tfrac15$.}  Here $\frac{2+1/\kappa}m\ge\frac7m$, and
by \eqref{eq:eps-smalle} $\eps\le0.068$.  Since
$1-\Gamma_0-0.068-\beta_m\ge1-0.378-0.068-0.35704>0.19>0$
($\beta_m\le\beta_5=1.0711/3=0.35703\ldots$), dropping $c_mz^2\ge0$,
\[
 F\ \ge\ z\bigl(1-\Gamma_0-0.068-\beta_m\bigr)+\frac7m-\Gamma_0-0.068-\sigma_m
 \ \ge\ \frac7{13}-0.378-0.068-0.0076\ >\ 0.084\ >\ 0 .
\]

\emph{Sub-case $\kappa\ge\tfrac15$.}  Here \eqref{eq:eps-smalle} gives
$\eps\le\frac{5e}{4(1-e)^2}\le\frac{4500}{3481}e\le1.293e$, so
$1-\Gamma_0-\eps\ge1-\Gamma_1>0$ by \eqref{eq:G-ends}.  Dropping
$c_mz^2\ge0$ and using $1/x\le1+\beta_mz$, $\frac{2+1/\kappa}m\ge\frac3m$
($\kappa<1$ on the domain):
\[
 F\ \ge\ z\bigl(1-\Gamma_1-\beta_m\bigr)+\frac3m-\Gamma_1-\sigma_m .
\]
The bracket is positive ($\ge1-0.3995-0.35704>0.24$), and
$z\ge1.9047\cdot1.2\,(m-2)e\ge2.2856(m-2)e$ (from $1+\kappa\ge1.2$), so
\begin{equation}\label{eq:Gm}
 F\ \ge\ G_m(e):=2.2856(m-2)e\bigl(1-\Gamma_1(e)-\beta_m\bigr)
 +\frac3m-\Gamma_1(e)-\sigma_m(e) .
\end{equation}
Bounding the bracket once more by $\Gamma_1\le0.3995$ and expanding
$\Gamma_1(e)=2.8524\sqrt e+1.803e+4.1616e^2$,
\begin{equation}\label{eq:Ghat}
 G_m(e)\ \ge\ \widehat G_m(e):=\frac3m+\bigl(c_{1,m}e-2.8524\sqrt e\bigr)
 -4.1616e^2-\sigma_m(e),
\end{equation}
\begin{equation}\label{eq:c1}
 c_{1,m}:=2.2856(m-2)\bigl(0.6005-\beta_m\bigr)-1.803
 \ =\ 1.3725028\,(m-2)-4.25110616,
\end{equation}
an exact terminating decimal for each $m$ (Table~\ref{tab:smalle}).

\emph{$\widehat G_m$ is strictly decreasing on $(0,1/60]$.}  The three terms
$-4.1616e^2$, $-\sigma_m(e)$, and the constant are nonincreasing in $e$; for
the remaining term, its derivative satisfies, for $0<e\le1/60$,
\[
 \frac{d}{de}\bigl(c_{1,m}e-2.8524\sqrt e\bigr)
 =c_{1,m}-\frac{1.4262}{\sqrt e}
 \ \le\ c_{1,m}-1.4262\sqrt{60}
 \ <\ 10.847-11.045\ <\ 0,
\]
because $c_{1,m}$ is increasing in $m$ with
$c_{1,13}=10.84642464<10.847$, and
$1.4262\sqrt{60}>1.4262\cdot7.745=11.045919>11.045$
($7.745^2=59.985025<60$).  Hence
$\widehat G_m(e)\ge\widehat G_m(1/60)$ on $(0,1/60]$, and evaluating with
the safe roundings above ($2.8524\sqrt{1/60}<0.36825$,
$4.1616/3600<0.001156\,1$, $\sigma_m(1/60)\le\bar\sigma_m$ of
Table~\ref{tab:smalle}, $c_{1,m}\ge\underline c_{1,m}$ of
Table~\ref{tab:smalle}):
\[
 \widehat G_m(1/60)\ \ge\ \frac3m+\frac{\underline c_{1,m}}{60}
 -0.36825-0.0011561-\bar\sigma_m\ \ge\ \widehat G_m^{\,\flat}
\]
with the values $\widehat G_m^{\,\flat}$ of Table~\ref{tab:smalle}, all
positive (minimum $0.0383$ at $m=11$).  This proves $F\ge0$ in both
sub-cases, hence \eqref{eq:red}, hence the proposition via
Lemma~\ref{lem:reduce} and Lemma~\ref{lem:relax}.
\end{proof}

\begin{table}\centering
\begin{tabular}{lccccc}\toprule
$m$&5&7&9&11&13\\\midrule
$\beta_m$&$0.35704$&$0.21422$&$0.15302$&$0.11902$&$0.09738$\\
$\underline c_{1,m}$ &$-0.13360$&$2.61140$&$5.35641$&$8.10141$&$10.84642$\\
$\bar\sigma_m$&$0.0076$&$0.000258$&$0.0000088$&$0.0000003$&$0.00000002$\\
$\widehat G_m^{\,\flat}$&$0.2207$&$0.1024$&$0.0531$&$0.0383$&$0.0421$\\
\bottomrule
\end{tabular}
\caption{Small-$e$ constants.  $\beta_m$ is rounded up (safe in
$1-\Gamma_1-\beta_m$ and in $c_{1,m}$ via \eqref{eq:c1});
$\underline c_{1,m}$ is the exact value \eqref{eq:c1} truncated toward
$-\infty$ at the fifth decimal; $\bar\sigma_m$ is an upper bound for
$\sigma_m(1/60)=1.21(0.034)^{(m-2)/2}$; $\widehat G_m^{\,\flat}$ is the
resulting lower bound for $\widehat G_m(1/60)$, truncated down.}
\label{tab:smalle}
\end{table}

\subsection{The Tur\'an band $\delta\le\Delta_m$: polynomial Bernstein certificates}

Here $\delta=\frac{1-3e}6\in(0,\Delta_m]$ with
$\Delta_m=\tfrac1{40}$ for $m\le9$ and $\Delta_m=\tfrac1{80}$ for
$m\in\{11,13\}$; write $s:=1-d/\delta\in[0,1)$, so
$\alpha=\tfrac13+\delta$, $e=\tfrac13-2\delta$, $q=\tfrac13+s\delta$,
$p=\tfrac23-s\delta$, and
$L^2=\tfrac19-\tfrac{(2-s)\delta}3-(1+s^2)\delta^2$, a polynomial in
$(\delta,s)$.  The band covers all admissible $\kappa$ (there
$\kappa_q<\kappa_{\max}$).

\begin{lemma}[Chord]\label{lem:chord}
$\sqrt{2\alpha}\ \ge\ w_\ell(\delta):=0.8164+1.2\,\delta$ for
$0\le\delta\le\tfrac1{40}$.
\end{lemma}

\begin{proof}
$w_\ell(\delta)^2-(2/3+2\delta)$ is a quadratic in $\delta$ with positive
leading coefficient, hence maximal at the endpoints, where it is
negative: $0.8164^2=0.66650\ldots<2/3$ and
$(0.8164+0.03)^2=0.71639296<2/3+1/20=0.71666\ldots$.
\end{proof}

\begin{lemma}[Tail]\label{lem:tail}
Let $t_m^*$ be as in Table~\ref{tab:band}.  If $R_m>0$ and
$\delta\le\Delta_m$, then $d/\delta>t^*_m$, i.e.\ $s<1-t^*_m$.
\end{lemma}

\begin{proof}
\emph{Step 1 (gate $\Rightarrow$ curve).}
By Lemma~\ref{lem:gate}, $\kappa(m-1-x)>x(1-\ell^{m-2})$; on the band
$x=\alpha/p\ge\tfrac12$ (as $\alpha\ge\tfrac13\ge p/2$) and $m-1-x<m-1$,
and $\ell^{m-2}\le w(\delta)^{\nu}$ where
\[
 w(\delta):=\frac e\alpha=\frac{1-6\delta}{1+3\delta},\qquad
 \nu:=\frac{m-2}2
\]
(from $\ell^2\le e/\alpha$, [R2, Lem.~elementary]).  Hence, for $R_m>0$,
\begin{equation}\label{eq:phi}
 \frac d\delta=\frac{\kappa e}\delta
 \ >\ \phi_m(\delta):=\frac{e}{2(m-1)\delta}\,h(\delta),
 \qquad h(\delta):=1-w(\delta)^{\nu}.
\end{equation}

\emph{Step 2 (endpoint reduction: $\phi_m(\delta)\ge\phi_m(\Delta_m)$ on
$(0,\Delta_m]$).}
We claim $h$ is concave on $[0,\Delta_m]$.  Writing
$w=-2+3/(1+3\delta)$, so $w'=-9/(1+3\delta)^2$ and $w''=54/(1+3\delta)^3$,
\[
 (w^{\nu})''=\nu w^{\nu-2}\bigl[(\nu-1)w'^2+w\,w''\bigr]
 =\nu w^{\nu-2}\,\frac{27\,\bigl(3\nu-1-12\delta\bigr)}{(1+3\delta)^4}
 \ \ge\ 0
\]
on $[0,\Delta_m]$: indeed $w>0$ there (as $\delta\le\tfrac1{40}<\tfrac16$),
and $3\nu-1-12\delta\ge0$ because for $m\ge5$, $3\nu\ge\tfrac92>1+\tfrac{12}{40}$,
while for $m=3$, $3\nu=\tfrac32\ge1+12\delta$ iff $\delta\le\tfrac1{24}$,
and $\Delta_3=\tfrac1{40}<\tfrac1{24}$.  So $w^{\nu}$ is convex, $h$ is
concave, and $h(0)=0$; concavity then gives, for $0<\delta\le\Delta_m$,
\[
 h(\delta)\ \ge\ \tfrac{\delta}{\Delta_m}h(\Delta_m)
 +\bigl(1-\tfrac{\delta}{\Delta_m}\bigr)h(0)
 =\tfrac{\delta}{\Delta_m}h(\Delta_m),
 \qquad\text{i.e.}\qquad
 \frac{h(\delta)}\delta\ \ge\ \frac{h(\Delta_m)}{\Delta_m}.
\]
Combining with $e=\tfrac13-2\delta\ge\tfrac13-2\Delta_m$ (both factors
positive),
\[
 \phi_m(\delta)=\frac{e}{2(m-1)}\cdot\frac{h(\delta)}{\delta}
 \ \ge\ \frac{\tfrac13-2\Delta_m}{2(m-1)}\cdot\frac{h(\Delta_m)}{\Delta_m}
 \ =\ \phi_m(\Delta_m).
\]

\emph{Step 3 (endpoint comparison, exact).}
It remains to check $\phi_m(\Delta_m)>t^*_m$, i.e.\ with the rationals
\[
 c_m:=\frac{\tfrac13-2\Delta_m}{2(m-1)\Delta_m},\qquad
 w_m:=w(\Delta_m),\qquad r_m:=1-\frac{t^*_m}{c_m},
\]
that $w_m^{\nu}<r_m$.  Table~\ref{tab:tailints} lists $c_m$, $w_m$, $r_m$;
each $r_m\in(0,1)$, so squaring gives the equivalent \emph{rational}
inequality $w_m^{m-2}<r_m^2$, i.e.\ the one-line integer comparison in the
last column of Table~\ref{tab:tailints}, which is verified exactly (and can
be checked by hand or by any exact integer arithmetic).  This proves
$d/\delta>\phi_m(\delta)\ge\phi_m(\Delta_m)>t^*_m$.
\end{proof}

\begin{table}\centering
\small
\begin{tabular}{cccccl}\toprule
$m$&$\Delta_m$&$w_m$&$c_m$&$r_m$&integer form of $w_m^{m-2}<r_m^2$\\\midrule
$3$&$1/40$&$34/43$&$17/6$&$757/850$&$34\cdot850^2=24\,565\,000<24\,641\,107=43\cdot757^2$\\
$5$&$1/40$&$34/43$&$17/12$&$299/425$&$34^3\!\cdot\!425^2=7\,099\,285\,000<7\,108\,005\,307=43^3\!\cdot\!299^2$\\
$7$&$1/40$&$34/43$&$17/18$&$481/850$&$34^5\!\cdot\!850^2=32\,827\,093\,840\,000<34\,012\,020\,380\,923=43^5\!\cdot\!481^2$\\
$9$&$1/40$&$34/43$&$17/24$&$191/425$&$34^7\!\cdot\!425^2=9\,487\,030\,119\,760\,000<9\,916\,214\,751\,794\,467=43^7\!\cdot\!191^2$\\
$11$&$1/80$&$74/83$&$37/30$&$223/370$&$74^9\!\cdot\!370^2<83^9\!\cdot\!223^2$ \ (see below)\\
$13$&$1/80$&$74/83$&$37/36$&$493/925$&$74^{11}\!\cdot\!925^2<83^{11}\!\cdot\!493^2$ \ (see below)\\
\bottomrule
\end{tabular}
\caption{Tail-lemma endpoint data.  All entries are exact rationals:
$w_m=e(\Delta_m)/\alpha(\Delta_m)$, $c_m$ as in Step~3, $r_m=1-t^*_m/c_m$.
The two large comparisons are
$74^9\cdot370^2=9\,109\,382\,235\,108\,373\,145\,600
<9\,296\,351\,954\,199\,516\,700\,787=83^9\cdot223^2$ and
$74^{11}\cdot925^2=311\,768\,606\,996\,584\,070\,908\,160\,000
<313\,006\,138\,444\,263\,224\,418\,858\,083=83^{11}\cdot493^2$.}
\label{tab:tailints}
\end{table}

\begin{proposition}[Band]\label{prop:band}
For odd $3\le m\le13$, $R_m>0$ and $\delta\le\Delta_m$ imply $(T_m)$.
\end{proposition}

\begin{proof}
Write $A_m=A^e+LA^o$, $B_m=B^e+LB^o$, $R_m=R^e+LR^o$, $f=\alpha-L$ with
$A^e,\dots$ polynomials in $(\delta,s)$ (possible since $m-1$ is even and
$p^{m-1}-L^{m-1}$ is divisible by $p^2-L^2$).  Substituting
$w_\ell\le\sqrt{2\alpha}$ (Lemma~\ref{lem:chord}) into \eqref{eq:Tm2} and
collecting $L$-parities, $(T_m)$ follows from $G:=P+LQ\ge0$, where
\[
 P+LQ\ :=\ w_\ell(\delta)\,4\alpha^2d\,(A B f)\big|_{\text{parity}}
 -e^2\,(B f)^2\big|_{\text{parity}}-4\alpha^2 (A R)\big|_{\text{parity}},
\]
two explicit polynomials in $(\delta,s)$ with rational coefficients
(generated and archived by the validation script).  Let
$H:=L^2Q^2-P^2$; then $H=c_m\,\delta\,K$ with $c_m>0$ and $K$ a polynomial.
On the box $\mathcal B_m:=[0,\Delta_m]\times[0,1-t^*_m]$ (which by
Lemma~\ref{lem:tail} contains all band points with $R_m>0$):
\emph{all two-dimensional Bernstein coefficients of $Q$ and of $K$ on
$\mathcal B_m$ are nonnegative} (exact rational computation; no subdivision
is needed).  Hence $Q\ge0$ and $H\ge0$ on $\mathcal B_m$; if $P\ge0$ then
$G\ge0$ trivially, and if $P<0$ then $L^2Q^2\ge P^2$ gives $LQ\ge|P|$, so
$G\ge0$.  By \eqref{eq:Tm2} and Lemma~\ref{lem:relax} this proves the claim.
\end{proof}

\begin{table}\centering
\begin{tabular}{lcccccc}\toprule
$m$&3&5&7&9&11&13\\\midrule
$\Delta_m$&$1/40$&$1/40$&$1/40$&$1/40$&$1/80$&$1/80$\\
$t^*_m$&$0.31$&$0.42$&$0.41$&$0.39$&$0.49$&$0.48$\\\bottomrule
\end{tabular}\caption{Band parameters.}\label{tab:band}
\end{table}

\subsection{The moderate range: certified subdivision}

\begin{proposition}\label{prop:moderate}
For odd $3\le m\le13$ and $e\in[\tfrac1{60},\tfrac13-2\Delta_m]$, either
$R_m\le0$ or $(T_m)$ holds.
\end{proposition}

\begin{proof}[Certificate]
Exact-rational interval subdivision over
$(e,\kappa)\in[\tfrac1{60},\tfrac13-2\Delta_m]\times[0,1]$: on each box,
directed rational bounds for $\alpha,d,q,p,L^2,f,k_m(\alpha),k_m(L),A_m,B_m,
R_m$ and integer-square-root brackets for $L$, $\sqrt{2\alpha}$ certify
either $R_m\le0$ or \eqref{eq:Tm2}; boxes wholly outside the admissible
domain ($q\le\tfrac13$ or $\kappa>\kappa_{\max}$) are discarded.  The run
(same protocol as [R2, \S certificate-protocol]) terminates with
$6757,3911,3567,4307,6487,8785$ boxes for $m=3,5,\dots,13$.  This step is
\emph{computational}; a prose replacement is open.
\end{proof}

\subsection{Assembly}

\begin{proof}[Proof of Theorem~\ref{thm:smallm-main}]
If $R_m\le0$, \eqref{eq:smallm-target} is trivial ($\psi\ge0$).  If $R_m>0$:
for $e\le1/60$, Proposition~\ref{prop:smalle} applies (Corollary
\ref{cor:gate-smalle} shows this range meets $\xi>1$);
for $1/60\le e\le\tfrac13-2\Delta_m$, Proposition~\ref{prop:moderate}
gives $(T_m)$; for $e\ge\tfrac13-2\Delta_m$ (i.e.\ $\delta\le\Delta_m$),
Proposition~\ref{prop:band} gives $(T_m)$.  In the latter two cases
Lemma~\ref{lem:relax} upgrades $(T_m)$ to \eqref{eq:smallm-target}.
Finally, Remark~\ref{rem:muniform} converts \eqref{eq:smallm-target} into the
odd-cycle bound in Region II.
\end{proof}

\begin{remark}[Computational components]\label{rem:comp}
Two steps are certificate-based rather than prose: (i) the Bernstein
coefficient tables of Proposition~\ref{prop:band} (a classical positivity
certificate: nonnegativity of all Bernstein coefficients of an explicit
rational polynomial on a box implies nonnegativity of the polynomial), and
(ii) the subdivision of Proposition~\ref{prop:moderate}.  Both are
exact-rational and small.  Everything else is analytic; in particular the
tail cut (Lemma~\ref{lem:tail}) is analytic up to the six displayed
one-line integer comparisons of Table~\ref{tab:tailints}.  The
per-cycle theorems $C_5,\dots,C_{13}$ are no longer used.
\end{remark}

\section{Zone B without boxes}\label{sec:zoneB}

\subsection*{Provenance and verification status}
Proved by Claude on 2026-07-16/17; independent verifier: \emph{pass}
(414/414 checks); adversarial reviewer: \emph{sound}.  The statement of
Lemma~\emph{zoneB-battle} of [R2] is unchanged, so the surrounding proof
of Lemma~\emph{zoneB} in [R2] goes through verbatim and
\texttt{zoneB\_certifier.py} can be deleted from the certificate
protocol.  The proof consists of a monotonicity toolkit, a gate argument
along two explicit curves, and $15$ displayed hand-checkable table rows.
Two by-products (Remarks at the end of this section): the
$17.37$-exponential in the battle statement and the $\max$ in the
definition of $\mathfrak K$ are both removable.  The content is the note
\texttt{claude\_jul17/zoneB\_analytic.tex}, inlined verbatim.

\noindent
\textbf{Scope.}  This note replaces the computational proof (the
$23$-box exhaustive-subdivision certificate \texttt{zoneB\_certifier.py})
of the Zone-B battle, Lemma~\texttt{lem:zoneB-battle} of
\texttt{paper\_region2\_v2.tex} (henceforth \RII), by a fully analytic
proof.  The \emph{statement} of the battle lemma is unchanged, so the proof
of Lemma~\texttt{lem:zoneB} in \RII{} (which combines the battle with
Lemma~\texttt{lem:Kmax}, the bounds
$\Pi\ge\underline\Pi$, $\Lambda\le\overline\Lambda$, and the certificate
\texttt{eq:cert-II}) goes through verbatim.  The only results quoted from
\RII{} are: the chart formulas of Lemma~\texttt{lem:chart}, and items (ii)
and (iv) of Lemma~\texttt{lem:Kmax}; everything else below is
self-contained.  Every displayed inequality is validated by
\texttt{validate\_zoneB.py} (dense float grids over the exact stated
domains, plus exact rational arithmetic for every table row, every
surrogate constant, and a family of corner spot checks).

\subsection{Setting and statement}

Throughout, $e\in[\tfrac1{60},\tfrac{2033}{10000}]$ and
\[
        \kxi(e)=\frac e{(1-e)^2},\qquad
        \kbar(e)=\min\Bigl(\frac{1-e}{1+e},\ \frac{1-3e}{6e}\Bigr),\qquad
        \kappa\in[\kxi(e),\kbar(e)]
\]
(the \emph{strip}; it is empty unless $\kxi\le\kbar$, and all statements are
vacuous at such $e$).  From the chart (Lemma~\texttt{lem:chart} of \RII):
\[
        x=x(e,\kappa)=\frac{1-e}{1+e+2\kappa e},\qquad
        u=1-x=\frac{2e(1+\kappa)}{1+e+2\kappa e},\qquad
        A=A(e,\kappa)=\frac{x(1+\kappa)}\kappa ,
\]
\[
        \Phi=\Phi(e,\kappa)=\frac{2(1-e)(1+\kappa)^2e}{\kappa(1+e+2\kappa e)^2},
        \qquad
        w(\nu)=\frac{\sqrt{\nu^2+4\nu}-\nu}2 ,
        \qquad
        \lambda=\ln(1/x).
\]
Since $\nu+2w(\nu)=\sqrt{\nu^2+4\nu}$, the quantity
$e^{-\Phi-2w(\Phi)}$ of \RII{} equals
\begin{equation}\label{eq:Kstar-def}
        \Kstar(e,\kappa):=\exp\Bigl(-\sqrt{\Phi^2+4\Phi}\Bigr).
\end{equation}
The battle lemma to be proved is
(Lemma~\texttt{lem:zoneB-battle} of \RII):
for all $e\in[\tfrac1{60},\tfrac{2033}{10000}]$ and
$\kappa\in[\kxi(e),\kbar(e)]$,
\begin{equation}\label{eq:zb-battle}
        \underline\Pi(e,\kappa)\ \ge\
        \frac{\Kfrak(e,\kappa)}{x^2}+\overline\Lambda(e,\kappa),
\end{equation}
where, with $\lbar=\sqrt{2e/(1-e)}$,
$\underline\rho=\tfrac14(1-e^{-17.37(1+\kappa)e})$ and
$\epsbar=\min\bigl(\tfrac14,\;\frac e{4(1-e)^2\kappa(1+\underline\rho)}\bigr)$,
\[
        \underline\Pi
        =\sqrt{1-e}\;\frac{(1-\lbar)(1-\lbar^{14})(1-\epsbar)}{1+\lbar},
        \qquad
        \overline\Lambda
        =x^{13}\Bigl(\frac{2e}{1-e}\Bigr)^{13/2}\frac{(1-e)^2}{15e},
\]
and $\Kfrak=x^{14}(14-A)_+/15$ if the \emph{gate}
``$A<14$ and $15(14-A)\lambda\ge A+1$'' holds, while
$\Kfrak=\max\bigl(x^{14}(14-A)_+/15,\;\Kstar\bigr)$ otherwise.

Two explicit curves organise the proof:
\[
        \gamma(e):=\frac7{25}-e
        \qquad\text{(the \emph{gate curve})},
        \qquad
        G(e):=\sqrt{1-e}\;\frac{(1-\lbar)(1-\lbar^{14})}{1+\lbar},
\]
and the branch point of $\max(\gamma,\kxi)$ lies in
$(\tfrac{111}{1000},\tfrac18)$; we only use the rational comparison
$\kxi(\tfrac18)=\tfrac8{49}>\tfrac{31}{200}=\gamma(\tfrac18)$
(equivalently $8\cdot200>49\cdot31$).
The $e$-range is split at $e=\tfrac18$: for $e\le\tfrac18$ we work with the
curve $\kappa=\gamma(e)$, for $e\ge\tfrac18$ with the curve
$\kappa=\kxi(e)$.

\begin{convention}[safe decimals]\label{conv:zb-safe}
Every terminating decimal below denotes the displayed exact rational
number.  In the verification tables, quantities that must be bounded
\emph{above} (marked~$\uparrow$) are exact rationals rounded up, and
quantities that must be bounded \emph{below} (marked~$\downarrow$) are
rounded down, so each table row is, as displayed, a finite exact rational
computation that certifies its inequality.  The script
\texttt{validate\_zoneB.py} re-verifies every row from the unrounded
rationals.
\end{convention}

\subsection{Monotonicity toolkit}

\begin{lemma}[elementary monotonicities]\label{lem:zb-mono}
For $e\in(0,\tfrac13)$ and $\kappa>0$, with $D:=1+e+2\kappa e$:
\begin{enumerate}[label=\textup{(\roman*)},itemsep=1pt]
\item $x=(1-e)/D\in(0,1)$ is strictly decreasing in $e$ and in $\kappa$.
\item $u=1-x=2e(1+\kappa)/D$ is strictly increasing in $e$ and in $\kappa$:
$\partial_e u=2(1+\kappa)/D^2>0$ and $\partial_\kappa u=2e(1-e)/D^2>0$.
\item $A=(1-e)(1+\kappa)/(\kappa D)$ is strictly decreasing in $e$ and in
$\kappa$.
\item $\Phi$ is strictly decreasing in $\kappa$ on $\kappa\in(0,1]$;
and for $e\in[\tfrac1{60},\tfrac18]$, $\kappa\le\tfrac{79}{300}$, $\Phi$ is
strictly increasing in $e$.
\item $t\mapsto\exp(-\sqrt{t^2+4t})$ is strictly decreasing on $t\ge0$;
hence by \eqref{eq:Kstar-def}, $\Kstar$ is increasing in $\kappa$ on
$\kappa\in(0,1]$ and, on the domain of \textup{(iv)}, decreasing in $e$.
\item $G$ is positive and strictly decreasing on $(0,\tfrac13)$.
\item $\kxi$ is strictly increasing; $\gamma$ is strictly decreasing;
$\kxi/(4\gamma)$ is strictly increasing on $[\tfrac1{60},\tfrac18]$, with
maximal value $\kxi(\tfrac18)/(4\gamma(\tfrac18))=\tfrac{400}{1519}<1$.
\item $e\mapsto x(e,\gamma(e))=\dfrac{1-e}{1+\tfrac{39}{25}e-2e^2}$ is
strictly decreasing on $[\tfrac1{60},\tfrac18]$.
\item $e\mapsto x(e,\kxi(e))$ is strictly decreasing on
$[\tfrac18,\tfrac{2033}{10000}]$.
\item $\lambda=\ln(1/x)\ge u=1-x$ for $x\in(0,1]$.
\end{enumerate}
\end{lemma}

\begin{proof}
(i) In $\kappa$: $\partial_\kappa x=-2e(1-e)/D^2<0$.  In $e$: the numerator
$1-e$ decreases and $D=1+e(1+2\kappa)$ increases.
(ii) Direct differentiation as displayed.
(iii) In $e$: by (i), $x$ decreases and the factor $(1+\kappa)/\kappa$ is
constant.  In $\kappa$: with $N=1+\kappa$ and $\tilde D=\kappa D
=\kappa(1+e)+2\kappa^2e$,
\[
        N'\tilde D-N\tilde D'
        =\kappa(1+e)+2\kappa^2e-(1+\kappa)\bigl((1+e)+4\kappa e\bigr)
        =-(1+e)-4\kappa e-2\kappa^2e<0 .
\]
(iv) $\partial_\kappa\ln\Phi=\frac2{1+\kappa}-\frac1\kappa-\frac{4e}D
=\frac{\kappa-1}{\kappa(1+\kappa)}-\frac{4e}D\le-\frac{4e}D<0$ for
$\kappa\le1$.  Next,
$\partial_e\ln\Phi=\frac1e-\frac1{1-e}-\frac{2(1+2\kappa)}D
\ge\frac1e-\frac1{1-e}-2(1+2\kappa)$ (as $D\ge1$); for $e\le\tfrac18$ and
$\kappa\le\tfrac{79}{300}$ this is at least
$8-\tfrac87-2\bigl(1+\tfrac{158}{300}\bigr)=\tfrac{1997}{525}>0$.
(v) $t^2+4t$ is increasing on $t\ge0$; combine with (iv).
(vi) $\lbar=\sqrt{2e/(1-e)}$ is increasing in $e$ and $\lbar<1$ iff
$e<\tfrac13$; the three factors $\sqrt{1-e}$, $(1-\lbar)/(1+\lbar)$ and
$1-\lbar^{14}$ are positive and decreasing.
(vii) $\kxi$: numerator increases, denominator decreases.  The quotient of
a positive increasing by a positive decreasing function increases; its
value at $e=\tfrac18$ is $\frac{8/49}{4\cdot31/200}=\frac{400}{1519}$.
(viii) With $\kappa=\gamma(e)$ one gets $2\kappa e=\tfrac{14}{25}e-2e^2$,
so $D=1+\tfrac{39}{25}e-2e^2$, whose derivative $\tfrac{39}{25}-4e$ is
positive for $e\le\tfrac18<\tfrac{39}{100}$; the numerator $1-e$ decreases.
(ix) Composition of (i) with $\kxi$ increasing (vii).
(x) $-\ln x\ge1-x$.
\end{proof}

\begin{lemma}[the $\Lambda$-bound]\label{lem:zb-lambda}
For all $(e,\kappa)$ in the strip,
\[
        \overline\Lambda(e,\kappa)\ \le\
        \Lambda_1(e):=\Bigl(\frac{1-e}{1+e}\Bigr)^{13}
        \Bigl(\frac{2e}{1-e}\Bigr)^{13/2}\frac{(1-e)^2}{15e}
        \ \le\ \widehat\Lambda:=\frac{13}{10^6}.
\]
\end{lemma}

\begin{proof}
The first inequality is $x\le(1-e)/(1+e)$ (Lemma~\ref{lem:zb-mono}(i) with
$\kappa\downarrow0$), all other factors of $\overline\Lambda$ being
$\kappa$-free.  Collecting exponents,
\[
        \Lambda_1(e)=\frac{2^{13/2}}{15}\,
        \frac{e^{11/2}(1-e)^{17/2}}{(1+e)^{13}},
        \qquad
        (\ln\Lambda_1)'(e)=\frac{11}{2e}-\frac{17}{2(1-e)}-\frac{13}{1+e}.
\]
For $e\le e_1:=\tfrac{2033}{10000}$ each term is bounded by its value at
$e_1$ resp.\ by $13$:
\[
        (\ln\Lambda_1)'(e)\ \ge\ \frac{11}{2e_1}-\frac{17}{2(1-e_1)}-13
        =\frac{55000}{2033}-\frac{85000}{7967}-13>3.38>0
\]
(exact rational comparison: $55000\cdot7967-85000\cdot2033-13\cdot2033\cdot
7967=54{,}820{,}157>0$).  Hence $\Lambda_1$ is increasing on $(0,e_1]$ and
\[
        \Lambda_1(e)\le\Lambda_1(e_1)
        =\frac{2^6}{15}\,\sqrt{2e_1(1-e_1)}\;
        \frac{e_1^5(1-e_1)^8}{(1+e_1)^{13}}
        \ \le\ \frac{64}{15}\cdot\frac{14229}{25000}\cdot
        \frac{e_1^5(1-e_1)^8}{(1+e_1)^{13}}
        \ \le\ \frac{124}{10^7}\ <\ \widehat\Lambda,
\]
using $\bigl(\tfrac{14229}{25000}\bigr)^2\ge2e_1(1-e_1)
=\tfrac{32393822}{10^8}$; the final two comparisons are exact rational
evaluations (script, \S1--2).
\end{proof}

\begin{lemma}[envelope and the $\Kfrak$-bounds]\label{lem:zb-envelope}
At every strip point:
\begin{enumerate}[label=\textup{(\alph*)},itemsep=1pt]
\item $\underline\Pi\ \ge\ G(e)\,(1-\epsbar)$ with
$\epsbar\le\min\bigl(\tfrac14,\ \tfrac{\kxi(e)}{4\kappa}\bigr)$;
in particular $\underline\Pi\ge\tfrac34G(e)$.
\item $\Kfrak\le\Kstar(e,\kappa)$ (in both gate branches); consequently the
$\max(\cdot,\cdot)$ in the definition of $\Kfrak$ is redundant.
\item If the gate holds at $(e,\kappa)$ then
$\Kfrak\le\tfrac{14}{15}x^{14}$.
\end{enumerate}
\end{lemma}

\begin{proof}
(a) $\underline\Pi=G(e)(1-\epsbar)$ by definition, and
$\epsbar\le\frac e{4(1-e)^2\kappa(1+\underline\rho)}
\le\frac e{4(1-e)^2\kappa}=\frac{\kxi(e)}{4\kappa}$ because
$\underline\rho\ge0$; also $\epsbar\le\tfrac14$ by its definition as a
minimum, and $\kappa\ge\kxi$ gives again $\epsbar\le\tfrac14$.

(b) Write $K(n)=x^n(n-A)/(n+1)$.  If $A\ge14$ then
$x^{14}(14-A)_+/15=0\le\Kstar$.  If $A<14$, then by
Lemma~\texttt{lem:Kmax}(ii) of \RII{} the function $K$ is log-concave on
$(A,\infty)$ with an interior stationary point $n_*$ (the unique root of
$\lambda=(A+1)/((n-A)(n+1))$, which exists because the right side decreases
from $+\infty$ to $0$), so $K(14)\le K(n_*)$; and by
Lemma~\texttt{lem:Kmax}(iv), $K(n_*)\le e^{-\Phi-2w(\Phi)}=\Kstar$.  Thus
$x^{14}(14-A)_+/15\le\Kstar$ always.  On the gate branch
$\Kfrak=x^{14}(14-A)_+/15\le\Kstar$; on the other branch
$\Kfrak=\max\bigl(x^{14}(14-A)_+/15,\Kstar\bigr)=\Kstar$.

(c) On the gate branch, $(14-A)_+\le14$.
\end{proof}

\subsection{The gate on the two curves}

\begin{lemma}[$u$-gate]\label{lem:zb-ugate}
Say the \emph{$u$-gate} holds at $(e,\kappa)$ if $A<14$ and
$15(14-A)u\ge A+1$.  Then:
\begin{enumerate}[label=\textup{(\alph*)},itemsep=1pt]
\item ($u$-gate $\Rightarrow$ gate)  If the $u$-gate holds at $(e,\kappa)$,
so does the gate of \eqref{eq:zb-battle}, since
$\lambda\ge u$ (Lemma~\ref{lem:zb-mono}(x)) and $14-A>0$ give
$15(14-A)\lambda\ge15(14-A)u\ge A+1$.
\item (upward extension)  If the $u$-gate holds at $(e,\kappa_0)$, it holds
at $(e,\kappa)$ for every $\kappa\ge\kappa_0$: by
Lemma~\ref{lem:zb-mono}(ii),(iii), $A(e,\kappa)\le A(e,\kappa_0)<14$ and
$u(e,\kappa)\ge u(e,\kappa_0)$, so
$15\bigl(14-A(e,\kappa)\bigr)u(e,\kappa)
\ge15\bigl(14-A(e,\kappa_0)\bigr)u(e,\kappa_0)\ge A(e,\kappa_0)+1\ge
A(e,\kappa)+1$. \qedhere
\end{enumerate}
\end{lemma}

\begin{proposition}[gate on the curves]\label{prop:zb-gate}
\begin{enumerate}[label=\textup{(\alph*)},itemsep=1pt]
\item For every $e\in[\tfrac1{60},\tfrac18]$ the $u$-gate holds at
$(e,\gamma(e))$, hence the gate holds at $(e,\kappa)$ for all
$\kappa\ge\gamma(e)$.
\item For every $e\in[\tfrac18,\tfrac{2033}{10000}]$ the $u$-gate holds at
$(e,\kxi(e))$, hence the gate holds at every strip point with
$e\ge\tfrac18$.
\end{enumerate}
\end{proposition}

\begin{proof}
(a) Fix a row $[a,b]$ of Table~\ref{tab:gate} (the three rows tile
$[\tfrac1{60},\tfrac18]$) and let $e\in[a,b]$, $\kappa=\gamma(e)\in
[\gamma(b),\gamma(a)]$.  By Lemma~\ref{lem:zb-mono}(i),(iii),
\[
        A(e,\gamma(e))=x(e,\gamma(e))\Bigl(1+\frac1{\gamma(e)}\Bigr)
        \ \le\ x\bigl(a,\gamma(b)\bigr)\Bigl(1+\frac1{\gamma(b)}\Bigr)
        =:\overline A,
\]
since $x$ decreases in both arguments and $1+1/\gamma(e)\le1+1/\gamma(b)$;
and by Lemma~\ref{lem:zb-mono}(ii),
$u(e,\gamma(e))\ge u(a,\gamma(b))=:\underline u$.
The row certifies $\overline A<14$ and
$15(14-\overline A)\,\underline u-(\overline A+1)>0$; monotonicity then
gives $15(14-A)u\ge15(14-\overline A)\underline u\ge\overline A+1\ge A+1$,
i.e.\ the $u$-gate at $(e,\gamma(e))$.  Lemma~\ref{lem:zb-ugate} extends it
upward and converts it to the gate.

(b) Same argument on the single row $[\tfrac18,\tfrac{2033}{10000}]$: by
Lemma~\ref{lem:zb-mono}(ii),(iii),(vii), for $e\ge a=\tfrac18$,
\[
        A(e,\kxi(e))\le A(e,\kxi(a))\le A(a,\kxi(a))=:\overline A,
        \qquad
        u(e,\kxi(e))\ge u(a,\kxi(a))=:\underline u ,
\]
and the row of Table~\ref{tab:gate} certifies the $u$-gate.  Every strip
point with $e\ge\tfrac18$ has $\kappa\ge\kxi(e)$, so
Lemma~\ref{lem:zb-ugate}(b),(a) applies.
\end{proof}

\begin{table}[h]\centering
\caption{Gate rows.  $\overline A{=}x(a,\kappa_-)(1{+}1/\kappa_-)\uparrow$,
$\underline u{=}u(a,\kappa_-)\downarrow$ with $\kappa_-=\gamma(b)$ (part a)
resp.\ $\kappa_-=\kxi(a)$ (part b); slack
$=15(14-\overline A)\underline u-(\overline A+1)\downarrow$.  All entries
exact rationals (Convention~\ref{conv:zb-safe}).}\label{tab:gate}
\begin{tabular}{lcccc}
\toprule
$[a,b]$ & $\kappa_-$ & $\overline A\uparrow$ & $\underline u\downarrow$ & slack$\downarrow$\\
\midrule
\multicolumn{5}{l}{(a) curve $\kappa=\gamma(e)$ on $[\tfrac1{60},\tfrac18]$}\\
$[1/60,\,21/1000]$ & $259/1000$ & $4.6621$ & $0.0409$ & $0.06$ \\
$[21/1000,\,8/125]$ & $27/125$ & $5.3506$ & $0.0495$ & $0.07$ \\
$[8/125,\,1/8]$ & $31/200$ & $6.4352$ & $0.1364$ & $8.04$ \\
\midrule
\multicolumn{5}{l}{(b) curve $\kappa=\kxi(e)$ on $[\tfrac18,\tfrac{2033}{10000}]$}\\
$[1/8,\,2033/10000]$ & $8/49$ & $5.3477$ & $0.2494$ & $26.02$ \\
\bottomrule
\end{tabular}
\end{table}

\subsection{The two one-dimensional inequalities}

For a rational $b$ let $c_\ell(b)$, $c_s(b)$ denote the $5$-digit decimals
of Table~\ref{tab:surr} with
\begin{equation}\label{eq:surr}
        c_\ell^2\ \ge\ \frac{2b}{1-b},
        \qquad
        c_s^2\ \le\ 1-b
        \qquad(\text{exact rational comparisons}),
\end{equation}
so that, by Lemma~\ref{lem:zb-mono}(vi) and since
$\lbar(b)^{14}=\bigl(\tfrac{2b}{1-b}\bigr)^7$ is rational,
\begin{equation}\label{eq:Glo}
        G(e)\ \ge\ G(b)\ \ge\
        G_\downarrow(b):=c_s\,
        \frac{\bigl(1-c_\ell\bigr)\Bigl(1-\bigl(\tfrac{2b}{1-b}\bigr)^{7}\Bigr)}
        {1+c_\ell}
        \qquad\text{for } e\le b .
\end{equation}

\begin{proposition}[Case-S inequality]\label{prop:caseS}
For every $e\in[\tfrac1{60},\tfrac18]$,
\begin{equation}\label{eq:caseS}
        \tfrac34\,G(e)\ \ge\
        \frac{\Kstar\bigl(e,\gamma(e)\bigr)}{x\bigl(e,\gamma(e)\bigr)^2}
        +\widehat\Lambda .
\end{equation}
\end{proposition}

\begin{proof}
Fix the row $[a,b]$ of Table~\ref{tab:caseS} containing $e$ (the seven rows
tile $[\tfrac1{60},\tfrac18]$).  Since $\gamma(e)\le\gamma(a)\le
\gamma(\tfrac1{60})=\tfrac{79}{300}$ and by
Lemma~\ref{lem:zb-mono}(iv) (both parts, on the stated domains),
\[
        \Phi\bigl(e,\gamma(e)\bigr)\ \ge\ \Phi\bigl(e,\gamma(a)\bigr)
        \ \ge\ \Phi\bigl(a,\gamma(a)\bigr)=:\Phi_a ,
\]
so by Lemma~\ref{lem:zb-mono}(v) and $e^t\ge1+t+\tfrac{t^2}2+\tfrac{t^3}6$
($t\ge0$),
\[
        \Kstar\bigl(e,\gamma(e)\bigr)
        \ \le\ \exp\Bigl(-\sqrt{\Phi_a^2+4\Phi_a}\Bigr)
        \ \le\ e^{-r}
        \ \le\ \frac1{1+r+\frac{r^2}2+\frac{r^3}6}=:K^{\uparrow},
\]
where $r=r(a)$ is the rational surrogate of Table~\ref{tab:caseS} with
$r^2\le\Phi_a^2+4\Phi_a$ (exact comparison).  By
Lemma~\ref{lem:zb-mono}(viii), $x(e,\gamma(e))^2\ge x(b,\gamma(b))^2$, and by
\eqref{eq:Glo}, $G(e)\ge G_\downarrow(b)$.  The row certifies
\[
        \tfrac34\,G_\downarrow(b)\ \ge\
        \frac{K^{\uparrow}}{x(b,\gamma(b))^2}+\widehat\Lambda,
\]
which combined with the three monotone bounds yields \eqref{eq:caseS}.
\end{proof}

\begin{table}[h]\centering
\caption{Case-S rows.  $r=r(a)$ satisfies $r^2\le\Phi_a(\Phi_a+4)$,
$\Phi_a=\Phi(a,\gamma(a))$;
$K^\uparrow=1/(1{+}r{+}\tfrac{r^2}2{+}\tfrac{r^3}6)\uparrow$;
$X_\downarrow=x(b,\gamma(b))^2\downarrow$;
margin $=\tfrac34G_\downarrow(b)-K^\uparrow/X_\downarrow-\widehat\Lambda
\downarrow$.}\label{tab:caseS}
\begin{tabular}{lccccc}
\toprule
$[a,b]$ & $r$ & $K^\uparrow$ & $X_\downarrow$ & $G_\downarrow(b)\downarrow$ & margin$\downarrow$\\
\midrule
$[1/60,\,23/1000]$ & $0.88959$ & $0.41622$ & $0.89136$ & $0.63595$ & $0.0100$ \\
$[23/1000,\,33/1000]$ & $1.04817$ & $0.35850$ & $0.84928$ & $0.57596$ & $0.0098$ \\
$[33/1000,\,6/125]$ & $1.26209$ & $0.29468$ & $0.79119$ & $0.50537$ & $0.0065$ \\
$[6/125,\,67/1000]$ & $1.53635$ & $0.23144$ & $0.72527$ & $0.43499$ & $0.0071$ \\
$[67/1000,\,89/1000]$ & $1.84283$ & $0.17909$ & $0.65808$ & $0.36930$ & $0.0048$ \\
$[89/1000,\,111/1000]$ & $2.17522$ & $0.13781$ & $0.59913$ & $0.31450$ & $0.0058$ \\
$[111/1000,\,1/8]$ & $2.51056$ & $0.10754$ & $0.56532$ & $0.28369$ & $0.0225$ \\
\bottomrule
\end{tabular}
\end{table}

\begin{proposition}[Case-L inequalities]\label{prop:caseL}
\begin{enumerate}[label=\textup{(\alph*)},itemsep=1pt]
\item For every $e\in[\tfrac1{60},\tfrac18]$,
\begin{equation}\label{eq:caseLg}
        G(e)\Bigl(1-\frac{\kxi(e)}{4\gamma(e)}\Bigr)
        \ \ge\ \tfrac{14}{15}\,x\bigl(e,\gamma(e)\bigr)^{12}+\widehat\Lambda .
\end{equation}
\item For every $e\in[\tfrac18,\tfrac{2033}{10000}]$,
\begin{equation}\label{eq:caseLx}
        \tfrac34\,G(e)
        \ \ge\ \tfrac{14}{15}\,x\bigl(e,\kxi(e)\bigr)^{12}+\widehat\Lambda .
\end{equation}
\end{enumerate}
\end{proposition}

\begin{proof}
(a) Fix the row $[a,b]$ of Table~\ref{tab:caseL}(a) containing $e$.  By
Lemma~\ref{lem:zb-mono}(vi),(vii) the left side of \eqref{eq:caseLg} is a
product of two positive decreasing functions of $e$ (positivity of the
second factor: $\kxi/(4\gamma)\le\tfrac{400}{1519}<1$ on
$[\tfrac1{60},\tfrac18]$ by (vii)), so it is at least
$G_\downarrow(b)\bigl(1-\kxi(b)/(4\gamma(b))\bigr)=:L_\downarrow$; by
Lemma~\ref{lem:zb-mono}(viii) the right side is at most
$\tfrac{14}{15}x(a,\gamma(a))^{12}+\widehat\Lambda=:R_\uparrow$ (the power
$x(a,\gamma(a))^{12}$ is an exact rational).  The row certifies
$L_\downarrow\ge R_\uparrow$.

(b) Same with the single row of Table~\ref{tab:caseL}(b), using
Lemma~\ref{lem:zb-mono}(vi) for the left side and (ix) for the right.
\end{proof}

\begin{table}[h]\centering
\caption{Case-L rows.  (a): $L_\downarrow=G_\downarrow(b)\bigl(1-
\kxi(b)/(4\gamma(b))\bigr)\downarrow$,
$R_\uparrow=\tfrac{14}{15}x(a,\gamma(a))^{12}\uparrow$.
(b): $L_\downarrow=\tfrac34G_\downarrow(b)\downarrow$,
$R_\uparrow=\tfrac{14}{15}x(a,\kxi(a))^{12}\uparrow$.
Margin $=L_\downarrow-R_\uparrow-\widehat\Lambda\downarrow$.}\label{tab:caseL}
\begin{tabular}{lccc}
\toprule
$[a,b]$ & $L_\downarrow$ & $R_\uparrow$ & margin$\downarrow$\\
\midrule
\multicolumn{4}{l}{(a) curve $\kappa=\gamma(e)$ on $[\tfrac1{60},\tfrac18]$}\\
$[1/60,\,29/1000]$ & $0.57998$ & $0.56429$ & $0.0156$ \\
$[29/1000,\,13/200]$ & $0.40350$ & $0.39307$ & $0.0104$ \\
$[13/200,\,1/8]$ & $0.20898$ & $0.14341$ & $0.0655$ \\
\midrule
\multicolumn{4}{l}{(b) curve $\kappa=\kxi(e)$ on $[\tfrac18,\tfrac{2033}{10000}]$}\\
$[1/8,\,2033/10000]$ & $0.11051$ & $0.02983$ & $0.0806$ \\
\bottomrule
\end{tabular}
\end{table}

\begin{table}[h]\centering
\caption{Surrogates for \eqref{eq:surr} (exact $5$-digit decimals;
each satisfies the displayed exact rational comparison).}\label{tab:surr}
\begin{tabular}{lcc@{\qquad}lcc}
\toprule
$b$ & $c_\ell$ & $c_s$ & $b$ & $c_\ell$ & $c_s$\\
\midrule
$23/1000$ & $0.21699$ & $0.98843$ & $89/1000$ & $0.44203$ & $0.95446$\\
$29/1000$ & $0.24441$ & $0.98539$ & $111/1000$ & $0.49972$ & $0.94286$\\
$33/1000$ & $0.26126$ & $0.98336$ & $1/8$ & $0.53453$ & $0.93541$\\
$6/125$ & $0.31756$ & $0.97570$ & $2033/10000$ & $0.71440$ & $0.89258$\\
$13/200$ & $0.37288$ & $0.96695$ & & &\\
$67/1000$ & $0.37898$ & $0.96591$ & & &\\
\bottomrule
\end{tabular}
\end{table}

\subsection{Proof of the battle}

\begin{proof}[Proof of \eqref{eq:zb-battle} \textup{(Lemma
\texttt{lem:zoneB-battle} of \RII)}]
Let $(e,\kappa)$ be a strip point.  By Lemma~\ref{lem:zb-lambda},
$\overline\Lambda\le\widehat\Lambda$ throughout, so it suffices to prove
$\underline\Pi\ge\Kfrak/x^2+\widehat\Lambda$.  We distinguish three cases.

\medskip\noindent
\emph{Case S: $e\le\tfrac18$ and $\kappa\le\gamma(e)$.}
By Lemma~\ref{lem:zb-envelope}(a),(b),
\[
        \underline\Pi\ \ge\ \tfrac34G(e),
        \qquad
        \frac{\Kfrak}{x^2}\ \le\ \frac{\Kstar(e,\kappa)}{x(e,\kappa)^2}.
\]
Since $\kappa\le\gamma(e)\le\tfrac{79}{300}<1$, Lemma~\ref{lem:zb-mono}(iv),(v)
give $\Kstar(e,\kappa)\le\Kstar(e,\gamma(e))$, and Lemma~\ref{lem:zb-mono}(i)
gives $x(e,\kappa)\ge x(e,\gamma(e))$; hence
$\Kfrak/x^2\le\Kstar(e,\gamma(e))/x(e,\gamma(e))^2$, and
Proposition~\ref{prop:caseS} closes this case.

\medskip\noindent
\emph{Case L-$\gamma$: $e\le\tfrac18$ and $\kappa\ge\gamma(e)$.}
By Proposition~\ref{prop:zb-gate}(a) the gate holds at $(e,\kappa)$, so
Lemma~\ref{lem:zb-envelope}(c) and Lemma~\ref{lem:zb-mono}(i) give
\[
        \frac{\Kfrak}{x^2}\ \le\ \tfrac{14}{15}\,x(e,\kappa)^{12}
        \ \le\ \tfrac{14}{15}\,x\bigl(e,\gamma(e)\bigr)^{12}.
\]
By Lemma~\ref{lem:zb-envelope}(a) and $\kappa\ge\gamma(e)$,
\[
        \underline\Pi\ \ge\ G(e)\Bigl(1-\frac{\kxi(e)}{4\kappa}\Bigr)
        \ \ge\ G(e)\Bigl(1-\frac{\kxi(e)}{4\gamma(e)}\Bigr),
\]
and Proposition~\ref{prop:caseL}(a) closes this case.

\medskip\noindent
\emph{Case L-$\xi$: $e\ge\tfrac18$.}
Here $\kappa\ge\kxi(e)$, so by Proposition~\ref{prop:zb-gate}(b) the gate
holds; as above
$\Kfrak/x^2\le\tfrac{14}{15}x(e,\kappa)^{12}
\le\tfrac{14}{15}x(e,\kxi(e))^{12}$, while
$\underline\Pi\ge\tfrac34G(e)$ by Lemma~\ref{lem:zb-envelope}(a).
Proposition~\ref{prop:caseL}(b) closes this case.

\medskip
The three cases cover the strip: for $e\le\tfrac18$ every $\kappa$
satisfies $\kappa\le\gamma(e)$ or $\kappa\ge\gamma(e)$, and $e\ge\tfrac18$
is Case L-$\xi$.
\end{proof}

\subsection{Remarks}

\begin{remark}[the $\underline\rho$-term is not needed]\label{rmk:rho}
The proof used only $\underline\rho\ge0$
(Lemma~\ref{lem:zb-envelope}(a)).  Hence the battle \eqref{eq:zb-battle} holds
with $\epsbar$ replaced by the larger
$\min\bigl(\tfrac14,\kxi(e)/(4\kappa)\bigr)$, i.e.\ the
$17.37$-exponential can be deleted from the statement of
Lemma~\texttt{lem:zoneB-battle}, and with it the rational comparison
$28(1+\kappa)e/(1+3e)\ge17.37(1+\kappa)e$ in the proof of
Lemma~\texttt{lem:zoneB} of \RII{} becomes unnecessary (only
$\rho\ge0$ is used there).
\end{remark}

\begin{remark}[the $\max$ is redundant]
By Lemma~\ref{lem:zb-envelope}(b), on the gate-failure branch
$\Kfrak=\Kstar$; the paper may state the dichotomy as: gate
$\Rightarrow\Kfrak=x^{14}(14-A)_+/15$, otherwise $\Kfrak=\Kstar$.
\end{remark}

\begin{remark}[margins]
The battle \eqref{eq:zb-battle} in fact holds with pointwise margin
$\ge0.116$ on the whole strip (minimum at the pinch corner
$e\approx0.2029$, $\kappa=\kxi(e)$, where $\kxi$ meets $\kappa_q$); the
analytic chain above retains a margin $\ge0.0048$ in the tightest table
row (Case S, $[67/1000,89/1000]$) and $\ge0.06$ in the tightest gate row.
The proof uses $15$ table rows in total.
\end{remark}

\begin{remark}[verification protocol]
\texttt{validate\_zoneB.py} checks, and exits $0$ only if all pass:
(1) exact rational verification of every row of
Tables~\ref{tab:gate}--\ref{tab:surr}, including every surrogate comparison
\eqref{eq:surr} and $r^2\le\Phi_a(\Phi_a+4)$;
(2) the exact rational comparisons quoted in Lemma~\ref{lem:zb-lambda} and
Lemma~\ref{lem:zb-mono}(iv),(vii);
(3) every monotonicity claim of Lemma~\ref{lem:zb-mono} on a dense grid over
its exact stated domain;
(4) an end-to-end pointwise replay of the three-case assembly on a
$1201\times1201$ grid over the strip (each intermediate inequality
checked pointwise, including $\Kfrak\le\Kstar$ and the direct battle
margin);
(5) exact-\texttt{Fraction} spot checks of \eqref{eq:zb-battle} at $25$
corner/extreme rational points via safe-direction brackets, independent of
the tables.
\end{remark}

\section{Zone C without boxes}\label{sec:zoneC}

\subsection*{Provenance and verification status}
Proved by Claude on 2026-07-17 (v2 after one repair round: three
unsafe-rounding links in the wide-corner comparison chains, found by the
adversarial reviewer and the independent verifier, were re-rounded and
re-certified); final verdicts: verifier \emph{pass} (10-part script,
$\approx1.1$M comparisons, zero failures), reviewer \emph{sound}.  The
theorem proves Lemma~\emph{zoneC-mod} of [R2] analytically on its full
domain, eliminating the $3001$-box certificate
\texttt{zoneC\_certifier.py} entirely; Lemma~\emph{turan} of [R2] becomes
the special case $\delta<\tfrac1{2000}$ of the widened corner
proposition.  The only finite checks left are $21$ displayed
hand-checkable table rows.  The content is the note
\texttt{claude\_jul17/proof\_zoneC\_analytic\_v2.tex}, inlined verbatim.

\subsection{Statement, set-up, and what is being replaced}\label{sec:zoneC-setup}

This note refers throughout to the Region-II manuscript
(\texttt{paper\_region2\_v2.tex}); results quoted from it are cited by their
labels.  We work on the Region-II domain
$\D=\{\,\tfrac13<q<\tfrac12,\ q<\alpha\le r(q),\ m\ \text{odd}\,\}$ in the
$(e,\kappa)$ chart of that paper:
\[
 \alpha=\tfrac{1-e}2,\quad d=\alpha-q\in(0,\delta),\quad
 \delta=\tfrac{1-3e}6,\quad \kappa=\tfrac de,\quad p=1-q,\quad
 \xi=\tfrac{4\alpha^2d}{e^2},
\]
\[
 L=\sqrt{pq-\alpha^2},\quad x=\tfrac\alpha p,\quad s=\tfrac qp,\quad
 y=\tfrac Lp,\quad \ell=\tfrac L\alpha,\quad f=\alpha-L,\quad
 R_m=\alpha^m+L^m-pq^{m-1}.
\]
Note $q>\tfrac13\iff d<\delta$ and $\xi\le1\iff d\le e^2/(1-e)^2$.  The goal
is Lemma~\emph{zoneC-mod} of the paper: for every
$(q,\alpha,m)\in\D$ with $m\ge15$, $\xi\le1$ and $\tfrac1{60}\le e<\tfrac13$,
\begin{equation}\label{eq:zoneC-target}
 R_m\ \le\ C_m\,\psi(\xi,\rho),
\end{equation}
with the Huber payment $C_m\psi$ of Lemma~\emph{huber} of the paper.  In the
paper, \eqref{eq:zoneC-target} is proved on $e\le\tfrac13-\tfrac1{1000}$ by the
exact interval certificate \texttt{zoneC\_certifier.py} ($3001$ boxes,
subdivision depth $28$, cycle lengths checked up to $m=1716$; the prose count
``$2997$'' at the certificate paragraph is stale), plus the analytic
Tur\'an-corner sliver $\delta<\tfrac1{2000}$ (Lemma~\emph{turan}).  Here we
prove \eqref{eq:zoneC-target} analytically on the whole range:

\begin{theorem}\label{thm:zoneC-main}
Let $(q,\alpha,m)\in\D$ with $m\ge15$ odd, $\xi\le1$, and
$\tfrac1{60}\le e<\tfrac13$.  Then \eqref{eq:zoneC-target} holds, with the dual
certificate $\lambda=2\rho_{\mathrm{lo}}\xi$ if
$2\rho_{\mathrm{lo}}\xi\le1$ and $\lambda=1$ otherwise, exactly as in
Lemma~\emph{threegeo}(iii).  Consequently Lemma~\emph{zoneC-mod} holds with
no computer certificate, and Lemma~\emph{turan} is the special case
$\delta<\tfrac1{2000}$ of Proposition~\ref{prop:wide} below.
\end{theorem}

The proof splits at $e=\tfrac3{10}$ (i.e.\ $\delta=\tfrac1{60}$):
Proposition~\ref{prop:strip} covers $\tfrac1{60}\le e\le\tfrac3{10}$ by a
two-variable analysis whose only case analysis is a $21$-row table of
one-variable evaluations at rational endpoints; Proposition~\ref{prop:wide}
covers $\tfrac3{10}\le e<\tfrac13$ (all $0<d<\delta$, hence in particular
all $\xi\le1$) by a widened Tur\'an-corner chain.  Everything below uses only
Lemma~\emph{threegeo} of the paper and elementary calculus; every decimal
constant is an exact rational rounded in the safe direction, and every
displayed inequality --- including each individual link of the two
comparison chains in Step~4 of Proposition~\ref{prop:wide} --- is checked
by the companion script \texttt{validate\_zoneC\_analytic\_v2.py} (dense
float grids over the exact stated domains, exact \texttt{Fraction}
arithmetic at corners and for every chain link, table rows recomputed to
$50$ digits).

\subsubsection*{Standing notation}
For $m\ge15$ put $n=m-2\ge13$ and
\[
 \lambda_x=\ln\tfrac1x=\ln\tfrac p\alpha,\qquad
 \Delta=\ln\tfrac xs=\ln\tfrac\alpha q>0,\qquad
 \lambda_s=\ln\tfrac1s=\lambda_x+\Delta,
\]
\[
 G_2=\ell^2\bigl(1-(y/s)^{13}\bigr),\qquad G=\ln(1+G_2),\qquad
 n_g=\tfrac G\Delta,\qquad
 \mathrm{Def}_m=\tfrac{R_m}{\alpha^3p^{\,m-2}} .
\]
On $\D$ one has $0<y<s$ and $0\le\ell<1$ (Lemma~\emph{frontier-gap} of the
paper; equivalently $L<q\iff q>\tfrac13$), so $0<G_2<1+G_2$ and $G>0$.

\subsubsection*{Elementary identities}
All are one-line expansions, checked in exact rational arithmetic in Part~1
of the validation script.  With $S:=q^2-L^2$:
\begin{align}
 &L^2=\alpha e-d(e+d),\qquad pq-\alpha^2=L^2,\qquad
 \alpha-\delta=\tfrac13,\qquad \alpha-e=\tfrac{1-3e}2=3\delta,
 \label{eq:id1}\\
 &S=3\alpha\delta-(2\alpha-e)d+2d^2,\qquad
 L^2-e^2=\delta-\tfrac d3-6\delta^2+2\delta d-d^2,\label{eq:id2}\\
 &(1-\alpha)^2-4\alpha e=\tfrac{(1-3e)^2}4\ \ge0
 \quad\Longrightarrow\quad 4L^2\le4\alpha e\le(1-\alpha)^2\le p^2
 \quad\Longrightarrow\quad y\le\tfrac12.\label{eq:yhalf}
\end{align}

\subsection{The defect majorant and its closed-form supremum}\label{sec:zoneC-majorant}

\begin{lemma}[majorant]\label{lem:zc-majorant}
For every $n\ge13$,
\(
 \alpha\,\mathrm{Def}_m
 = x^n-s^n-\ell^2(s^n-y^n)
 \le \varphi(n):=x^n-(1+G_2)s^n .
\)
\end{lemma}

\begin{proof}
The identity is Lemma~\emph{threegeo}(i) (divide $R_m$ by $\alpha^2p^{m-2}$
and use $pq-\alpha^2=L^2$).  Since $0<y/s<1$ and $n\ge13$,
$\ell^2(s^n-y^n)=\ell^2 s^n(1-(y/s)^n)\ge s^nG_2$.
\end{proof}

\begin{lemma}[closed-form supremum]\label{lem:zc-sup}
Extend $\varphi(t)=x^t-(1+G_2)s^t$ to real $t$.  Then $\varphi$ is increasing
on $(-\infty,n_*]$ and decreasing on $[n_*,\infty)$, where
$n_*=\bigl(G+\ln(\lambda_s/\lambda_x)\bigr)/\Delta>n_g$, and
\begin{equation}\label{eq:phistar}
 \sup_{t\in\mathbb R}\varphi(t)=\varphi(n_*)
 =\frac{\Delta}{\lambda_s}\;h\;x^{\,n_g},
 \qquad
 h:=\exp\Bigl(-\frac{\lambda_x}{\Delta}\,
 \ln\bigl(1+\tfrac\Delta{\lambda_x}\bigr)\Bigr)\in(\eu^{-1},1).
\end{equation}
Moreover, with $v:=\Delta/\lambda_x$,
\begin{equation}\label{eq:hbound}
 h\ \le\ \exp\bigl(-1+\tfrac v2\bigr).
\end{equation}
\end{lemma}

\begin{proof}
$\varphi'(t)=x^t\bigl[-\lambda_x+(1+G_2)\lambda_s(s/x)^t\bigr]$ and
$(s/x)^t$ is strictly decreasing, so $\varphi'$ changes sign exactly once,
from $+$ to $-$, at the root $n_*$ of
$(x/s)^{n_*}=(1+G_2)\lambda_s/\lambda_x$; taking logarithms gives the stated
$n_*$, and $n_*>n_g=G/\Delta$ because $\lambda_s>\lambda_x$.  At $t=n_*$,
$(1+G_2)s^{n_*}=x^{n_*}\lambda_x/\lambda_s$, hence
$\varphi(n_*)=x^{n_*}(1-\lambda_x/\lambda_s)=x^{n_*}\Delta/\lambda_s$; and
$x^{n_*}=\exp(-\lambda_xn_*)
 =x^{n_g}\exp\bigl(-\tfrac{\lambda_x}\Delta\ln\tfrac{\lambda_s}{\lambda_x}\bigr)
 =x^{n_g}h$.
For \eqref{eq:hbound}: $-\ln h=\tfrac1v\ln(1+v)\ge\tfrac1v(v-\tfrac{v^2}2)
=1-\tfrac v2$, by $\ln(1+v)\ge v-v^2/2$ for $v\ge0$.
\end{proof}

\begin{lemma}[pointwise scalar facts]\label{lem:zc-scalar}
On $\D$:
\begin{enumerate}
\item[(i)] $d\ge q\Delta$ and $\kappa\ge q\Delta/e$
  \quad(from $d=q(\eu^\Delta-1)\ge q\Delta$);
\item[(ii)] $\Delta\le d/q$ \quad(from $\ln(1+d/q)\le d/q$);
\item[(iii)] $\lambda_x\ge\ln\frac{1+e}{1-e}\ge2e$
  \quad(from $p/\alpha=1+\frac{e+d}\alpha\ge1+\frac{2e}{1-e}=\frac{1+e}{1-e}$
  and $\ln\frac{1+e}{1-e}=2\operatorname{artanh}e\ge2e$);
\item[(iv)] $G_2\le\ell^2\le\frac e\alpha=\frac{2e}{1-e}$
  \quad(from $L^2\le\alpha e$ by \eqref{eq:id1});
\item[(v)] if $R_m>0$ and $m\ge15$, then $n=m-2>n_g$, i.e.\ $m>2+n_g$
  \quad(Lemma~\emph{threegeo}(ii): $x^n>(1+G_2)s^n$, take logarithms).
\end{enumerate}
\end{lemma}

\subsection{Payment floors and the three pointwise master inequalities}
\label{sec:zoneC-master}

Define the $m$-free payment coefficients
\[
 c_e:=\frac{2(1-x^{14})(1-y^{14})}{(1+x)(1+y)},
 \qquad
 K_2:=\frac{\sqrt{2\alpha}\,(1-y^{14})}{2\alpha^3(1+y)} .
\]

\begin{lemma}[payment floors]\label{lem:zc-payfloor}
Let $m\ge15$.  If $2\rho_{\mathrm{lo}}\xi\le1$ then
$\dfrac{C_m\psi(\xi,\rho)}{\alpha^3p^{m-2}}\ \ge\ m\,c_e\,\kappa^2$;
if $2\rho_{\mathrm{lo}}\xi>1$ then
$\dfrac{C_m\psi(\xi,\rho)}{\alpha^3p^{m-2}}\ \ge\ m\,K_2\,f\,d$.
\end{lemma}

\begin{proof}
Relax Lemma~\emph{threegeo}(iii): in the first branch use
$(1+4\xi)/(1+2\xi)\ge1$ and $1-y^{m-1}\ge1-y^{14}$ ($m\ge15$); in the second
use $(4\xi+1)/(4\xi+2)\ge\tfrac12$ and $1-y^{m-1}\ge1-y^{14}$.
\end{proof}

\begin{lemma}[pointwise master inequalities]\label{lem:zc-master}
Let $(q,\alpha,m)\in\D$, $m\ge15$ odd, $R_m>0$, and write $w:=\lambda_xG/\Delta$
(so $x^{n_g}=\eu^{-w}$).  Then \eqref{eq:zoneC-target} holds as soon as the
applicable one of the following scalar inequalities does:
\begin{align}
 &\textup{(deep, case I: }2\rho_{\mathrm{lo}}\xi\le1,\ n_g\ge15\textup{):}
 &h\,\eu^{-w}e^2\ &\le\ \lambda_x\,\alpha\,c_e\,q^2\,G,
 \label{eq:SCdeep}\\
 &\textup{(shallow, case I: }2\rho_{\mathrm{lo}}\xi\le1,\ n_g<15\textup{):}
 &h\,\eu^{-w}e^2\ &\le\ 15\,\lambda_x\,\alpha\,c_e\,q^2\,\Delta,
 \label{eq:SCsh}\\
 &\textup{(case II: }2\rho_{\mathrm{lo}}\xi>1\textup{):}
 &h\,\eu^{-w}\ &\le\ 15\,\lambda_x\,\alpha\,K_2\,f\,q.
 \label{eq:SCII}
\end{align}
\end{lemma}

\begin{proof}
By Lemmas~\ref{lem:zc-majorant} and~\ref{lem:zc-sup},
$\alpha\,\mathrm{Def}_m\le\varphi(n)\le\varphi(n_*)
=(\Delta/\lambda_s)h\,\eu^{-w}\le(\Delta/\lambda_x)h\,\eu^{-w}$.
It therefore suffices that
$(\Delta/\lambda_x)\,h\,\eu^{-w}\le\alpha\,m\cdot(\text{payment floor of
Lemma~\ref{lem:zc-payfloor}})$.

\emph{Deep.}  By Lemma~\ref{lem:zc-scalar}(v), $m>n_g$, so
$m\,c_e\kappa^2\ge n_g\,c_e\,q^2\Delta^2/e^2$ using
Lemma~\ref{lem:zc-scalar}(i).  Since $n_g\Delta^2=G\Delta$, the sufficient
condition becomes
$(\Delta/\lambda_x)h\eu^{-w}\le\alpha c_eq^2G\Delta/e^2$, which is
\eqref{eq:SCdeep}.

\emph{Shallow.}  Use $m\ge15$ and $\kappa\ge q\Delta/e$:
$(\Delta/\lambda_x)h\eu^{-w}\le15\alpha c_eq^2\Delta^2/e^2$ is
\eqref{eq:SCsh}.

\emph{Case II.}  Use $m\ge15$ and $d\ge q\Delta$:
$(\Delta/\lambda_x)h\eu^{-w}\le15\alpha K_2fq\Delta$ is \eqref{eq:SCII}.

Finally, if $R_m\le0$ then \eqref{eq:zoneC-target} is trivial because
$C_m\ge0$ and $\psi\ge0$.
\end{proof}

\subsection{Interval brackets and the $21$-row table
($\tfrac1{60}\le e\le\tfrac3{10}$)}\label{sec:zoneC-table}

Fix a closed interval $E=[e_L,e_R]\subseteq[\tfrac1{60},\tfrac3{10}]$ and let
$(q,\alpha,m)$ range over all points of $\D$ with $e\in E$, $m\ge15$ odd and
\begin{equation}\label{eq:dcap}
 d\ \le\ d_{\max}(e):=\min\Bigl(\frac{e^2}{(1-e)^2},\ \delta(e)\Bigr),
\end{equation}
which contains the Zone-C constraint set $\{\xi\le1,\ q>\tfrac13\}$.
Define the \emph{interval constants} (elementary evaluations at the rational
endpoints; $\alpha(u)=\tfrac{1-u}2$, $\delta(u)=\tfrac{1-3u}6$):
\[
 \bar q:=\max\Bigl(\tfrac13,\ \alpha(e_R)-\tfrac{e_R^2}{(1-e_R)^2}\Bigr),
 \qquad
 D:=\min\Bigl(\tfrac{e_R^2}{(1-e_R)^2},\ \delta(e_L)\Bigr),
 \qquad
 \Lambda:=\frac D{\bar q},
\]
\[
 \lambda_L:=\ln\frac{1+e_L}{1-e_L},\qquad
 \widehat G_2:=\frac{2e_R}{1-e_R},\qquad
 Y:=\frac{\sqrt{e_R(1-e_R)/2}}{\bar q},
\]
\[
 \gamma:=\Bigl[\frac2{1-e_L}-\frac{D(e_R+D)}{\alpha(e_R)^2e_L}\Bigr]
 \,(1-Y^{13})\,\frac{\ln(1+\widehat G_2)}{\widehat G_2},
 \qquad
 a:=\lambda_L\,\gamma\,e_L,
\]
\[
 w^-:=2\gamma\max\Bigl((1-e_R)^2\bar q,\ \frac{e_L^2}{3\delta(e_L)}\Bigr),
 \qquad
 v^+:=\min\Bigl(\frac{e_R}{2(1-e_R)^2\bar q},\ \frac{3\delta(e_L)}{\lambda_L}
 \Bigr),\qquad
 h^+:=\eu^{-1+v^+/2},
\]
\[
 x_U:=\frac{1-e_L}{1+e_L},\qquad
 y_U:=\frac{\sqrt{e_R(1-e_R)/2}}{\alpha(e_R)+e_R},\qquad
 c^-:=\frac{2(1-x_U^{14})(1-y_U^{14})}{(1+x_U)(1+y_U)},
\]
\[
 f^-:=\alpha(e_R)-\sqrt{e_R(1-e_R)/2},\qquad
 K^-:=\frac{\sqrt{2\alpha(e_R)}\,(1-y_U^{14})}{2\,\alpha(e_L)^3\,(1+y_U)}.
\]

\begin{lemma}[brackets]\label{lem:zc-bracket}
For every point as above with $e\in E$:
\begin{enumerate}
\item[(a)] $q\ge\bar q$;\ \ $d\le d_{\max}(e)\le D$;\ \
  $\Delta\le d/q\le\Lambda$;\ \ $\lambda_x\ge\max(\lambda_L,2e)$.
\item[(b)] $G\ge\gamma\,e>0$, and hence
  $\lambda_xG\ge a$ and $\lambda_xG\ge2\gamma e^2$.
\item[(c)] $w=\lambda_xG/\Delta\ \ge\ w^-$; if moreover $n_g\ge15$ then also
  $w\ge15\lambda_L$.
\item[(d)] $v=\Delta/\lambda_x\le v^+$ and $h\le h^+$.
\item[(e)] $c_e\ge c^-$,\quad $f\ge f^-$,\quad $K_2\ge K^-$.
\item[(f)] if $n_g<15$ at some point of $E$, then $\Lambda>\gamma e_L/15$.
\end{enumerate}
\end{lemma}

\begin{proof}
(a) $q=\alpha(e)-d\ge\alpha(e)-d_{\max}(e)
=\max\bigl(\tfrac13,\alpha(e)-\tfrac{e^2}{(1-e)^2}\bigr)$ (using
$\alpha-\delta=\tfrac13$), and both expressions in the maximum are
nonincreasing in $e$, so $q\ge\bar q$.  $d_{\max}$ is a minimum of an
increasing and a decreasing function of $e$, whence $d_{\max}(e)\le D$.
$\Delta\le d/q$ is Lemma~\ref{lem:zc-scalar}(ii); then $\Delta\le D/\bar q$.
$\lambda_x\ge\ln\frac{1+e}{1-e}\ge\lambda_L$ (monotone) and $\ge2e$ by
Lemma~\ref{lem:zc-scalar}(iii).

(b) Write $G_2=\ell^2(1-(y/s)^{13})$.  First,
$\ell^2/e=\frac1\alpha-\frac{d(e+d)}{\alpha^2e}
\ge\frac2{1-e_L}-\frac{D(e_R+D)}{\alpha(e_R)^2e_L}$,
bounding each monotone atom at its worst endpoint
($\tfrac1\alpha=\tfrac2{1-e}\ge\tfrac2{1-e_L}$; $d\le D$; $e+d\le e_R+D$;
$\alpha\ge\alpha(e_R)$; $e\ge e_L$).  Second,
$y/s=L/q\le\sqrt{\alpha e}/\bar q\le Y$ since
$\alpha e=\tfrac{e(1-e)}2$ is increasing on $e\le\tfrac12$; hence
$1-(y/s)^{13}\ge1-Y^{13}$.  Third, by concavity of $\ln(1+\cdot)$ and
Lemma~\ref{lem:zc-scalar}(iv) ($G_2\le\widehat G_2$),
$G=\ln(1+G_2)\ge G_2\,\ln(1+\widehat G_2)/\widehat G_2$.  Multiplying the
three bounds gives $G\ge\gamma e$.  Then
$\lambda_xG\ge\lambda_L\cdot\gamma e_L=a$ and
$\lambda_xG\ge2e\cdot\gamma e$.

(c) $w=\lambda_xG/\Delta\ge2\gamma e^2/\Delta$.  Bound $\Delta$ twice:
$\Delta\le\frac{e^2}{(1-e)^2q}$ gives
$w\ge2\gamma(1-e)^2q\ge2\gamma(1-e_R)^2\bar q$;
$\Delta\le d/q\le\delta(e)\cdot3$ gives
$w\ge\frac{2\gamma e^2}{3\delta(e)}\ge\frac{2\gamma e_L^2}{3\delta(e_L)}$
because $e^2/\delta(e)$ is increasing.  Hence $w\ge w^-$.  If $n_g\ge15$,
$w=\lambda_xn_g\ge15\lambda_x\ge15\lambda_L$.

(d) As in (c), $v=\Delta/\lambda_x
\le\min\bigl(\frac{e^2}{(1-e)^2q}\cdot\frac1{2e},\
\frac{3\delta(e_L)}{\lambda_L}\bigr)$, and
$\frac{e}{2(1-e)^2q}\le\frac{e_R}{2(1-e_R)^2\bar q}$ (increasing numerator
atoms, decreasing $q$).  Then $h\le\eu^{-1+v/2}\le h^+$ by
\eqref{eq:hbound}.

(e) $x=\alpha/p\le\alpha/(\alpha+e)=\frac{1-e}{1+e}\le x_U$ and
$t\mapsto(1-t^{14})/(1+t)$ is decreasing on $(0,1)$, so
$(1-x^{14})/(1+x)\ge(1-x_U^{14})/(1+x_U)$.  Similarly
$y\le\sqrt{\alpha e}/(\alpha+e)=:y_{\mathrm{up}}(e)\le y_U$
($y_{\mathrm{up}}^2=\frac{2e(1-e)}{(1+e)^2}$ is increasing for
$e<\tfrac13$), giving $(1-y^{14})/(1+y)\ge(1-y_U^{14})/(1+y_U)$; multiply
for $c_e\ge c^-$.  Next $f=\alpha-L\ge\alpha(e)-\sqrt{\alpha(e)e}$, which is
decreasing in $e$, so $f\ge f^-$.  Finally
$\sqrt{2\alpha}\ge\sqrt{2\alpha(e_R)}$, $\alpha^3\le\alpha(e_L)^3$, and the
$y$-factor as before give $K_2\ge K^-$.

(f) $n_g<15$ means $\Delta>G/15\ge\gamma e/15\ge\gamma e_L/15$; combined
with $\Delta\le\Lambda$ from (a), such a point can exist only if
$\Lambda>\gamma e_L/15$.
\end{proof}

\begin{lemma}[assembled row bounds]\label{lem:zc-rows}
With the constants above, define
\[
 \Psi_{\mathrm{deep}}(E):=
 \frac{h^+\exp\bigl(-\max(w^-,\,15\lambda_L)\bigr)}
      {2\gamma\,\alpha(e_R)\,c^-\,\bar q^{\,2}},
 \qquad
 \Psi_{\mathrm{II}}(E):=
 \frac{h^+\,\eu^{-a/\Lambda}}
      {15\,\lambda_L\,\alpha(e_R)\,K^-\,f^-\,\bar q},
\]
and, when $\Lambda>\gamma e_L/15$ (otherwise the shallow case is empty
on $E$ by Lemma~\ref{lem:zc-bracket}(f)),
\[
 \Psi_{\mathrm{sh}}(E):=
 \frac{h^+\,(e_R/2)\,\eu^{-a/\Delta^*}}
      {15\,\alpha(e_R)\,c^-\,\bar q^{\,2}\,\Delta^*},
 \qquad
 \Delta^*:=\min\Bigl(\Lambda,\ \max\bigl(a,\ \tfrac{\gamma e_L}{15}\bigr)\Bigr).
\]
If $\Psi_{\mathrm{deep}}(E)\le1$, $\Psi_{\mathrm{II}}(E)\le1$, and (when
applicable) $\Psi_{\mathrm{sh}}(E)\le1$, then \eqref{eq:zoneC-target} holds at
every point of $\D$ with $e\in E$, \eqref{eq:dcap}, and $m\ge15$ odd.
\end{lemma}

\begin{proof}
Assume $R_m>0$ (else trivial) and verify the applicable master inequality of
Lemma~\ref{lem:zc-master}.

\emph{Deep} \eqref{eq:SCdeep}: by Lemma~\ref{lem:zc-bracket},
$h\eu^{-w}e^2\le h^+\exp(-\max(w^-,15\lambda_L))\,e^2$ and
$\lambda_x\alpha c_eq^2G
\ge(\lambda_xG)\,\alpha(e_R)c^-\bar q^2
\ge2\gamma e^2\,\alpha(e_R)c^-\bar q^2$; the $e^2$ cancels and the claim is
$\Psi_{\mathrm{deep}}(E)\le1$.

\emph{Shallow} \eqref{eq:SCsh}: divide \eqref{eq:SCsh} by
$\lambda_x\Delta$ and use $h\le h^+$,
$\eu^{-w}=\eu^{-\lambda_xG/\Delta}\le\eu^{-a/\Delta}$
(Lemma~\ref{lem:zc-bracket}(b)), $e^2/\lambda_x\le e/2\le e_R/2$
(Lemma~\ref{lem:zc-scalar}(iii)), $\alpha\ge\alpha(e_R)$, $c_e\ge c^-$ and
$q\ge\bar q$: \eqref{eq:SCsh} follows from
\[
 h^+\,\frac{e_R}2\,\frac{\eu^{-a/\Delta}}{\Delta}
 \ \le\ 15\,\alpha(e_R)\,c^-\,\bar q^{\,2}.
\]
The function $\Delta\mapsto\eu^{-a/\Delta}/\Delta$ increases on $(0,a]$ and
decreases on $[a,\infty)$; the shallow case confines $\Delta$ to
$(\gamma e_L/15,\ \Lambda]$ (Lemma~\ref{lem:zc-bracket}(a),(f) and
$\Delta>G/15\ge\gamma e_L/15$), on which its supremum is attained at
$\Delta^*$.  The claim is $\Psi_{\mathrm{sh}}(E)\le1$.

\emph{Case II} \eqref{eq:SCII}: $h\eu^{-w}\le h^+\eu^{-a/\Delta}\le
h^+\eu^{-a/\Lambda}$ ($\Delta\le\Lambda$, and $\eu^{-a/\Delta}$ is
increasing in $\Delta$), while
$15\lambda_x\alpha K_2fq\ge15\lambda_L\alpha(e_R)K^-f^-\bar q$; the claim is
$\Psi_{\mathrm{II}}(E)\le1$.
\end{proof}

\begin{proposition}[the strip $\tfrac1{60}\le e\le\tfrac3{10}$]
\label{prop:strip}
For every $(q,\alpha,m)\in\D$ with $m\ge15$ odd,
$e\in[\tfrac1{60},\tfrac3{10}]$ and $d\le d_{\max}(e)$ (in particular for
all of Zone C there), \eqref{eq:zoneC-target} holds.
\end{proposition}

\begin{proof}
Apply Lemma~\ref{lem:zc-rows} to the $21$ intervals of
Table~\ref{tab:rows}, whose union is $[\tfrac1{60},\tfrac3{10}]$.  Each
table entry is an upper bound (rounded up in the last digit) for the
corresponding $\Psi(E)$, an explicit elementary expression in the rational
endpoints; all entries are $<1$.  The entries are recomputed with $50$-digit
arithmetic in Part~4 of the validation script, and every bracket of
Lemma~\ref{lem:zc-bracket} is checked on dense $(e,\kappa,m)$ grids in Part~3.
In the first four rows $\Lambda\le\gamma e_L/15$, so the shallow case is
empty there by Lemma~\ref{lem:zc-bracket}(f).
\end{proof}

\begin{table}[h]
\centering\small
\caption{Row bounds for Lemma~\ref{lem:zc-rows}.  ``--'' means
$\Lambda\le\gamma e_L/15$ (shallow case empty).  $\gamma$ is rounded down;
$\Psi$ entries are rounded up.}\label{tab:rows}
\begin{tabular}{llcccc}
\toprule
$e_L$ & $e_R$ & $\gamma\ge$ & $\Psi_{\mathrm{deep}}\le$ &
$\Psi_{\mathrm{sh}}\le$ & $\Psi_{\mathrm{II}}\le$\\
\midrule
$1/60$ & $1/48$ & $1.9894$ & $0.391$ & -- & $0.698$\\
$1/48$ & $1/40$ & $1.9885$ & $0.348$ & -- & $0.537$\\
$1/40$ & $1/32$ & $1.9822$ & $0.333$ & -- & $0.545$\\
$1/32$ & $1/25$ & $1.9733$ & $0.322$ & -- & $0.521$\\
$1/25$ & $1/20$ & $1.9654$ & $0.315$ & $0.643$ & $0.441$\\
$1/20$ & $1/16$ & $1.9508$ & $0.331$ & $0.497$ & $0.416$\\
$1/16$ & $1/13$ & $1.9326$ & $0.248$ & $0.390$ & $0.388$\\
$1/13$ & $1/10$ & $1.8759$ & $0.177$ & $0.368$ & $0.477$\\
$1/10$ & $1/8$  & $1.8173$ & $0.100$ & $0.308$ & $0.476$\\
$1/8$  & $3/20$ & $1.7382$ & $0.057$ & $0.284$ & $0.506$\\
$3/20$ & $17/100$ & $1.6632$ & $0.032$ & $0.268$ & $0.520$\\
$17/100$ & $9/50$ & $1.6382$ & $0.019$ & $0.241$ & $0.472$\\
$9/50$ & $19/100$ & $1.5462$ & $0.017$ & $0.263$ & $0.541$\\
$19/100$ & $1/5$ & $1.4072$ & $0.015$ & $0.302$ & $0.638$\\
$1/5$ & $21/100$ & $1.3018$ & $0.013$ & $0.324$ & $0.690$\\
$21/100$ & $11/50$ & $1.2826$ & $0.009$ & $0.305$ & $0.630$\\
$11/50$ & $6/25$ & $1.1403$ & $0.008$ & $0.338$ & $0.703$\\
$6/25$ & $13/50$ & $1.0113$ & $0.005$ & $0.335$ & $0.639$\\
$13/50$ & $7/25$ & $0.8181$ & $0.004$ & $0.368$ & $0.640$\\
$7/25$ & $29/100$ & $0.7254$ & $0.002$ & $0.329$ & $0.436$\\
$29/100$ & $3/10$ & $0.5823$ & $0.002$ & $0.379$ & $0.487$\\
\bottomrule
\end{tabular}
\end{table}

\subsection{The widened Tur\'an corner ($\tfrac3{10}\le e<\tfrac13$)}
\label{sec:zoneC-wide}

Here $\delta=\tfrac{1-3e}6\in(0,\tfrac1{60}]$, $\alpha=\tfrac13+\delta$,
$e=\tfrac13-2\delta$, and $d\in(0,\delta)$ is unrestricted (no use of
$\xi\le1$ is made).  All constants below are for this range; each is checked
on dense $(\delta,d,m)$ grids and, where rational, exactly, in Part~5 of the
validation script.

\begin{lemma}[corner geometry]\label{lem:zc-corner-geo}
For $0<\delta\le\tfrac1{60}$ and $0<d<\delta$:
\begin{enumerate}
\item[(C1)] $\tfrac23\delta\le S\le(1+3\delta)\delta\le1.05\,\delta$, where
 $S=q^2-L^2=3\alpha\delta-(2\alpha-e)d+2d^2$ by \eqref{eq:id2}.
 \emph{Proof:} $S/\delta=1-\tfrac t3+\delta(3-4t+2t^2)$ with
 $t=d/\delta\in(0,1)$, and $1\le3-4t+2t^2\le3$ on $[0,1]$.
\item[(C2)] $e<L\le\sqrt{\alpha e}\le\tfrac13-\tfrac\delta2$.
 \emph{Proof:} $L^2-e^2$ in \eqref{eq:id2} is decreasing in $d$; at
 $d=\delta$ it equals $\tfrac23\delta-5\delta^2>0$.  And
 $\alpha e=\tfrac19-\tfrac\delta3-2\delta^2
 \le(\tfrac13-\tfrac\delta2)^2$.
\item[(C3)] $\tfrac{19}{30}\le\tfrac23-2\delta\le q+L\le\tfrac23
 +\tfrac\delta2\le0.675$;\quad
 $\alpha+L\le0.35+0.3415651\le0.6915651$
 (using $\sqrt{\alpha e}\le\sqrt{0.35/3}\le0.3415651$);\quad
 $x\le\tfrac\alpha{1-\alpha}\le\tfrac7{13}$;\quad $y\le\tfrac12$
 by \eqref{eq:yhalf}.
\item[(C4)] $f=\dfrac{3\alpha\delta+d(e+d)}{\alpha+L}$
 (identity, from $\alpha^2-L^2=3\alpha\delta+d(e+d)$), and
 \[
  1.4459\,\delta\le\frac{3\alpha\delta}{0.6915651}\le f
  \le\frac{(1.05+\tfrac13)\delta}{2/3-\delta}\le
  \frac{1.383334\,\delta}{0.65}\le2.1283\,\delta,
 \]
 whence $1-\ell=f/\alpha\le3f\le6.3849\,\delta$ and
 $\ell^2\ge1-12.7698\,\delta\ge0.78717$.
\item[(C5)] $t_0:=1-\tfrac ys=\tfrac{q-L}q=\tfrac S{(q+L)q}$ satisfies
 \[
  2.8218\,\delta\le\frac{(2/3)\delta}{0.675\cdot0.35}\le t_0\le
  \frac{1.05\,\delta}{(19/30)(1/3)}\le4.9737\,\delta .
 \]
\item[(C6)] $G_2\ge18.09\,\delta$ and $G_2\le13\,t_0\le64.6581\,\delta
 \le1.0777$; consequently $\ln(1+G_2)\ge12.27\,\delta$.
 \emph{Proof:} $t\mapsto1-(1-t)^{13}$ is concave, so on
 $[0,T]$ with $T:=4.9737/60$ it dominates its chord:
 $1-(1-t_0)^{13}\ge t_0\,(1-(1-T)^{13})/T\ge8.146\,t_0$.  Hence
 $G_2=\ell^2(1-(1-t_0)^{13})\ge0.78717\cdot8.146\cdot2.8218\,\delta
 \ge18.09\,\delta$.  The upper bound is Bernoulli
 ($1-(1-t_0)^{13}\le13t_0$, $\ell\le1$).  Concavity of $\ln(1+\cdot)$ on
 $[0,1.0777]$ gives
 $\ln(1+G_2)\ge G_2\,\tfrac{\ln2.0777}{1.0777}\ge0.6785\,G_2
 \ge12.27\,\delta$.
\end{enumerate}
\end{lemma}

\begin{proposition}[wide Tur\'an corner]\label{prop:wide}
For every $(q,\alpha,m)\in\D$ with $m\ge15$ odd and
$0<\delta\le\tfrac1{60}$ (i.e.\ $\tfrac3{10}\le e<\tfrac13$),
\eqref{eq:zoneC-target} holds.
\end{proposition}

\begin{proof}
Assume $R_m>0$ (else trivial).

\emph{Step 1 (forcing).}  Lemma~\emph{threegeo}(ii) with
$\ln(\alpha/q)\le d/q$ and $q>\tfrac13$ gives, by (C6),
\begin{equation}\label{eq:wide-forcing}
 (m-2)\,d\ >\ q\,\ln(1+G_2)\ \ge\ \tfrac13\cdot12.27\,\delta\ =\ 4.09\,\delta .
\end{equation}

\emph{Step 2 (defect).}  With $n=m-2$, the mean value theorem on
$[s,x]$ and $[y,s]$ in Lemma~\ref{lem:zc-majorant}'s identity, together with
$y=\ell x$, $x-s=d/p$, $s-y=(q-L)/p$ and $q-L=S/(q+L)$, gives
\[
 \alpha\,\mathrm{Def}_m\ \le\ \frac{n\,x^{n-1}}p\,N,
 \qquad
 N:=d-\ell^{\,m-1}\frac S{q+L}.
\]
Insert $S=3\alpha\delta-(2\alpha-e)d+2d^2$, drop the term
$-2d^2\ell^{m-1}/(q+L)\le0$, use $\ell^{m-1}\le1$ on the positive $d$-term
and Bernoulli $\ell^{m-1}\ge1-(m-1)(1-\ell)$ on the $3\alpha\delta$-term:
\[
 N\le d\Bigl(1+\frac{2\alpha-e}{q+L}\Bigr)-\frac{3\alpha\delta}{q+L}
 +\frac{3\alpha\delta\,(m-1)(1-\ell)}{q+L}.
\]
Since $d\le\delta$ and $q-\alpha=-d$, the first two terms combine to
\[
 \frac{d(q+L+2\alpha-e)-3\alpha\delta}{q+L}
 \le\delta\,\frac{q+L-\alpha-e}{q+L}
 =\delta\,\frac{L-e-d}{q+L}
 \le\frac{\delta\,(L-e)}{q+L}
 \le\frac{\delta}{2e}\cdot\frac{\delta}{19/30}
 \le2.6316\,\delta^2,
\]
using $L-e=(L^2-e^2)/(L+e)\le\delta/(2e)\le\delta/0.6$ from \eqref{eq:id2}
(note $L^2-e^2\le\delta$ because $2\delta d\le\tfrac d3$).  For the third
term, (C4) and (C3) give
$3\alpha\cdot6.3849/(19/30)\le1.05\cdot6.3849\cdot\tfrac{30}{19}
\le10.5855$, so
\[
 N\ \le\ 2.6316\,\delta^2+10.5855\,(m-1)\,\delta^2
 \ \le\ \Bigl(\tfrac{2.6316}{15}+10.5855\Bigr)m\,\delta^2
 \ \le\ 10.761\,m\,\delta^2 \qquad(m\ge15).
\]
With $\alpha p\ge\tfrac13\cdot0.65=\tfrac{13}{60}$,
\begin{equation}\label{eq:wide-defect}
 \mathrm{Def}_m\ \le\ \tfrac{60}{13}\cdot10.761\,m(m-2)\,\delta^2x^{\,m-3}
 \ \le\ 49.67\,m(m-2)\,\delta^2x^{\,m-3}.
\end{equation}

\emph{Step 3 (payments).}  From Lemma~\ref{lem:zc-payfloor} with
$\kappa=d/e\ge3d$, $x\le\tfrac7{13}$, $y\le\tfrac12$ and $e^2\le\tfrac19$:
\begin{align}
 2\rho_{\mathrm{lo}}\xi\le1:\quad
 &\frac{C_m\psi}{\alpha^3p^{m-2}}\ \ge\
 \frac{2\bigl(1-(\tfrac7{13})^{14}\bigr)(1-2^{-14})\cdot9\cdot13}
      {20\cdot\tfrac32}\;d^2m\ \ge\ 7.798\,d^2m,
 \label{eq:wide-payI}\\
 2\rho_{\mathrm{lo}}\xi>1:\quad
 &\frac{C_m\psi}{\alpha^3p^{m-2}}\ \ge\
 \frac{0.8164\,(1-2^{-14})\cdot1.4459}
      {0.042875\cdot\tfrac32\cdot2}\;\delta\,d\,m\ \ge\ 9.176\,\delta\,d\,m,
 \label{eq:wide-payII}
\end{align}
using in the second line $\sqrt{2\alpha}\ge\sqrt{2/3}\ge0.8164$,
$f\ge1.4459\,\delta$ from (C4), $\alpha^3\le0.35^3=0.042875$, and
$\tfrac{4\xi+1}{4\xi+2}\ge\tfrac12$.

\emph{Step 4 (comparison).}  For $j\in\{2,3\}$ the sequence
$(m-2)^jx^{m-3}$ is decreasing in $m\ge15$, because
$x\,\bigl(\tfrac{14}{13}\bigr)^3\le\tfrac7{13}\bigl(\tfrac{14}{13}\bigr)^3
<1$.  In the branch \eqref{eq:wide-payI}, insert
\eqref{eq:wide-forcing} as $d>4.09\,\delta/(m-2)$; comparing with
\eqref{eq:wide-defect} it suffices that
\[
 49.67\,(m-2)^3x^{\,m-3}\ \le\ 49.67\cdot13^3\cdot
 \bigl(\tfrac7{13}\bigr)^{12}\ \le\ 64.84\ \le\ 130.44\ \le\
 7.798\cdot4.09^2 .
\]
In the branch \eqref{eq:wide-payII} it suffices that
\[
 49.67\,(m-2)^2x^{\,m-3}\ \le\ 49.67\cdot169\cdot
 \bigl(\tfrac7{13}\bigr)^{12}\ \le\ 4.99\ \le\ 37.52\ \le\
 9.176\cdot4.09 .
\]
Every link in both chains is an exact rational inequality, with the
decimal buffers rounded in the safe direction:
$49.67\cdot13^3\cdot(7/13)^{12}=64.8306\ldots\le64.84$ and
$49.67\cdot169\cdot(7/13)^{12}=4.9869\ldots\le4.99$ (upper bounds of the
left sides, rounded up), while $130.44\le7.798\cdot4.09^2=130.4457\ldots$
and $37.52\le9.176\cdot4.09=37.5298\ldots$ (lower bounds of the right
sides, rounded down).  Part~5 of the validation script checks each link
separately in \texttt{Fraction} arithmetic.
\end{proof}

\subsection{Assembly and remarks}\label{sec:zoneC-assembly}

\begin{proof}[Proof of Theorem~\ref{thm:zoneC-main}]
If $e\le\tfrac3{10}$, the Zone-C constraints $\xi\le1$ and $q>\tfrac13$
imply $d\le d_{\max}(e)$, so Proposition~\ref{prop:strip} applies.  If
$e\ge\tfrac3{10}$, i.e.\ $\delta\le\tfrac1{60}$,
Proposition~\ref{prop:wide} applies (it needs neither $\xi\le1$ nor
$e<\tfrac13-\tfrac1{1000}$).  The dual-certificate clause is inherited from
Lemma~\emph{threegeo}(iii), which is the only payment bound used.
\end{proof}

\begin{remark}[what this deletes from the paper]
Theorem~\ref{thm:zoneC-main} replaces, in the proof of Lemma~\emph{zoneC-mod}:
the exact interval certificate \texttt{zoneC\_certifier.py} ($3001$ boxes,
depth $28$, explicit cycle lengths up to $m=1716$), the strict-gate bound
$M_+$, the per-box $m$-loop with its geometric tail argument, and the
$\kappa\to0$ bottom-out.  Lemma~\emph{turan} ($\delta<\tfrac1{2000}$)
becomes the trivial sub-case of Proposition~\ref{prop:wide}.  The only
remaining checks are: the $21\times3$ evaluations of Table~\ref{tab:rows}
(explicit elementary expressions at rational endpoints) and the
finitely many decimal constants of Section~\ref{sec:zoneC-wide} — all
hand-checkable, none requiring subdivision, interval arithmetic, or a scan
over $m$.  After this change, the only computer certificate left in
Region~II is Zone B's $23$ boxes.
\end{remark}

\begin{remark}[where the slack lives]
The tightest entries are $\Psi_{\mathrm{deep}}=0.391$ on
$[\tfrac1{60},\tfrac1{48}]$ (the ridge at the junction with the small-$e$
lemma, where the certificate itself was tightest: the true pointwise ratio
there is $\approx0.136$) and $\Psi_{\mathrm{II}}=0.703$ on
$[\tfrac{11}{50},\tfrac6{25}]$; the wide-corner branch~I comparison closes
with factor $130.44/64.84\approx2.0$.  The three structural losses
accepted for the sake of $m$-freeness are: the unconstrained supremum
$\varphi(n_*)$ (worth $\le\eu$ when $n_*<13$), the crude
$\tfrac{1+4\xi}{1+2\xi}\ge1$, and $h\le\eu^{-1+v/2}$; none is close to
critical.
\end{remark}

\begin{remark}[hypotheses actually used]
The proof uses $m\ge15$ only through $n\ge13$ (the exponent $13$ in $G_2$),
$1-y^{m-1}\ge1-y^{14}$, and the constant $15$ in
Lemma~\ref{lem:zc-master}; the oddness of $m$ is never used.  The domain
$\alpha\le r(q)$ is used only through $L^2>0$, $\ell<1$, $y<s$
(Lemma~\emph{frontier-gap}).
\end{remark}

\section{The residual strip, analytically}\label{sec:strip}


\subsection*{Provenance and verification status}
Proved by Claude on 2026-07-17, in four staged parts (left-surplus
criterion; generalized upper-end reflection; $\Lambda$-cap and
monotonicity reductions; tail and finite pairs) merged into a single
self-contained note; independent verifier: \emph{pass} ($\approx70$
exact-arithmetic check groups, including the criterion at the worst point
$(q,\ell)=(\tfrac13,\lbar)$ for \emph{all} $2500$ residual pairs with
$m\le201$, and certified-quadrature positivity probes of the kernel
integral itself); adversarial reviewer: \emph{sound}.  The theorem proves
diagonal positivity on the entire residual strip --- including $r=1$ ---
for $q\in[0,\tfrac13]$, $\ell\in(0,\tfrac12]$, $\ell<q+r/m$: the tail
$m\ge31$ closes analytically, and the $49$ remaining pairs split into
groups of $22+6$ (closed by displayed small comparisons) and $21$
(closed by the one-line exact rational inequalities of the displayed
table --- the only computer-assisted residue, in the sanctioned
hand-checkable style, though the $21$ cross-multiplications involve
large integers and are more comfortably checked by the companion
script).  Consequently it supersedes at one stroke: Prop.~\emph{M61}
($\operatorname{StripCert}(61)$), the residual-strip content of
Thm.~\emph{r1}, the three-case \S\emph{analytic-strip}, and ---
by inlining self-contained copies of the four constant lemmas ---
Appendix~\emph{app:constants} of [R1] (whose standalone analytic
replacement is nevertheless kept as \S\ref{sec:appconst}, since it is a
drop-in edit for [R1] itself).  The content is the note
\texttt{claude\_jul17/strip\_final.tex}, inlined verbatim (sectioning
demoted; \emph{thm:main} renamed \emph{thm:strip-main}).

\subsection{Setting and main result}\label{sec:setting}

A \emph{residual pair} is a pair of integers $(m,r)$ with $r\ge1$ and $n:=m-2r$ odd with $n>2r$ (so $m\ge5$ is odd).  Throughout, $q\in[0,1/3]$, $p=1-q$, and
\[
        \theta=\frac rm\in\bigl(0,\tfrac14\bigr),\quad
        \nu=\frac nm=1-2\theta,\quad
        a=\frac12-q,\quad
        b=\nu-q,\quad
        L=b-a=\nu-\frac12>0,\quad
        \varepsilon=\frac{1-2q}{4}=\frac a2 .
\]
The \emph{strip} of the pair is $\{(q,\ell):q\in[0,\tfrac13],\ \ell\in(0,\tfrac12],\ \ell<q+\tfrac rm\}$.  For $\ell>0$ put $\kappa(s)=s^{r-1}(\ell+s)^{-m}$ and
\[
        I=I_{m,r}(q,\ell)
        =\int_0^\infty \kappa(s)\,\rho_{n,m}(q+s)\,ds,
        \qquad
        \rho_{n,m}(u)=\frac mn\bigl(u^n+(1-u)^n\bigr)-u^{n-1}.
\]
By the one-dimensional kernel form (\texttt{paper\_new.tex}, Prop.~\emph{kernel}), $\wtP_{m,r}(q,\ell)=C_{m,r}\ell^{\,n+r}I$ with $C_{m,r}>0$, so $I\ge0$ is exactly diagonal positivity on the strip.  We write $\rho=\rho_{n,m}$.  Define the \emph{handoff function}
\begin{equation}\label{eq:Lambda}
        \Lambda(q)\;=\;\frac{m}{\dfrac{r-1}{b}+\dfrac{n-1}{\nu}}-b
        \qquad(\text{well defined: }n\ge3\text{ gives }\tfrac{n-1}\nu>0,\ \tfrac{r-1}b\ge0).
\end{equation}

\begin{theorem}[Unified strip theorem]\label{thm:strip-main}
For every residual pair $(m,r)$ with $r\ge1$, every $q\in[0,1/3]$, and every $\ell\in(0,\tfrac12]$ with $\ell<q+\tfrac rm$,
\[
        I_{m,r}(q,\ell)\;\ge\;0,
        \qquad\text{hence}\qquad
        \wtP_{m,r}(q,\ell)\;\ge\;0 .
\]
\end{theorem}

The proof occupies \S\ref{sec:rho}--\S\ref{sec:assembly}: a left-surplus criterion (\S\ref{sec:leftsurplus}), an upper-end reflection lemma (\S\ref{sec:reflection}), a closed-form cap on $\Lambda$ (\S\ref{sec:cap}), and an analytic verification of the criterion at the single worst point $q=\tfrac13$ for every pair (\S\ref{sec:prod}--\S\ref{sec:finite}), assembled in \S\ref{sec:assembly}.

\begin{corollary}[Supersession]\label{cor:supersede}
In \texttt{paper\_new.tex}, keep Theorem~\textup{pointwise} (cases $\ell\le0$ or $2r\ge n$, valid for $q\le\tfrac12$) and Theorem~\textup{ibp} (case $\ell\ge q+r/m$, valid for $q\le\tfrac12$).  Together with these two, Theorem~\ref{thm:strip-main} covers every diagonal case $\wtP_{m,r}(q,\ell)\ge0$ with $q\in[0,1/3]$, $\ell\in[-\tfrac12,\tfrac12]$, and therefore:
\begin{enumerate}[label=\textup{(\alph*)}]
\item \emph{Proposition~M61 / $\operatorname{StripCert}(61)$ is superseded}: its scope ($r\ge2$, odd $m\le61$, $(q,\ell)\in[0,\tfrac13]\times[0,\tfrac12]$, $0<\ell<q+r/m$) is strictly contained in Theorem~\ref{thm:strip-main}; the exact-rational interval checker leaves the proof chain.
\item \emph{Theorem~r1 is superseded}: its new content was the range $0<\ell<q+\tfrac1m$ at $r=1$, which is the $r=1$ instance of Theorem~\ref{thm:strip-main}; its other cases were already delegated to Theorems pointwise/ibp.
\item \emph{\S sec:analytic-strip is superseded} (Lemmas left-surplus-tail, upper-threshold, upper-reflection-tail, and the $m\ge63$ case of prop:residual-all): Theorem~\ref{thm:strip-main} needs neither $m\ge63$, nor $n\ge33$, nor $\theta\ge\tfrac16$, nor $\ell>\tfrac25$.
\item \emph{Proposition residual-all becomes a one-line corollary} of Theorem~\ref{thm:strip-main} (its hypotheses $r\ge2$, $n>2r$, $0<\ell<q+r/m$, $q\le\tfrac13$ define a subset of our strip), and the high-density theorem's case list shrinks to: pointwise, ibp, Theorem~\ref{thm:strip-main}.
\item \emph{Appendix app:constants is superseded}: its only clients are the lemmas of \S sec:analytic-strip; the four constant lemmas this note does use (Lemmas~\ref{lem:tools}, \ref{lem:B0linear}, \ref{lem:cubic}, \ref{lem:P72}) are restated here with complete proofs, so this note is self-contained.
\end{enumerate}
What Theorem~\ref{thm:strip-main} does \emph{not} cover: $q>\tfrac13$ (Region II), $\ell\le0$ and $2r\ge n$ (Theorem pointwise), $\ell\ge q+r/m$ (Theorem ibp) --- exactly the intended division of labour.
\end{corollary}

\begin{proof}[Proof of the coverage claim]
Fix $q\le\tfrac13$ and a diagonal pair $(r,\ell)$, $\ell\in[-\tfrac12,\tfrac12]$.  If $\ell\le0$ or $2r\ge n$, Theorem pointwise applies.  Otherwise $\ell>0$ and $n>2r$; if $n$ were even, $m$ would be even --- excluded, so $(m,r)$ is a residual pair.  If $\ell\ge q+r/m$, Theorem ibp applies; else $0<\ell<q+r/m$ and $\ell\le\tfrac12$, a strip point, and Theorem~\ref{thm:strip-main} applies.
\end{proof}

\subsection{The one-sided kernel bounds}\label{sec:rho}

All four facts below have one-line proofs from the identity $\nu\rho(u)=u^n+(1-u)^n-\nu u^{n-1}$ ($n$ odd).

\begin{fact}[$\rho$-lemma]\label{fact:rho}
\begin{enumerate}[label=\textup{(\roman*)}]
\item $\rho(u)\ge0$ for all $u\ge0$ with $u\notin(\tfrac12,\nu)$.
\item $-\rho(u)\le\dfrac{\nu-u}{\nu}\,u^{n-1}$ for $\tfrac12<u<\nu$.
\item $\rho(u)\ge\dfrac mn(1-u)^n-u^{n-1}$ for $u\ge0$.
\item $\rho(u)\ge\dfrac{u-\nu}{\nu}\,u^{n-1}$ for $\nu<u\le1$.
\end{enumerate}
\end{fact}

\begin{proof}
(i) For $0\le u\le\tfrac12$: since $1-u\ge u\ge0$ and $n-1$ is even, $(1-u)^{n-1}\ge u^{n-1}$, so $(1-u)^n\ge(1-u)u^{n-1}$ and
$\nu\rho(u)\ge u^n+(1-u)u^{n-1}-\nu u^{n-1}=(1-\nu)u^{n-1}\ge0$.
For $\nu\le u\le1$: $\nu\rho(u)=(u-\nu)u^{n-1}+(1-u)^n\ge0$, both terms nonnegative ($n$ odd).
For $u>1$: $(1-u)^n=-(u-1)^n\ge-(u-1)u^{n-1}$ since $0\le u-1<u$, so $\nu\rho(u)\ge(u-\nu)u^{n-1}-(u-1)u^{n-1}=(1-\nu)u^{n-1}\ge0$.

(ii) and (iv): from $\nu\rho(u)=(u-\nu)u^{n-1}+(1-u)^n$ and $(1-u)^n\ge0$ --- valid \emph{exactly when} $u\le1$, since $n$ is odd --- we get $\nu\rho(u)\ge(u-\nu)u^{n-1}$ for all $u\le1$; rearranged, this is (ii) on $(\tfrac12,\nu)$ and (iv) on $(\nu,1]$.  The restriction $u\le1$ is where the side condition of \S\ref{sec:reflection} enters.

(iii) Drop $\tfrac mn u^n\ge0$ (valid as $u\ge0$) from $\rho(u)=\tfrac mn u^n+\tfrac mn(1-u)^n-u^{n-1}$.

Negative $u$ is never used ($u=q+s\ge0$).
\end{proof}

\subsection{The left-surplus criterion}\label{sec:leftsurplus}

This section proves the case-(i) mechanism.  With
\begin{equation}\label{eq:cn}
        c_n=c_n(m,r,q)=1-\nu\,\frac{(q+\varepsilon)^{n-1}}{(p-\varepsilon)^{n}},
\end{equation}
define the \emph{criterion ratio}
\begin{equation}\label{eq:R}
        R(m,r,q,\ell)
        \;=\;
        c_n\cdot\frac{b}{r L^{2}}\cdot
        \bigl(2(p-\varepsilon)\bigr)^{n}
        \left(\frac{\varepsilon}{b}\right)^{r}
        \left(\frac{\ell+a}{\ell+\varepsilon}\right)^{m}.
\end{equation}

\begin{theorem}[Left-surplus criterion]\label{thm:leftsurplus}
Let $(m,r)$ be a residual pair, $q\in[0,1/3]$, and $0<\ell<q+\tfrac rm$.  If $R(m,r,q,\ell)\ge1$, then $I_{m,r}(q,\ell)\ge0$.
\end{theorem}

No upper bound on $\ell$, no lower bound on $n$ or $r$, and no reference to $\Lambda$ enters the hypotheses: the chain below uses only $q\in[0,\tfrac13]$ and $0<\ell<q+r/m$.  (The stage-A1 note stated this theorem with the case label $\ell\le\min(\max(\Lambda(q),q),\tfrac12)$ attached; as observed in stage B1, that label is not a hypothesis of the mechanism, and we drop it here.)

\begin{lemma}[Deficit monotonicity, all $r\ge1$]\label{lem:defmon}
Let $(m,r)$ be a residual pair, $q\in[0,1/3]$, and $0<\ell<q+r/m$.  Then for every $s\ge a$,
\begin{equation}\label{eq:defmon}
        (n-1)(\ell-q)\;\le\;(2r+1)(q+s),
\end{equation}
and consequently $s\mapsto(q+s)^{n-1}/(\ell+s)^{m}$ is nonincreasing on $[a,\infty)$.
\end{lemma}

\begin{proof}
First \eqref{eq:defmon}.  If $\ell\le q$ the left side is nonpositive and there is nothing to prove, so assume $\ell>q$.  From $\ell<q+r/m$,
$(n-1)(\ell-q)<(n-1)r/m$, while $s\ge a$ gives $q+s\ge q+a=\tfrac12$, so $(2r+1)(q+s)\ge\tfrac{2r+1}{2}$.  Hence \eqref{eq:defmon} follows from
$2r(n-1)\le(2r+1)m$, an identity-level triviality: with $m=n+2r$,
\[
        (2r+1)m-2r(n-1)=(2r+1)(n+2r)-2r(n-1)=n+4r^{2}+4r\;>\;0
\]
for all integers $r\ge1$, $n\ge1$.  (Exact check: \texttt{validate\_strip\_final.py}, step~A1-1.)

For the monotonicity: the logarithmic derivative of $(q+s)^{n-1}(\ell+s)^{-m}$ is $\frac{n-1}{q+s}-\frac{m}{\ell+s}$, which is $\le0$ iff $(n-1)(\ell+s)\le m(q+s)$, i.e.\ iff $(n-1)(\ell-q)\le(m-n+1)(q+s)=(2r+1)(q+s)$, which is \eqref{eq:defmon}.
\end{proof}

By Fact~\ref{fact:rho}(i), the integrand of $I$ can be negative only for $q+s\in(\tfrac12,\nu)$, i.e.\ $s\in(a,b)$.  Define the \emph{deficit} and the \emph{surplus} (collected on the band $(0,\varepsilon)$, which is disjoint from $(a,b)$ since $\varepsilon=a/2<a$; $a\ge\tfrac16$ for $q\le\tfrac13$):
\[
        D=\int_a^b \kappa(s)\bigl(-\rho(q+s)\bigr)_+\,ds\ \ge0,
        \qquad
        \Sigma=\int_0^{\varepsilon}\kappa(s)\,\rho(q+s)\,ds .
\]

\begin{lemma}[Deficit upper bound]\label{lem:deficit}
Let $(m,r)$ be a residual pair, $q\in[0,1/3]$, $0<\ell<q+r/m$.  Then
\begin{equation}\label{eq:Dbar}
        D\;\le\;\overline D
        \;:=\;
        b^{\,r-1}\,\frac{(1/2)^{n-1}}{(\ell+a)^{m}}\,\frac{L^{2}}{2\nu}
        \;>\;0.
\end{equation}
\end{lemma}

\begin{proof}
On $(a,b)$, Fact~\ref{fact:rho}(ii) at $u=q+s\in(\tfrac12,\nu)$ gives $(-\rho(q+s))_+\le(q+s)^{n-1}\,\frac{\nu-(q+s)}{\nu}$, the right side nonnegative on the whole window.  Hence
\[
        D\;\le\;\int_a^b s^{r-1}\,\frac{(q+s)^{n-1}}{(\ell+s)^{m}}\,
        \frac{\nu-q-s}{\nu}\,ds .
\]
Bound the three factors separately: $s^{r-1}\le b^{\,r-1}$ (for $r=1$ both sides are $1$); by Lemma~\ref{lem:defmon} the middle factor is nonincreasing on $[a,b]$, hence at most $(q+a)^{n-1}/(\ell+a)^{m}=(1/2)^{n-1}/(\ell+a)^{m}$; and the remaining integral is exact:
$\int_a^b\frac{\nu-q-s}{\nu}\,ds=\frac1\nu\int_0^{L}t\,dt=\frac{L^{2}}{2\nu}$ (substitute $t=b-s$, $b-a=L$).  Multiplying gives \eqref{eq:Dbar}; positivity is clear ($b\ge\nu-\tfrac13>\tfrac16$, $L>0$).
\end{proof}

\begin{lemma}[Surplus lower bound]\label{lem:surplus}
Let $(m,r)$ be a residual pair, $q\in[0,1/3]$, $\ell>0$.  With $c_n$ as in \eqref{eq:cn},
\begin{equation}\label{eq:Sbar}
        \Sigma\;\ge\;\underline\Sigma
        \;:=\;
        \frac mn\,(p-\varepsilon)^{n}\,c_n\,
        \frac{\varepsilon^{r}}{r\,(\ell+\varepsilon)^{m}} .
\end{equation}
\end{lemma}

\begin{proof}
For $s\in(0,\varepsilon)$ we have $u=q+s\in[q,q+\varepsilon]\subset[0,\tfrac12)$, so Fact~\ref{fact:rho}(iii) applies: $\rho(q+s)\ge\frac mn(p-s)^{n}-(q+s)^{n-1}$.  On $(0,\varepsilon)$, $s\mapsto(p-s)^n$ is decreasing (as $p-s\ge p-\varepsilon>0$) and $s\mapsto(q+s)^{n-1}$ is increasing, so
\[
        \rho(q+s)\;\ge\;\frac mn(p-\varepsilon)^{n}-(q+\varepsilon)^{n-1}
        =\frac mn(p-\varepsilon)^{n}
        \Bigl(1-\nu\,\frac{(q+\varepsilon)^{n-1}}{(p-\varepsilon)^{n}}\Bigr)
        =\frac mn(p-\varepsilon)^{n}c_n,
\]
using $\tfrac nm=\nu$.  Also $(\ell+s)^{-m}\ge(\ell+\varepsilon)^{-m}$ on $(0,\varepsilon)$, and $\int_0^{\varepsilon}s^{r-1}ds=\varepsilon^{r}/r$ exactly.  Since $c_n>0$ (Lemma~\ref{lem:cn} below), the three bounds multiply to \eqref{eq:Sbar}.
\end{proof}

\begin{lemma}[Lower bounds for $c_n$]\label{lem:cn}
For every residual pair and every $q\in[0,1/3]$,
\[
        c_n\;\ge\;1-\frac{12}{7}\Bigl(\frac57\Bigr)^{n-1}.
\]
In particular $c_n\ge\frac{43}{343}>0$ for all $n\ge3$, and $c_n\ge1-\frac{7500}{16807}=\frac{9307}{16807}>\frac12$ for all $n\ge5$.
\end{lemma}

\begin{proof}
Write $t=q+\varepsilon=\frac{1+2q}{4}$ and $w=p-\varepsilon=\frac{3-2q}{4}$, so $c_n=1-\nu\,t^{n-1}/w^{n}$.  On $q\in[0,1/3]$, $t$ is increasing and $w$ is decreasing in $q$ with $w\ge\tfrac7{12}>0$, so $t^{n-1}/w^{n}$ is maximal at $q=\tfrac13$, where $t=\tfrac5{12}$, $w=\tfrac7{12}$, and
$t^{n-1}/w^{n}\le(5/12)^{n-1}/(7/12)^{n}=\frac{12}{7}(\frac57)^{n-1}$.
Together with $\nu\le1$ this gives the display.  The two evaluations are exact rational comparisons:
$n=3$: $\frac{12}{7}\cdot\frac{25}{49}=\frac{300}{343}<1$, so $c_3\ge\frac{43}{343}$;
$n=5$: $\frac{12}{7}\cdot\frac{625}{2401}=\frac{7500}{16807}\le\frac12\iff15000\le16807$.
(Both checked exactly; step~A1-1.)  For the single pair with $n=3$, namely $(5,1)$, the exact value at $q=\tfrac13$ is $c_3=\tfrac{163}{343}\approx0.475<\tfrac12$, so $c_n$ must be kept as an explicit factor in \eqref{eq:R}.
\end{proof}

\begin{proposition}[Ratio identity]\label{prop:ratio}
With $\underline\Sigma$, $\overline D$, $R$ as in \eqref{eq:Sbar}, \eqref{eq:Dbar}, \eqref{eq:R}: $\underline\Sigma/\overline D=R(m,r,q,\ell)$.
\end{proposition}

\begin{proof}
Direct computation:
\[
        \frac{\underline\Sigma}{\overline D}
        =\frac{\frac mn(p-\varepsilon)^{n}c_n\,\dfrac{\varepsilon^{r}}{r(\ell+\varepsilon)^{m}}}
              {b^{\,r-1}\dfrac{(1/2)^{n-1}}{(\ell+a)^{m}}\dfrac{L^{2}}{2\nu}}
        =c_n\cdot\frac mn\cdot\frac{2\nu}{L^{2}}\cdot
         \frac{(p-\varepsilon)^{n}}{(1/2)^{n-1}}\cdot
         \frac{\varepsilon^{r}}{r\,b^{\,r-1}}\cdot
         \Bigl(\frac{\ell+a}{\ell+\varepsilon}\Bigr)^{m}.
\]
Now $\frac mn\cdot\nu=1$, $\frac{(p-\varepsilon)^n}{(1/2)^{n-1}}=\tfrac12(2(p-\varepsilon))^n$, and $\frac{\varepsilon^r}{r\,b^{r-1}}=\frac br(\varepsilon/b)^r$; collecting gives $R$.
\end{proof}

\begin{proof}[Proof of Theorem~\ref{thm:leftsurplus}]
Decompose $I=\int_0^\infty\kappa(s)\rho(q+s)\,ds$ over the pairwise disjoint bands $(0,\varepsilon]$, $(\varepsilon,a]$, $(a,b)$, $[b,\infty)$ ($\varepsilon=a/2<a<b$).  On $(\varepsilon,a]$ and $[b,\infty)$, $u=q+s\ge0$ and $u\notin(\tfrac12,\nu)$, so the integrand is nonnegative by Fact~\ref{fact:rho}(i); on $(a,b)$ the integrand is $\ge-\kappa(s)(-\rho(q+s))_+$.  Hence
\[
        I\;\ge\;\Sigma-D\;\ge\;\underline\Sigma-\overline D
        \;=\;\overline D\,\bigl(R(m,r,q,\ell)-1\bigr)\;\ge\;0
\]
by Lemmas~\ref{lem:surplus}, \ref{lem:deficit}, Proposition~\ref{prop:ratio}, $\overline D>0$, and the hypothesis.  The only two bands carrying estimates, $(0,\varepsilon)$ and $(a,b)$, are disjoint by construction; no reflection is used, so no reflection/surplus overlap can arise.
\end{proof}

\begin{proposition}[Monotonicity of $R$]\label{prop:mono}
Fix a residual pair $(m,r)$.
\begin{enumerate}[label=\textup{(\alph*)}]
\item For fixed $q\in[0,1/3]$, $\ell\mapsto R(m,r,q,\ell)$ is positive and strictly decreasing on $(0,\infty)$.
\item For fixed $\ell>0$, $q\mapsto R(m,r,q,\ell)$ is strictly decreasing on $[0,1/3]$.
\end{enumerate}
\end{proposition}

\begin{proof}
(a) Only the last factor of \eqref{eq:R} depends on $\ell$.  Since $\varepsilon=a/2<a$ (as $a\ge\tfrac16>0$),
$\frac{d}{d\ell}\,\frac{\ell+a}{\ell+\varepsilon}=\frac{\varepsilon-a}{(\ell+\varepsilon)^{2}}<0$,
and the $m$-th power of a positive decreasing function is decreasing; positivity holds as $\ell+a>\ell+\varepsilon>0$ and $c_n>0$.

(b) We show each factor of \eqref{eq:R} is positive and nonincreasing in $q$, at least one strictly.  Recall $\tfrac{da}{dq}=-1$, $\tfrac{d\varepsilon}{dq}=-\tfrac12$, $\tfrac{db}{dq}=-1$; $L,\nu$ do not depend on $q$.
\begin{itemize}
\item $c_n=1-\nu t^{n-1}/w^{n}$ with $t=\frac{1+2q}{4}$ increasing, $w=\frac{3-2q}{4}$ decreasing and positive: $t^{n-1}/w^{n}$ is increasing, so $c_n$ is decreasing; $c_n>0$ by Lemma~\ref{lem:cn}.
\item $b/(rL^{2})$: $b$ is decreasing and positive ($b>\tfrac16$), $rL^{2}$ constant.
\item $(2(p-\varepsilon))^{n}=\bigl(\tfrac{3-2q}{2}\bigr)^{n}$: decreasing, positive.
\item $(\varepsilon/b)^{r}$: $\frac{d}{dq}\log\frac{\varepsilon}{b}=-\frac{2}{1-2q}+\frac{1}{\nu-q}<0$ because $2(\nu-q)>1-2q\iff\nu>\tfrac12$, true for every residual pair.
\item $\bigl(\frac{\ell+a}{\ell+\varepsilon}\bigr)^{m}$:
$\frac{d}{dq}\log\frac{\ell+a}{\ell+\varepsilon}
=\frac{-1}{\ell+a}+\frac{1/2}{\ell+\varepsilon}<0
\iff\ell+a<2(\ell+\varepsilon)=2\ell+a\iff0<\ell$: strictly decreasing, positive.
\end{itemize}
A product of positive nonincreasing factors, one strictly decreasing, is strictly decreasing.  (Re-verified by random finite differences; step~A1-3.)
\end{proof}

\begin{corollary}[Two-variable comparison]\label{cor:twovar}
For any residual pair, any $q\in[0,1/3]$, $\ell>0$, and any $\lbar\ge\ell$:
$R(m,r,q,\ell)\ge R(m,r,\tfrac13,\lbar)$.
\end{corollary}

\begin{proof}
$R(m,r,q,\ell)\ge R(m,r,q,\lbar)\ge R(m,r,\tfrac13,\lbar)$ by Proposition~\ref{prop:mono}(a),(b).
\end{proof}

\subsection{The upper-end reflection lemma}\label{sec:reflection}

This is the case-(ii) mechanism.  The reflection sends the negative window $u\in(\tfrac12,\nu)$, parametrized as $u=\nu-e$ with $e\in(0,L)$, to the reflected points $u'=\nu+e$; Fact~\ref{fact:rho}(iv) is applied at $u'$, so we need $u'\le1$ throughout the band.

\begin{lemma}[Side condition]\label{lem:side}
For a residual pair $(m,r)$, the following are equivalent: \textup{(i)} $\nu+e\le1$ for every $e\in(0,\nu-\tfrac12)$; \textup{(ii)} $2\nu-\tfrac12\le1$; \textup{(iii)} $\theta\ge\tfrac18$.  Under \textup{(iii)} every reflected point satisfies the strict bound $u'=\nu+e<2\nu-\tfrac12\le1$, so Fact~\ref{fact:rho}(iv) applies at every $u'$ in the open band.  For $\theta<\tfrac18$ the reflected band leaves $[0,1]$, so $\tfrac18$ is sharp for this argument.
\end{lemma}

\begin{proof}
The supremum of $\nu+e$ over $e\in(0,\nu-\tfrac12)$ is $2\nu-\tfrac12$ (not attained), so (i)$\Leftrightarrow$(ii); and $2\nu-\tfrac12\le1\iff\nu\le\tfrac34\iff1-2\theta\le\tfrac34\iff\theta\ge\tfrac18$.  The final claims hold because $e<\nu-\tfrac12$ is strict.
\end{proof}

\begin{theorem}[Generalized upper-end reflection]\label{thm:reflection}
Let $(m,r)$ be a residual pair with $\theta\ge\tfrac18$, let $q\in[0,\tfrac13]$, and let $\ell\in(0,\tfrac12]$ satisfy $\ell\ge\max(\Lambda(q),q)$.  Then $I_{m,r}(q,\ell)\ge0$.
\end{theorem}

\begin{proof}
\emph{Step 0: consequences of the hypotheses.}  Put $C=\ell+b$.  From $\ell\ge q$,
\begin{equation}\label{eq:chain}
        b\;\le\;\nu\;\le\;C,
\end{equation}
since $b=\nu-q\le\nu$ (as $q\ge0$) and $C-\nu=\ell-q\ge0$.  From $\ell\ge\Lambda(q)$ and \eqref{eq:Lambda},
$C=\ell+b\ge\Lambda(q)+b=m\big/\bigl(\tfrac{r-1}{b}+\tfrac{n-1}{\nu}\bigr)$,
and since $C>0$ and the denominator is positive, this is equivalent to the \emph{payment condition}
\begin{equation}\label{eq:payment}
        \frac{r-1}{b}+\frac{n-1}{\nu}\;\ge\;\frac mC ,
\end{equation}
which thus holds \emph{by definition} of $\Lambda$; no auxiliary threshold estimate is needed.

\emph{Step 1: band decomposition.}  The negative window of $\rho(q+s)$ is $s\in(a,b)$ by Fact~\ref{fact:rho}(i).  Split $I$ over $(0,a)$, $(a,b)$, $(b,b+L)$, $(b+L,\infty)$ --- pairwise disjoint, so no mass is used twice; in particular the reflection partner band $(b,b+L)$ is disjoint from the deficit band $(a,b)$.  On $(0,a)$, $q+s\le\tfrac12$; on $(b+L,\infty)$, $q+s\ge\nu$: both contributions are $\ge0$ by Fact~\ref{fact:rho}(i).

\emph{Step 2: reflection pairing.}  Substituting $s=b\mp e$, $e\in(0,L)$:
\[
        I_-+I_+=\int_0^{L}
        \Bigl[\kappa(b+e)\,\rho(\nu+e)+\kappa(b-e)\,\rho(\nu-e)\Bigr]\,de .
\]
Fix $e\in(0,L)$.  At the deficit point $u=\nu-e\in(\tfrac12,\nu)$, Fact~\ref{fact:rho}(ii) gives $-\rho(\nu-e)\le\frac e\nu(\nu-e)^{n-1}$.  At the reflected point $u'=\nu+e$, Lemma~\ref{lem:side} (here $\theta\ge\tfrac18$ is used, and only here) gives $u'<1$, so Fact~\ref{fact:rho}(iv) yields $\rho(\nu+e)\ge\frac e\nu(\nu+e)^{n-1}$.  Hence the integrand is at least
$\frac e\nu[\kappa(b+e)(\nu+e)^{n-1}-\kappa(b-e)(\nu-e)^{n-1}]$,
which is $\ge0$ provided
\begin{equation}\label{eq:product}
        \Bigl(\frac{b+e}{b-e}\Bigr)^{r-1}
        \Bigl(\frac{\nu+e}{\nu-e}\Bigr)^{n-1}
        \;\ge\;
        \Bigl(\frac{C+e}{C-e}\Bigr)^{m}
        \qquad\bigl(\kappa(b\pm e)=(b\pm e)^{r-1}(C\pm e)^{-m}\bigr).
\end{equation}

\emph{Step 3: odd-log series.}  All three base points exceed $e$: $e<L=\nu-\tfrac12<\nu-q=b\le\nu\le C$, using $q<\tfrac12$ and \eqref{eq:chain}.  Therefore
$\log\frac{X+e}{X-e}=2\sum_{j\ge0}\frac{e^{2j+1}}{(2j+1)X^{2j+1}}$ ($0<e<X$) applies with $X\in\{b,\nu,C\}$, and \eqref{eq:product} is equivalent to
\[
        \sum_{j\ge0}\frac{2\,e^{2j+1}}{2j+1}
        \left[\frac{r-1}{b^{2j+1}}+\frac{n-1}{\nu^{2j+1}}-\frac{m}{C^{2j+1}}\right]
        \;\ge\;0 .
\]
Each bracket is nonnegative: fix $j\ge0$; by \eqref{eq:chain}, $b^{-2j}\ge\nu^{-2j}\ge C^{-2j}$, so
\[
        \frac{r-1}{b^{2j+1}}+\frac{n-1}{\nu^{2j+1}}
        =\frac{r-1}{b}\cdot\frac1{b^{2j}}+\frac{n-1}{\nu}\cdot\frac1{\nu^{2j}}
        \;\ge\;
        \Bigl(\frac{r-1}{b}+\frac{n-1}{\nu}\Bigr)\frac1{C^{2j}}
        \;\ge\;\frac{m}{C^{2j+1}},
\]
the last inequality being \eqref{eq:payment}.  (For $r=1$ the first summand is $0$ on both sides.)  Hence \eqref{eq:product} holds for every $e\in(0,L)$, so $I_-+I_+\ge0$ and $I\ge0$.
\end{proof}

\begin{remark}[Where each hypothesis is used]
$\theta\ge\tfrac18$ only places the reflected band inside $[0,1]$ (Lemma~\ref{lem:side}), which is what Fact~\ref{fact:rho}(iv) requires.  $\ell\ge\Lambda(q)$ enters only through \eqref{eq:payment}, at its exact threshold.  $\ell\ge q$ enters only through $\nu\le C$ in \eqref{eq:chain} ($\nu\le C\iff\ell\ge q$ exactly).  No lower bound on $m$ and no upper bound on $\ell$ is needed.  This theorem frees Lemma upper-reflection-tail of \texttt{paper\_new.tex} from its hypotheses $m\ge63$, $\theta\ge\tfrac16$, $\ell>\tfrac25$: the constant $\tfrac25$ and the residue count $(r-1)/m\ge\tfrac2{13}$ disappear because the payment condition is taken at $\ell=\Lambda(q)$ rather than at a uniform numeric surrogate.
\end{remark}

\subsection{The $\Lambda$-cap}\label{sec:cap}

$\Lambda$ is not monotone in $q$ for $r\ge2$, and the inequality $\Lambda(q)\ge\tfrac12$ is \emph{false} in general (e.g.\ $\Lambda(0)=\tfrac4{17}$ for $(17,2)$; even $\Lambda(q)\ge q+\theta$ fails, e.g.\ $(49,6)$, $q=\tfrac13$: $\Lambda=\tfrac{54808}{136563}<\tfrac{67}{147}$; step~B1-3).  What holds is a sharp $q$-free \emph{upper} cap, in closed form.

\begin{lemma}[M\"obius form and $\Lambda$-cap]\label{lem:cap}
Let $(m,r)$ be a residual pair ($r\ge1$) and write $E=(r-1)\nu$, $B=n-1$, $A=n-E$.  Then for every $q$ with $b=\nu-q>0$ (in particular for every $q\in[0,1/3]$),
\begin{equation}\label{eq:moebius}
        \Lambda(q)\;=\;\frac{b\,(A-Bb)}{E+Bb}
        \;=\;\frac1B\Bigl[(n+E)-t-\frac{En}{t}\Bigr]
        \Big|_{t=E+Bb},
\end{equation}
and consequently
\begin{equation}\label{eq:cap}
        \Lambda(q)\;\le\;\Lhat=\Lhat(m,r)
        \;:=\;\frac{\bigl(\sqrt n-\sqrt E\bigr)^2}{n-1}
        \;=\;\frac{n}{n-1}\Bigl(1-\sqrt{\tfrac{r-1}{m}}\Bigr)^{\!2}
        \;=\;\nu\,\frac{\bigl(\sqrt m-\sqrt{r-1}\bigr)^2}{n-1}.
\end{equation}
\end{lemma}

\begin{proof}
Since $\frac{r-1}b+\frac{n-1}\nu=\frac{(r-1)\nu+(n-1)b}{b\nu}=\frac{E+Bb}{b\nu}$ and $m\nu=n$,
\[
        \Lambda(q)=\frac{mb\nu}{E+Bb}-b
        =\frac{b\bigl(n-E-Bb\bigr)}{E+Bb}
        =\frac{b(A-Bb)}{E+Bb}.
\]
Substituting $t=E+Bb$ (so $t>E\ge0$, $b=(t-E)/B$, $A-Bb=A+E-t=n-t$):
$\Lambda=\frac{(t-E)(n-t)}{Bt}=\frac1B[(n+E)-t-\frac{En}{t}]$, which is \eqref{eq:moebius}.  By AM--GM, $t+\frac{En}t\ge2\sqrt{En}$ for every $t>0$ (for $E=0$, i.e.\ $r=1$, this reads $t\ge0$, trivial), hence
$\Lambda\le\frac{(n+E)-2\sqrt{En}}{B}=\frac{(\sqrt n-\sqrt E)^2}{n-1}$.
Finally $E=(r-1)\nu=n\cdot\frac{r-1}m$ gives $(\sqrt n-\sqrt E)^2=n(1-\sqrt{(r-1)/m})^2$, and $n=m\nu$ gives the third form.  (Identities and the rationalized cap $(t^2+En)^2\ge4En\,t^2$ checked exactly at $10^4$ random rational points; step~B1-1/2.)
\end{proof}

\begin{remark}[Sharpness]\label{rmk:sharp}
The cap is attained at $b^*=(\sqrt{En}-E)/B$ whenever $b^*\in[\nu-\tfrac13,\nu]$, which happens for genuine pairs: for $(113,26)$ the scanned $\max_{q\in[0,1/3]}\Lambda(q)$ agrees with $\Lhat$ to $10^{-13}$.  So no uniform-in-$q$ cap can improve on \eqref{eq:cap}.  For $r=1$, $\Lhat=\frac n{n-1}>1$ and the cap is vacuous --- $r=1$ pairs are never routed through Theorem~\ref{thm:reflection} below.
\end{remark}

\begin{corollary}[Window caps]\label{cor:capwindows}
Let $(m,r)$ be a residual pair with $r\ge2$.  Then:
\begin{enumerate}[label=\textup{(\alph*)}]
\item $\Lhat\le\dfrac{m}{m-2}\Bigl(1-\sqrt{\tfrac{r-1}{m}}\Bigr)^{2}$.
\item If $m\ge31$ and $\theta\ge\tfrac16$, then $r\ge6$ and
\[
        \Lhat\;\le\;\frac{31}{29}\Bigl(1-\sqrt{\tfrac{5\theta}6}\Bigr)^{2}
        \;\le\;
        \begin{cases}
        \dfrac{54901}{130500}\;\le\;\dfrac{43}{100}, & \theta\ge\tfrac16,\\[2mm]
        \dfrac{84289}{217500}\;\le\;\dfrac{39}{100}, & \theta\ge\tfrac{19}{100}.
        \end{cases}
\]
\end{enumerate}
\end{corollary}

\begin{proof}
(a) Since $\nu>\tfrac12$, $1\le2\nu$, so $n-1=\nu m-1\ge\nu m-2\nu=\nu(m-2)$ and the third form of \eqref{eq:cap} gives
$\Lhat\le\frac{\nu(\sqrt m-\sqrt{r-1})^2}{\nu(m-2)}=\frac{m}{m-2}(1-\sqrt{\tfrac{r-1}m})^2$.

(b) $r=\theta m\ge\tfrac{31}6>5$, so $r\ge6$ and $\frac{r-1}m=\theta(1-\tfrac1r)\ge\tfrac56\theta$.  The map $t\mapsto(1-\sqrt t\,)^2$ is nonincreasing on $[0,1]$ and $\tfrac m{m-2}\le\tfrac{31}{29}$ for $m\ge31$, giving the first bound.  For $\theta\ge\tfrac16$: $\tfrac{5\theta}6\ge\tfrac5{36}$ and, using $\sqrt5\ge\tfrac{559}{250}$ ($559^2=312481\le312500$),
\[
        \Bigl(1-\tfrac{\sqrt5}6\Bigr)^2=\frac{41}{36}-\frac{\sqrt5}{3}
        \le\frac{41}{36}-\frac{559}{750}=\frac{1771}{4500},
        \qquad
        \frac{31}{29}\cdot\frac{1771}{4500}=\frac{54901}{130500}
        \le\frac{43}{100}
\]
($\iff5490100\le5611500$).  For $\theta\ge\tfrac{19}{100}$: $\tfrac{5\theta}6\ge\tfrac{19}{120}$ and, using $\sqrt{19/120}\ge\tfrac{3979}{10000}$ ($3979^2\cdot120=1899892920\le19\cdot10^8$),
\[
        \Bigl(1-\sqrt{\tfrac{19}{120}}\Bigr)^2
        =\frac{139}{120}-2\sqrt{\tfrac{19}{120}}
        \le\frac{139}{120}-\frac{3979}{5000}=\frac{2719}{7500},
        \qquad
        \frac{31}{29}\cdot\frac{2719}{7500}=\frac{84289}{217500}
        \le\frac{39}{100}
\]
($\iff8428900\le8482500$).  (All comparisons exact; step~B2-1.)
\end{proof}

\subsection{The product form at the worst point}\label{sec:prod}

By Corollary~\ref{cor:twovar}, the criterion of Theorem~\ref{thm:leftsurplus} only needs to be verified at $q=\tfrac13$ and at one level $\lbar$ per pair.  At $q=\tfrac13$: $\varepsilon=\tfrac1{12}$, $a=\tfrac16$, $p-\varepsilon=\tfrac7{12}$, $b=\tfrac23-2\theta$, $L=\tfrac12-2\theta$, and
\begin{equation}\label{eq:cn-exact}
        c_n\;=\;1-\nu\,\frac{(5/12)^{n-1}}{(7/12)^{n}}
        \;=\;1-\nu\,\frac{12}{7}\Bigl(\frac57\Bigr)^{n-1}.
\end{equation}
Define, for $0<\theta<\tfrac14$ and $\ell>0$,
\begin{equation}\label{eq:PBF}
        P(\theta)=\frac{2/3-2\theta}{\theta\,(1/2-2\theta)^2},
        \qquad
        F(\ell)=\frac{\ell+1/6}{\ell+1/12}=\frac{12\ell+2}{12\ell+1},
        \qquad
        B_\ell(\theta)=\Bigl(\frac76\Bigr)^{1-2\theta}(8-24\theta)^{-\theta}F(\ell).
\end{equation}

\begin{lemma}[Product form]\label{lem:prod}
For every residual pair $(m,r)$ and every $\ell>0$,
\[
        R\bigl(m,r,\tfrac13,\ell\bigr)
        \;=\;c_n\cdot\frac{P(\theta)}{m}\cdot B_\ell(\theta)^m .
\]
\end{lemma}

\begin{proof}
Term by term in \eqref{eq:R}, using $r=m\theta$ and $n=m(1-2\theta)$:
$\frac{b}{rL^2}=\frac{2/3-2\theta}{m\theta(1/2-2\theta)^2}=\frac{P(\theta)}m$;
$(2(p-\varepsilon))^n=(\tfrac76)^n=((\tfrac76)^{1-2\theta})^m$;
$\frac{\varepsilon}{b}=\frac1{12b}=\frac1{8-24\theta}$, so $(\varepsilon/b)^r=((8-24\theta)^{-\theta})^m$; and
$\frac{\ell+a}{\ell+\varepsilon}=F(\ell)$.  Multiply and insert \eqref{eq:cn-exact}.  (For rational $\ell$ every factor of $R$ is rational: $B_\ell(\theta)^m=(\tfrac76)^n(\tfrac{m}{8m-24r})^rF(\ell)^m$.)
\end{proof}

With $\ell=\tfrac13+\theta$ (so $12\ell+2=6+12\theta$, $12\ell+1=5+12\theta$),
\begin{equation}\label{eq:B0}
        F\bigl(\tfrac13+\theta\bigr)=\frac{1/2+\theta}{5/12+\theta},
        \qquad
        B_0(\theta):=B_{\frac13+\theta}(\theta)
        =\Bigl(\frac76\Bigr)^{1-2\theta}\frac{(8-24\theta)^{-\theta}\,(\tfrac12+\theta)}{5/12+\theta}.
\end{equation}
$P$ and $B_0$ are exactly the functions of appendix app:constants of \texttt{paper\_new.tex}; its Case-A inequality was engineered for the left-surplus estimate at $q=\tfrac13$, $\ell=q+r/m$, and \S\ref{sec:tools} reproves the three lemmas it rests on, making this note independent of that appendix.

\subsection{Elementary tools and constant lemmas}\label{sec:tools}

\begin{lemma}[Exponential and logarithm bounds]\label{lem:tools}
\begin{enumerate}[label=\textup{(\roman*)}]
\item $e^{y}\ge1+y$ for every real $y$.
\item $B^m/m\ge e\log B$ for every real $B>0$, $m>0$.
\item $\log u\le u-1$ for every real $u>0$.
\item For every real $t\ge1$:
$L(t):=\dfrac{3(t^2-1)}{t^2+4t+1}\;\le\;\log t\;\le\;\dfrac{(t-1)(t+5)}{2(2t+1)}=:U(t)$.
\item $e\ge65/24$.
\end{enumerate}
\end{lemma}

\begin{proof}
(i) $y\mapsto e^y-1-y$ has derivative $e^y-1$, negative for $y<0$, positive for $y>0$; its minimum, at $y=0$, is $0$.
(ii) Apply (i) with $y=m\log B-1$: $B^m=e\cdot e^{\,m\log B-1}\ge e\,m\log B$; divide by $m>0$.
(iii) Apply (i) with $y=u-1$ and take logarithms.
(iv) Both bounds are equalities at $t=1$, so compare derivatives on $t\ge1$:
\[
        \frac{1}{t}-L'(t)
        =\frac{(t^2+4t+1)^2-12t(t^2+t+1)}{t\,(t^2+4t+1)^2}
        =\frac{(t-1)^4}{t\,(t^2+4t+1)^2}\ \ge\ 0,
\qquad
        U'(t)-\frac1t
        =\frac{4(t-1)^3}{t\,(4t+2)^2}\ \ge\ 0,
\]
using $L'(t)=12(t^2+t+1)/(t^2+4t+1)^2$ and, with $U(t)=(t^2+4t-5)/(4t+2)$, $U'(t)=(4t^2+4t+28)/(4t+2)^2$; integrate from $1$ to $t$.
(v) $e=\sum_{k\ge0}1/k!\ge1+1+\frac12+\frac16+\frac1{24}=\frac{65}{24}$.
\end{proof}

The rational Pad\'e values used below, each obtained from Lemma~\ref{lem:tools}(iv) by clearing denominators (exact re-check: step~B2-1):
\begin{equation}\label{eq:pade}
\begin{aligned}
        \log\frac{64}{63}&\ge L\Bigl(\frac{64}{63}\Bigr)=\frac{3\cdot127}{24193},
        &\log\frac73&\ge L\Bigl(\frac73\Bigr)=\frac{60}{71},
        &\log\frac76&\ge L\Bigl(\frac76\Bigr)=\frac{39}{253},
        \\[1mm]
        \log\frac76&\le U\Bigl(\frac76\Bigr)=\frac{37}{240},
        &\log2&\le U(2)=\frac7{10},
        &&
        \\[1mm]
        \log\frac{179}{154}&\ge L\Bigl(\frac{179}{154}\Bigr)=\frac{24975}{166021},
        &\log\frac{167}{142}&\ge L\Bigl(\frac{167}{142}\Bigr)=\frac{23175}{142909}.
        &&
\end{aligned}
\end{equation}
Set
\begin{equation}\label{eq:Kbar}
        \overline K\;=\;1+2\cdot\frac7{10}+2\cdot\frac{37}{240}
        \;=\;\frac{65}{24}
        \;\ge\;K:=1+2\log2+2\log\frac76 ,
        \qquad
        \lambda_*=\frac{127}{24193},
        \qquad
        s=\frac{89}{100}.
\end{equation}

\begin{lemma}[Linear minorant of $\log B_0$]\label{lem:B0linear}
For every $0\le\theta\le\frac16$,
$\log B_0(\theta)\ \ge\ \lambda_*+s\bigl(\frac16-\theta\bigr)$.
\end{lemma}

\begin{proof}
\emph{Value at the right endpoint.}  Since $4^{-1/6}=2^{-1/3}$,
\[
        B_0(1/6)
        =\Bigl(\frac76\Bigr)^{2/3}4^{-1/6}\,\frac{2/3}{7/12}
        =7^{2/3}6^{-2/3}\cdot2^{-1/3}\cdot\frac{8}{7}
        =\frac{8}{504^{1/3}}
        =\frac{4}{63^{1/3}},
\]
so $B_0(1/6)^3=\tfrac{64}{63}$ and, by \eqref{eq:pade}, $\log B_0(1/6)=\frac13\log\frac{64}{63}\ge\frac{127}{24193}=\lambda_*$.

\emph{Slope.}  From $\log B_0(\theta)=(1-2\theta)\log\frac76-\theta\log(8-24\theta)+\log(\frac12+\theta)-\log(\frac5{12}+\theta)$, the negated derivative $c(\theta):=-\frac{d}{d\theta}\log B_0(\theta)$ is
\[
        c(\theta)
        =2\log\frac76+\log(8-24\theta)-\frac{24\theta}{8-24\theta}
        -\frac{1}{1/2+\theta}+\frac{1}{5/12+\theta},
\]
\[
        c'(\theta)
        =-\frac{24}{8-24\theta}-\frac{192}{(8-24\theta)^2}
        +\frac{1}{(1/2+\theta)^2}-\frac{1}{(5/12+\theta)^2}.
\]
On $[0,\tfrac16]$, $8-24\theta\ge4>0$ makes the first two terms negative, and $0<\frac5{12}+\theta<\frac12+\theta$ makes the sum of the last two negative; so $c$ is strictly decreasing and, for every $\theta\in[0,\tfrac16]$,
\[
        c(\theta)\ \ge\ c(1/6)
        =2\log\frac76+\log4-1-\frac32+\frac{12}7
        =2\log\frac73-\frac{11}{14}
        \ \ge\ \frac{120}{71}-\frac{11}{14}
        =\frac{899}{994}
        \ \ge\ \frac{89}{100}=s,
\]
using \eqref{eq:pade} at $t=\tfrac73$ and $899\cdot100=89900\ge88466=89\cdot994$.

\emph{Conclusion.}  $\log B_0(\theta)=\log B_0(1/6)+\int_\theta^{1/6}c(u)\,du\ge\lambda_*+s(\frac16-\theta)$.
\end{proof}

\begin{lemma}[Weighted polynomial inequality]\label{lem:cubic}
For every $0<\theta\le\frac16$,
$P(\theta)\bigl[\lambda_*+s\bigl(\frac16-\theta\bigr)\bigr]\ \ge\ 72\,\lambda_*$.
\end{lemma}

\begin{proof}
For $0<\theta<\tfrac14$, $\theta(\tfrac12-2\theta)^2>0$, so the claim is equivalent to $G(\theta)\ge0$ on $(0,\tfrac16]$, where
\[
        G(\theta)
        :=\Bigl(\frac23-2\theta\Bigr)
        \Bigl[\lambda_*+s\Bigl(\frac16-\theta\Bigr)\Bigr]
        -72\lambda_*\,\theta\Bigl(\frac12-2\theta\Bigr)^2 .
\]
At $\theta=\tfrac16$ both terms equal $\frac13\lambda_*$, so $G(\tfrac16)=0$ and $(\tfrac16-\theta)$ divides $G$; polynomial division gives the exact factorisation
\[
        G(\theta)=\Bigl(\frac16-\theta\Bigr)Q(\theta),
        \qquad
        Q(\theta)
        =288\lambda_*\theta^2-(2s+96\lambda_*)\theta
        +\Bigl(\frac{2s}3+4\lambda_*\Bigr),
\]
verified by expanding both sides to
$(\frac{2\lambda_*}3+\frac s9)-(s+20\lambda_*)\theta+(2s+144\lambda_*)\theta^2-288\lambda_*\theta^3$.
On $[0,\tfrac16]$, $Q'(\theta)=576\lambda_*\theta-(2s+96\lambda_*)\le96\lambda_*-(2s+96\lambda_*)=-2s<0$, so
\[
        Q(\theta)\ \ge\ Q(1/6)
        =\frac s3-4\lambda_*
        =\frac{89}{300}-\frac{508}{24193}
        =\frac{2000777}{7257900}\ >\ 0,
\]
hence $Q>0$ and $G\ge0$ on $(0,\tfrac16]$.
\end{proof}

\begin{lemma}[Prefactor bound]\label{lem:P72}
For every real $\frac16\le\theta<\frac14$, $\;P(\theta)\ge72$.
\end{lemma}

\begin{proof}
One checks by expansion the polynomial identity
\[
        \frac23-2\theta-72\,\theta\Bigl(\frac12-2\theta\Bigr)^2
        =-\frac23\,(6\theta-1)(72\theta^2-24\theta+1).
\]
On $[\tfrac16,\tfrac14]$, $6\theta-1\ge0$; the convex quadratic $72\theta^2-24\theta+1$ has endpoint values $-1$ and $-\tfrac12$, so it is $\le\max\{-1,-\tfrac12\}<0$ there.  Hence the right side is nonnegative, i.e.\ $\frac23-2\theta\ge72\,\theta(\frac12-2\theta)^2$, and dividing by $\theta(\frac12-2\theta)^2>0$ gives the claim.
\end{proof}

\begin{lemma}[Quadratic minorant]\label{lem:quad}
For every real $\theta\in[\tfrac16,\tfrac14)$ and every fixed $\ell>0$,
\[
        \log B_\ell(\theta)\;\ge\;
        \log F(\ell)+\log\frac76-\overline K\theta+6\theta^2,
        \qquad \overline K=\frac{65}{24}.
\]
\end{lemma}

\begin{proof}
Write $8-24\theta=4(2-6\theta)$ with $2-6\theta\in(\tfrac12,1]$ on the stated range, so $\log(2-6\theta)\le1-6\theta$ by Lemma~\ref{lem:tools}(iii), and therefore, using $\theta>0$,
\[
        \log B_\ell(\theta)
        =\log F(\ell)+\log\frac76
        -\theta\Bigl(2\log\frac76+2\log2\Bigr)-\theta\log(2-6\theta)
        \;\ge\;\log F(\ell)+\log\frac76-K\theta+6\theta^2
\]
with $K$ as in \eqref{eq:Kbar}; replacing $K$ by $\overline K\ge K$ decreases the right side further.
\end{proof}

\subsection{Lower bounds for $c_n$ at $q=1/3$}\label{sec:cn}

\begin{lemma}[$c_n$ thresholds]\label{lem:cnthr}
Let $(m,r)$ be a residual pair and $c_n$ as in \eqref{eq:cn-exact} (the exact value at $q=\tfrac13$).  Then:
\begin{enumerate}[label=\textup{(\alph*)}]
\item if $n\ge14$, then $c_n\ge\dfrac{24193}{24765}$;
\item if $n\ge17$, then $c_n\ge\dfrac{99}{100}$.
\end{enumerate}
\end{lemma}

\begin{proof}
Since $\nu<1$ and $\tfrac57<1$, \eqref{eq:cn-exact} gives $c_n\ge1-\tfrac{12}7(\tfrac57)^{n-1}$, with deficit decreasing in $n$.
(a) It suffices that $\tfrac{12}7(\tfrac57)^{13}\le\tfrac{572}{24765}=1-\tfrac{24193}{24765}$, i.e.
\[
        24765\cdot12\cdot5^{13}
        =362\,768\,554\,687\,500\;\le\;387\,943\,597\,669\,628
        =572\cdot7^{14}.
\]
(b) It suffices that $\tfrac{12}7(\tfrac57)^{16}\le\tfrac1{100}$, i.e.
\[
        1200\cdot5^{16}=183\,105\,468\,750\,000\;\le\;232\,630\,513\,987\,207=7^{17}.
\]
Both integer comparisons are exact (step~B2-1).
\end{proof}

\subsection{The three $\theta$-windows}\label{sec:windows}

Each window proposition has \emph{scalar} hypotheses only (a $c_n$ level and a $\theta$ range); $\theta$ and $m$ are treated as real numbers, so the propositions apply both to the uniform tail and to individual small pairs.  The evaluation levels are $\tfrac13+\theta$ (W1), $\tfrac{43}{100}$ (W2), $\tfrac{39}{100}$ (W3).

\begin{proposition}[Window W1: $\theta\le1/6$]\label{prop:W1}
Let $(m,r)$ be a residual pair with $\theta\le\tfrac16$ and $c_n\ge\tfrac{24193}{24765}$.  Then
\[
        R\bigl(m,r,\tfrac13,\tfrac13+\tfrac rm\bigr)\;\ge\;
        c_n\cdot\frac{24765}{24193}\;\ge\;1 .
\]
\end{proposition}

\begin{proof}
By Lemma~\ref{lem:prod} and \eqref{eq:B0}, $R=c_n\,\frac{P(\theta)}m\,B_0(\theta)^m$.  By Lemma~\ref{lem:B0linear}, $\log B_0(\theta)\ge\lambda_*+s(\tfrac16-\theta)>0$; by Lemma~\ref{lem:tools}(ii), $B_0(\theta)^m/m\ge e\log B_0(\theta)$; and $P(\theta)>0$ on $(0,\tfrac14)$.  Hence, by Lemma~\ref{lem:cubic} and Lemma~\ref{lem:tools}(v),
\[
        R\;\ge\;c_n\,e\,P(\theta)\Bigl[\lambda_*+s\Bigl(\tfrac16-\theta\Bigr)\Bigr]
        \;\ge\;c_n\,e\cdot72\lambda_*
        \;\ge\;c_n\cdot\frac{65}{24}\cdot\frac{72\cdot127}{24193}
        \;=\;c_n\cdot\frac{24765}{24193}\;\ge\;1 . \qedhere
\]
\end{proof}

\begin{proposition}[Window W2: $1/6\le\theta\le19/100$]\label{prop:W2}
Let $(m,r)$ be a residual pair with $\tfrac16\le\theta\le\tfrac{19}{100}$ and $c_n\ge\tfrac{99}{100}$.  Then
\[
        R\bigl(m,r,\tfrac13,\tfrac{43}{100}\bigr)
        \;\ge\;\frac{3861}{20}\,V_2\;=\;
        \frac{1\,945\,991\,604\,141}{1\,527\,393\,200\,000}\;>\;\frac54\;>\;1,
        \quad
        V_2=\frac{24975}{166021}+\frac{39}{253}
        -\frac{65}{24}\cdot\frac{19}{100}+6\Bigl(\frac{19}{100}\Bigr)^2>0 .
\]
\end{proposition}

\begin{proof}
$F(\tfrac{43}{100})=\tfrac{179}{154}$, so $\log F\ge\tfrac{24975}{166021}$ and $\log\tfrac76\ge\tfrac{39}{253}$ by \eqref{eq:pade}.  By Lemma~\ref{lem:quad},
$\log B_{43/100}(\theta)\ge g(\theta):=\tfrac{24975}{166021}+\tfrac{39}{253}-\tfrac{65}{24}\theta+6\theta^2$.
Since $g'(\theta)=12\theta-\tfrac{65}{24}\le12\cdot\tfrac{19}{100}-\tfrac{65}{24}=-\tfrac{257}{600}<0$ on the window, $g(\theta)\ge g(\tfrac{19}{100})=V_2=\tfrac{16\,632\,406\,873}{2\,520\,198\,780\,000}>0$ there.  By Lemma~\ref{lem:prod}, Lemma~\ref{lem:tools}(ii),(v) and Lemma~\ref{lem:P72},
\[
        R\;=\;c_n\,\frac{P(\theta)}m\,B_{43/100}(\theta)^m
        \;\ge\;\frac{99}{100}\cdot72\cdot e\,\log B_{43/100}(\theta)
        \;\ge\;\frac{99}{100}\cdot72\cdot\frac{65}{24}\,V_2
        \;=\;\frac{3861}{20}\,V_2 ,
\]
and $\tfrac{3861}{20}V_2>\tfrac54\iff4\cdot1945991604141=7783966416564\ge7636966000000=5\cdot1527393200000$.
\end{proof}

\begin{proposition}[Window W3: $19/100\le\theta<1/4$]\label{prop:W3}
Let $(m,r)$ be a residual pair with $\tfrac{19}{100}\le\theta<\tfrac14$ and $c_n\ge\tfrac{99}{100}$.  Then
\[
        R\bigl(m,r,\tfrac13,\tfrac{39}{100}\bigr)
        \;\ge\;\frac{3861}{20}\,V_3\;=\;
        \frac{5\,342\,297\,399}{2\,589\,071\,360}\;>\;2\;>\;1,
        \quad
        V_3=\frac{23175}{142909}+\frac{39}{253}-\frac{4225}{13824}>0 .
\]
\end{proposition}

\begin{proof}
$F(\tfrac{39}{100})=\tfrac{167}{142}$, so $\log F\ge\tfrac{23175}{142909}$ by \eqref{eq:pade}.  By Lemma~\ref{lem:quad} and completion of the square, for \emph{every} real $\theta$,
\[
        \log B_{39/100}(\theta)
        \;\ge\;\frac{23175}{142909}+\frac{39}{253}
        -\frac{\overline K^2}{24}
        +6\Bigl(\theta-\frac{\overline K}{12}\Bigr)^2
        \;\ge\;\frac{23175}{142909}+\frac{39}{253}-\frac{4225}{13824}
        \;=\;V_3\;=\;\frac{5\,342\,297\,399}{499\,820\,226\,048}\;>\;0,
\]
using $\overline K^2/24=4225/13824$.  Exactly as in Proposition~\ref{prop:W2},
$R\ge\tfrac{99}{100}\cdot72\cdot\tfrac{65}{24}V_3=\tfrac{3861}{20}V_3>2$
($\iff5342297399\ge5178142720=2\cdot2589071360$).
\end{proof}

\begin{remark}
The window margins are genuine: the true minimum of the rate $\theta\mapsto\log B_{\lbar(\theta)}(\theta)$ over the ridge $\theta\in[0.15,0.22]$ is $\approx0.042$ (near $\theta=0.2$), while the minorants certify $V_2\approx0.0066$ and $V_3\approx0.0107$ against the required $\tfrac{20}{3861}\approx0.00518$ --- covered with factors $1.27$ and $2.06$.  The crude level $\lbar=\tfrac12$ would \emph{fail} here at large $m$ (the rate $(1-2\theta)\log\tfrac76+\log\tfrac87-\theta\log(8-24\theta)$ is negative near $\theta=0.2$); the caps $\tfrac{43}{100},\tfrac{39}{100}$ from Corollary~\ref{cor:capwindows} are what make the ridge windows close.
\end{remark}

\subsection{The criterion theorem}\label{sec:crit}

Let $\mathcal T$ denote the set of $21$ pairs listed in Table~\ref{tab:pairs} below.  For a residual pair $(m,r)$ define the \emph{evaluation level}
\begin{equation}\label{eq:lbar}
        \lbar_{m,r}=
        \begin{cases}
        \min\bigl(\tfrac12,\ \tfrac13+\tfrac rm\bigr), & \theta\le\tfrac16
                \text{ or } (m,r)\in\mathcal T,\\[1mm]
        \tfrac{43}{100}, & \theta\in(\tfrac16,\tfrac{19}{100}],\ (m,r)\notin\mathcal T,\\[1mm]
        \tfrac{39}{100}, & \theta\in(\tfrac{19}{100},\tfrac14),\ (m,r)\notin\mathcal T .
        \end{cases}
\end{equation}

\begin{theorem}[Criterion, all residual pairs]\label{thm:crit}
For every residual pair $(m,r)$:
\begin{enumerate}[label=\textup{(\alph*)}]
\item $R\bigl(m,r,\tfrac13,\lbar_{m,r}\bigr)\;\ge\;1$;
\item if $\theta>\tfrac16$ and $(m,r)\notin\mathcal T$, then $\Lambda(q)\le\lbar_{m,r}$ for every $q\in[0,\tfrac13]$.
\end{enumerate}
\end{theorem}

The proof is a uniform tail argument valid for all $m\ge31$ (Theorem~\ref{thm:tail}; in the range $\theta\le\tfrac16$ it is valid for all real $m>0$), plus finitely many exact rational checks for the $49$ residual pairs with $m\le29$ (\S\ref{sec:finite}), each displayed.

\begin{theorem}[Tail: $m\ge31$]\label{thm:tail}
Every residual pair with $m\ge31$ satisfies Theorem~\ref{thm:crit}\,(a),(b).
\end{theorem}

\begin{proof}
Two integrality facts: $n>2r$ and $n+2r=m$ give $n>\tfrac m2$; and $n$ is odd.  Thus $m\ge31$ implies $n\ge17$, so $c_n\ge\tfrac{99}{100}\ge\tfrac{24193}{24765}$ by Lemma~\ref{lem:cnthr} (the last comparison: $99\cdot24765=2451735\ge2419300=100\cdot24193$).  No pair with $m\ge31$ lies in $\mathcal T$.

\emph{Case $\theta\le\tfrac16$.}  Then $\lbar=\tfrac13+\tfrac rm$ (as $\tfrac13+\theta\le\tfrac12$) and Proposition~\ref{prop:W1} gives (a); (b) is vacuous.

\emph{Case $\tfrac16<\theta\le\tfrac{19}{100}$.}  By Corollary~\ref{cor:capwindows}(b) ($m\ge31$, $\theta\ge\tfrac16$, and $r\ge6\ge2$), $\Lambda(q)\le\Lhat\le\tfrac{54901}{130500}\le\tfrac{43}{100}=\lbar$ for all $q$, which is (b); and Proposition~\ref{prop:W2} gives (a).

\emph{Case $\tfrac{19}{100}<\theta<\tfrac14$.}  Identical via Corollary~\ref{cor:capwindows}(b) with $\theta\ge\tfrac{19}{100}$ ($\Lhat\le\tfrac{84289}{217500}\le\tfrac{39}{100}=\lbar$) and Proposition~\ref{prop:W3}.

Every residual pair has $\theta\in(0,\tfrac14)$, so the three cases are exhaustive.
\end{proof}

\begin{remark}\label{rmk:realm}
In the window $\theta\le\tfrac16$ no lower bound on $m$ is used: Proposition~\ref{prop:W1} needs only $c_n\ge\tfrac{24193}{24765}$, which Lemma~\ref{lem:cnthr}(a) reduces to $n\ge14$.  This already disposes of most pairs with $m\le29$.  The genuine role of $m\ge31$ is confined to Corollary~\ref{cor:capwindows}(b) (through $r\ge6$ and $\tfrac m{m-2}\le\tfrac{31}{29}$) and to $n\ge17$.
\end{remark}

\subsection{The finite pairs: $5\le m\le29$}\label{sec:finite}

There are exactly $49$ residual pairs with $m\le29$.  We close them in three groups.

\subsubsection{Group (a): window W1 by the $c_n$ threshold ($22$ pairs)}\label{sec:grpA}
By Lemma~\ref{lem:cnthr}(a), Proposition~\ref{prop:W1} applies whenever $\theta\le\tfrac16$ and $n\ge14$, i.e.\ to the $21$ pairs
\[
        (17,1);\ (19,1),(19,2);\ (21,1),(21,2),(21,3);\ (23,1),(23,2),(23,3);\
        (25,1),\dots,(25,4);\ (27,1),\dots,(27,4);\ (29,1),\dots,(29,4)
\]
(each has $6r\le m$ and $n\ge15$).  One further pair, $(19,3)$ ($\theta=\tfrac3{19}\le\tfrac16$, $n=13$), passes the threshold through the exact value \eqref{eq:cn-exact} thanks to $\nu=\tfrac{13}{19}$:
\begin{equation}\label{eq:193}
        c_{13}(19,3)\ge\frac{24193}{24765}
        \iff
        24765\cdot156\cdot5^{12}=943\,198\,242\,187\,500
        \;\le\;1\,052\,989\,765\,103\,276=572\cdot19\cdot7^{13}.
\end{equation}
For all $22$ pairs, $\lbar=\tfrac13+\tfrac rm=\min(\tfrac12,\tfrac13+\tfrac rm)$, so Theorem~\ref{thm:crit}(a) holds; (b) is vacuous ($\theta\le\tfrac16$).

\subsubsection{Group (b): windows W2/W3 by explicit handoff caps ($6$ pairs)}\label{sec:grpB}
For six pairs with $\theta>\tfrac16$ we check the hypotheses of Propositions~\ref{prop:W2}/\ref{prop:W3} directly, using Lemma~\ref{lem:cap} in the form $\sup_q\Lambda(q)\le\Lhat=\nu\,\frac{m+r-1-2\sqrt{m(r-1)}}{n-1}$ and a rational $\sigma\le\sqrt{m(r-1)}$ (safe direction: it \emph{increases} the bound; each $\sigma^2\le m(r-1)$ is one integer comparison).  The $c_n$ hypothesis holds by Lemma~\ref{lem:cnthr}(b) for $n\ge17$; for $n=15$ it is checked exactly via \eqref{eq:cn-exact}, each display being $\nu\frac{12}7(\tfrac57)^{14}\le\tfrac1{100}$ cleared of denominators:
\begin{equation}\label{eq:c15}
\begin{aligned}
        (25,5):&\ \ 3600\cdot5^{14}=21\,972\,656\,250\,000
                \le35\cdot7^{14}=23\,737\,807\,549\,715,\\
        (27,6):&\ \ 6000\cdot5^{14}=36\,621\,093\,750\,000
                \le63\cdot7^{14}=42\,728\,053\,589\,487,\\
        (29,7):&\ \ 18000\cdot5^{14}=109\,863\,281\,250\,000
                \le203\cdot7^{14}=137\,679\,283\,788\,347.
\end{aligned}
\end{equation}
The caps ($\sigma$ first, then the resulting bound on $\Lhat$):
\begin{itemize}[leftmargin=2em]
\item $(27,5)\in\mathrm{W2}$: $\sigma=\tfrac{1039}{100}$ ($1039^2=1079521\le1080000$),\quad
 $\Lhat\le\tfrac{17}{27}\cdot\tfrac{31-1039/50}{16}=\tfrac{8687}{21600}\le\tfrac{43}{100}$ ($868700\le928800$);
\item $(29,5)\in\mathrm{W2}$: $\sigma=\tfrac{1077}{100}$ ($1077^2=1159929\le1160000$),\quad
 $\Lhat\le\tfrac{19}{29}\cdot\tfrac{33-1077/50}{18}=\tfrac{3629}{8700}\le\tfrac{43}{100}$ ($362900\le374100$);
\item $(25,5)\in\mathrm{W3}$: $\sigma=10$ ($10^2=100=25\cdot4$),\quad
 $\Lhat\le\tfrac{15}{25}\cdot\tfrac{29-20}{14}=\tfrac{27}{70}\le\tfrac{39}{100}$ ($2700\le2730$);
\item $(27,6)\in\mathrm{W3}$: $\sigma=\tfrac{1161}{100}$ ($1161^2=1347921\le1350000$),\quad
 $\Lhat\le\tfrac{15}{27}\cdot\tfrac{32-1161/50}{14}=\tfrac{439}{1260}\le\tfrac{39}{100}$ ($43900\le49140$);
\item $(29,6)\in\mathrm{W3}$: $\sigma=\tfrac{301}{25}$ ($301^2=90601\le90625$),\quad
 $\Lhat\le\tfrac{17}{29}\cdot\tfrac{34-602/25}{16}=\tfrac{527}{1450}\le\tfrac{39}{100}$ ($52700\le56550$);
\item $(29,7)\in\mathrm{W3}$: $\sigma=\tfrac{1319}{100}$ ($1319^2=1739761\le1740000$),\quad
 $\Lhat\le\tfrac{15}{29}\cdot\tfrac{35-1319/50}{14}=\tfrac{1293}{4060}\le\tfrac{39}{100}$ ($129300\le158340$).
\end{itemize}
Window membership: $\theta=\tfrac5{27},\tfrac5{29}\le\tfrac{19}{100}$ ($500\le513$, $500\le551$) for the two W2 pairs, and $\theta=\tfrac15,\tfrac29,\tfrac6{29},\tfrac7{29}\ge\tfrac{19}{100}$ ($500\ge475$, $200\ge171$, $600\ge551$, $700\ge551$) for the four W3 pairs.  So Theorem~\ref{thm:crit}(a) holds by Propositions~\ref{prop:W2}/\ref{prop:W3}, and (b) holds because $\Lambda(q)\le\Lhat\le$ cap $=\lbar$.

\subsubsection{Group (c): the remaining $21$ pairs $\mathcal T$, one exact evaluation each}\label{sec:grpC}
For the pairs of $\mathcal T$ we evaluate $R$ at $\lbar=\min(\tfrac12,\tfrac13+\tfrac rm)$ and check $R\ge1$ in exact rational arithmetic.  By Lemma~\ref{lem:prod}, for each pair
\begin{equation}\label{eq:row}
        R\bigl(m,r,\tfrac13,\lbar\bigr)
        \;=\;c_n\cdot\frac b{rL^2}\cdot\Bigl(\frac76\Bigr)^{n}
        \Bigl(\frac{m}{8m-24r}\Bigr)^{r} F(\lbar)^{m},
        \qquad F(\lbar)=\frac{12\lbar+2}{12\lbar+1},
\end{equation}
a product of explicit rationals: each row of Table~\ref{tab:pairs} is one exact rational inequality.  The stated lower bounds $R_{\mathrm{lb}}$ are safe roundings (rounded \emph{down} to two decimals); every row, cleared of denominators, is a single comparison of two explicit integers, performed exactly in \texttt{validate\_strip\_final.py}, step~B2-5.  The smallest margin, row $(23,4)$, asserts
\[
        \frac{108095281916189}{109193914728689}\cdot\frac{506}{147}\cdot
        \Bigl(\frac76\Bigr)^{15}\Bigl(\frac{23}{88}\Bigr)^{4}
        \Bigl(\frac87\Bigr)^{23}
        =\frac ND\;\ge\;\frac{173}{50}\;>\;1
\]
with, in reduced form, $N=266\,077\,343\,630\,402\,608\,833\,493\,983\,035\,392$ and $D=76\,836\,775\,534\,377\,178\,226\,484\,864\,720\,357$ (so $50N\ge173D$); the smallest pair, row $(5,1)$, asserts
\[
        \frac{163}{343}\cdot\frac{80}{3}\cdot\Bigl(\frac76\Bigr)^{3}
        \Bigl(\frac5{16}\Bigr)^{1}\Bigl(\frac87\Bigr)^{5}
        =\frac{16\,691\,200}{1\,361\,367}\;\ge\;\frac{613}{50}\;>\;1 .
\]
Theorem~\ref{thm:crit}(b) is vacuous for $\mathcal T$ by definition of $\lbar$.

\begin{table}[ht]
\caption{Group (c): exact evaluation \eqref{eq:row} at $\lbar=\min(\tfrac12,\tfrac13+\tfrac rm)$.  All entries are exact rationals; $R_{\mathrm{lb}}$ is a safe (rounded-down) certified lower bound on $R(m,r,\tfrac13,\lbar)$, verified by exact cross-multiplication.}
\label{tab:pairs}
\centering
\footnotesize
\setlength{\tabcolsep}{4pt}
\begin{tabular}{@{}rrrllllr@{}}
\toprule
$m$ & $r$ & $n$ & $\lbar$ & $c_n$ & $\dfrac b{rL^2}$ & $\dfrac m{8m-24r},\ F(\lbar)$ & $R_{\mathrm{lb}}$\\
\midrule
 5&1& 3&$1/2$  &$163/343$                        &$80/3$   &$5/16,\ 8/7$    &$12.26$\\
 7&1& 5&$10/21$&$80149/117649$                   &$224/27$ &$7/32,\ 54/47$  &$7.06$\\
 9&1& 7&$4/9$  &$290447/352947$                  &$144/25$ &$3/16,\ 22/19$  &$9.78$\\
 9&2& 5&$1/2$  &$37921/50421$                    &$36$     &$3/8,\ 8/7$     &$27.37$\\
11&1& 9&$14/33$&$401702177/443889677$            &$704/147$&$11/64,\ 78/67$ &$15.87$\\
11&2& 7&$1/2$  &$1106639/1294139$                &$220/27$ &$11/40,\ 8/7$   &$6.73$\\
13&1&11&$16/39$&$24416185159/25705247659$        &$1040/243$&$13/80,\ 90/77$&$27.36$\\
13&2& 9&$19/39$&$482409391/524596891$            &$364/75$ &$13/56,\ 102/89$&$5.66$\\
13&3& 7&$1/2$  &$1341937/1529437$                &$416/9$  &$13/32,\ 8/7$   &$45.38$\\
15&1&13&$2/5$  &$94349947907/96889010407$        &$480/121$&$5/32,\ 34/29$  &$48.67$\\
15&2&11&$7/15$ &$1891389243/1977326743$          &$180/49$ &$5/24,\ 38/33$  &$6.89$\\
15&3& 9&$1/2$  &$37541107/40353607$              &$80/9$   &$5/16,\ 8/7$    &$7.48$\\
17&2&13&$23/51$&$1609027239419/1647113176919$    &$748/243$&$17/88,\ 126/109$&$9.78$\\
17&3&11&$1/2$  &$32325492131/33614554631$        &$1088/225$&$17/64,\ 8/7$  &$4.59$\\
17&4& 9&$1/2$  &$643823819/686011319$            &$170/3$  &$17/40,\ 8/7$   &$67.25$\\
19&4&11&$1/2$  &$36280145617/37569208117$        &$266/27$ &$19/56,\ 8/7$   &$8.68$\\
21&4&13&$1/2$  &$665527760349/678223072849$      &$126/25$ &$7/24,\ 8/7$    &$4.38$\\
21&5&11&$1/2$  &$13411599701/13841287201$        &$336/5$  &$7/16,\ 8/7$    &$93.93$\\
23&4&15&$1/2$  &$108095281916189/109193914728689$&$506/147$&$23/88,\ 8/7$   &$3.46$\\
23&5&13&$1/2$  &$2190361301861/2228447239361$    &$1472/135$&$23/64,\ 8/7$  &$10.27$\\
25&6&13&$1/2$  &$95365572907/96889010407$        &$700/9$  &$25/56,\ 8/7$   &$126.64$\\
\bottomrule
\end{tabular}
\end{table}

\subsubsection{Completeness}\label{sec:complete}
Groups (a), (b), (c) contain $22+6+21=49$ pairs; they are pairwise disjoint and their union is the full list of residual pairs with $m\le29$, as one checks by enumeration (step~B2-6).

\begin{proof}[Proof of Theorem~\ref{thm:crit}]
For $m\ge31$: Theorem~\ref{thm:tail}.  For $m\le29$: \S\ref{sec:grpA}--\S\ref{sec:grpC} cover all $49$ pairs by \S\ref{sec:complete}.
\end{proof}

\subsection{Assembly: proof of the unified strip theorem}\label{sec:assembly}

\begin{proof}[Proof of Theorem~\ref{thm:strip-main}]
Fix a residual pair $(m,r)$ and a strip point $(q,\ell)$: $q\in[0,\tfrac13]$, $\ell\in(0,\tfrac12]$, $\ell<q+\tfrac rm$.  Let $\lbar=\lbar_{m,r}$ be the evaluation level \eqref{eq:lbar}.

\emph{Route 1: $\ell\le\lbar$.}  By Corollary~\ref{cor:twovar} ($q\le\tfrac13$, $\ell\le\lbar$) and Theorem~\ref{thm:crit}(a),
$R(m,r,q,\ell)\ge R(m,r,\tfrac13,\lbar)\ge1$,
and Theorem~\ref{thm:leftsurplus} (hypotheses $q\in[0,\tfrac13]$, $0<\ell<q+\tfrac rm$: both hold at a strip point) gives $I\ge0$.

\emph{Route 2: $\ell>\lbar$.}  We first claim $\theta>\tfrac16$ and $(m,r)\notin\mathcal T$.  Indeed, in the opposite case $\lbar=\min(\tfrac12,\tfrac13+\tfrac rm)$, while every strip point satisfies $\ell\le\tfrac12$ and $\ell<q+\tfrac rm\le\tfrac13+\tfrac rm$, hence $\ell\le\lbar$ --- contradiction.  So $\lbar\in\{\tfrac{43}{100},\tfrac{39}{100}\}$, and:
\begin{itemize}
\item $\theta>\tfrac16>\tfrac18$;
\item $q\le\tfrac13<\tfrac{39}{100}\le\lbar<\ell$ (as $100<117$);
\item $\Lambda(q)\le\lbar<\ell$ by Theorem~\ref{thm:crit}(b);
\item $\ell\in(0,\tfrac12]$.
\end{itemize}
Thus $\ell\ge\max(\Lambda(q),q)$ and Theorem~\ref{thm:reflection} gives $I\ge0$.

In both routes $I_{m,r}(q,\ell)\ge0$, hence $\wtP_{m,r}(q,\ell)=C_{m,r}\ell^{\,n+r}I\ge0$.
\end{proof}

\begin{remark}[Structure of the case split]\label{rmk:structure}
The two routes overlap harmlessly and their union is the whole strip by the trivial dichotomy $\ell\le\lbar$ or $\ell>\lbar$; no covering lemma involving $\ell_{\mathrm{hand}}=\min(\max(\Lambda(q),q),\tfrac12)$ is needed (the stage A1/A2 handoff at $\ell_{\mathrm{hand}}$ is replaced by the handoff at the $q$-free level $\lbar$, which Theorem~\ref{thm:crit} makes admissible on both sides).  The two estimate-carrying band systems are internally disjoint: $(0,\varepsilon)$ vs.\ $(a,b)$ in Route~1, and $(a,b)$ vs.\ $(b,b+L)$ in Route~2; Route~1 uses no reflection and Route~2 uses no surplus, so the band-accounting pitfall of the July-4 $r=1$ proof cannot arise in either.
\end{remark}

\subsection{Validation}\label{sec:valid}

Margins (float; step~B1-4): over all $2500$ residual pairs with $5\le m\le201$, $\min\log R(m,r,\tfrac13,\lbar_{m,r})=0.1695$ at $(187,31)$ (a W1 pair with $\theta$ just below $\tfrac16$); over $203\le m\le1001$, the minimum is $0.1729$ at $(205,34)$.  At the finer stage-B1 evaluation level $\min(\max(\Lhat,\tfrac13),\tfrac12,\tfrac13+\theta)$ the corresponding minimum is $1.4810$, at $(19,3)$ --- the coarser $q$-free windows of \eqref{eq:lbar} trade margin for a fully analytic tail, and the criterion stays comfortably true either way.  As a sanity anchor, three handoff margins of the cartography are reproduced exactly: $R(9,2,\tfrac13,\tfrac12)\ge27$, $R(19,3,\tfrac13,\tfrac{470}{1007})\ge\tfrac{489}{100}$, $R(61,11,\tfrac13,\tfrac{107492}{301767})\ge14$ (log-margins $3.3095$, $1.5879$, $2.6532$).

\smallskip
\noindent\textbf{Validation script} (\texttt{validate\_strip\_final.py}, same directory; merges the four stage scripts; must exit $0$):
\begin{enumerate}[label=\textup{\arabic*.}]
\item[A1-1/2/3.] Exact sublemma checks ($2r(n-1)\le(2r+1)m$; $c_n$ comparisons); chain soundness $\underline\Sigma\le\Sigma$, $\overline D\ge D$, $\underline\Sigma/\overline D=R$, $I\ge\underline\Sigma-\overline D$ by 40-digit quadrature at twelve stress points; monotonicity of $R$ in $q,\ell$ by random finite differences.
\item[A2-1/2/3.] Side-condition equivalence on all pairs $m\le101$; payment product inequality \eqref{eq:product} on a float grid plus \emph{exact} rational verification at $\ell=\max(\Lambda(q),q)$ ($j=0$ equality at $\ell=\Lambda$, bracket inequality $j\le40$, product at rational $e$); quadrature $I\ge0$ at sample points of Route~2.
\item[B1-1/2/3/4.] M\"obius identities \eqref{eq:moebius} and rationalized cap exact at $10^4$ random rational points; $\Lhat-\max_q\Lambda\ge0$ float sweep (sharp at $(113,26)$); falsification values $\Lambda(0)=\tfrac4{17}$ for $(17,2)$, $\Lambda(\tfrac13)=\tfrac{54808}{136563}<\tfrac{67}{147}$ for $(49,6)$; the criterion scan and tail probe above.
\item[B2-1\dots7.] Every displayed constant of \S\ref{sec:tools}--\S\ref{sec:finite} in exact arithmetic (Pad\'e values, $\overline K$, $V_2$, $V_3$, final fractions, $c_n$ thresholds, \eqref{eq:193}, \eqref{eq:c15}, all group-(b) integer comparisons); minorant soundness on dense float grids; the $\Lambda$-cap lemma symbolically/numerically; tail simulation $31\le m\le501$ (window membership, hypotheses, $R\ge1$ float everywhere, exact on a subsample); all $21$ table rows recomputed exactly, the two reduced fractions digit-for-digit; completeness $49=22+6+21$; the three handoff anchors.
\item[M-1/2.] \emph{Merged assembly checks} (new): for every pair $m\le201$, $R(m,r,\tfrac13,\lbar_{m,r})\ge1$ (exact rational for $m\le61$, float beyond) and, for every $\theta>\tfrac16$ pair $\notin\mathcal T$, $\sup_q\Lambda(q)\le\lbar_{m,r}$ on a dense $q$-grid (exact at $q\in\{0,\tfrac16,\tfrac13\}$) --- i.e.\ both inputs of the Route-1/Route-2 dichotomy.
\end{enumerate}

\begin{remark}[Remaining computer-assisted residue; honest scope]\label{rmk:residue}
Everything above is analytic except one class of statements: the $21$ table rows are one-line exact rational inequalities whose large-integer cross-multiplications are verified by \texttt{validate\_strip\_final.py} rather than by hand (the remaining displayed integer comparisons --- \eqref{eq:pade}, \eqref{eq:193}, \eqref{eq:c15}, the group-(b) caps and window memberships, and the window constants --- are small enough to check manually).  This is the only computer-assisted residue in the strip proof.  The theorem covers exactly $q\in[0,\tfrac13]$; nothing here touches Region II ($q>\tfrac13$), where the diagonal-positivity route is known to fail structurally beyond $q\approx0.41175$.
\end{remark}



\section{The Appendix-B constants of [R1], analytically}\label{sec:appconst}

\subsection*{Provenance and verification status}
Proved by Claude on 2026-07-16; independent verifier: \emph{pass};
adversarial reviewer: \emph{sound}.  This is a drop-in replacement for
Appendix~B of [R1] (labels \emph{app:constants}, \emph{eq:app-caseA},
\emph{eq:app-caseB} preserved): both constant inequalities are proved for
\emph{all real} $\theta$ in their ranges and \emph{all real} $m>0$, so the
scripted sweep over odd $63\le m\le499$ and the twelve scripted integer
comparisons disappear.  The whole argument rests on the elementary bound
$B^m/m\ge e\log B$ plus rational Pad\'e minorants of the logarithm.  The
content is the note \texttt{claude\_jul17/app\_constants\_analytic.tex},
inlined verbatim.

\section{Constant inequalities for the residual-strip closure}\label{app:constants}

This appendix proves the two numerical inequalities used in
Lemma \emph{left-surplus-tail} of [R1].  Let
\[
        P(\theta)=\frac{2/3-2\theta}{\theta(1/2-2\theta)^2},
\]
\[
        B_0(\theta)=
        \left(\frac76\right)^{1-2\theta}
        (8-24\theta)^{-\theta}
        \frac{1/2+\theta}{5/12+\theta},
\]
and
\[
        B_1(\theta)=
        \left(\frac76\right)^{1-2\theta}
        (8-24\theta)^{-\theta}
        \frac{34}{29}.
\]
The required inequalities are
\begin{align}
        \frac{99}{100m}P(\theta)B_0(\theta)^m&\ge1
        &&(\theta=r/m\le1/6),\label{eq:app-caseA}\\
        \frac{99}{100m}P(\theta)B_1(\theta)^m&\ge1
        &&(1/6\le\theta=r/m<1/4),\label{eq:app-caseB}
\end{align}
for the integer pairs appearing in the residual strip with $m\ge63$ and
$r\ge2$.  We prove the following stronger statements, in which the parameters
are \emph{real} and unconstrained: \eqref{eq:app-caseA} holds for every real
$0<\theta\le1/6$ and every real $m>0$
(Proposition~\ref{prop:app-caseA-full}), and \eqref{eq:app-caseB} holds for
every real $1/6\le\theta<1/4$ and every real $m>0$
(Proposition~\ref{prop:app-caseB-full}).  In particular no finite
verification, and no use of the integrality of $r$ or of the parity or size
of $m$, remains in this appendix.

The elimination of the parameter $m$ rests on the elementary bound
$B^m/m\ge e\log B$ (Lemma~\ref{lem:app-tools}(ii) below): it replaces the
minimisation of $B^m/m$ over $m$ by its exact infimum over real $m>0$, which
is attained as $m\log B\to1$ and equals $e\log B$.  After this step both
inequalities become one-variable statements about $\theta$ alone, and these
are proved by exhibiting rational tangent/secant minorants of the logarithms
involved.

\subsection*{Elementary tools}

\begin{lemma}[Exponential and logarithm bounds]\label{lem:app-tools}
\begin{enumerate}[label=\textup{(\roman*)}]
\item $e^{y}\ge1+y$ for every real $y$.
\item For every real $B>0$ and every real $m>0$,
\[
        \frac{B^m}{m}\ \ge\ e\log B .
\]
\item $\log u\le u-1$ for every real $u>0$.
\item For every real $t\ge1$,
\[
        \frac{3(t^2-1)}{t^2+4t+1}
        \ \le\ \log t\ \le\
        \frac{(t-1)(t+5)}{2(2t+1)} .
\]
We write $L(t)$ and $U(t)$ for the left and right members.
\item $e\ge 65/24$.
\end{enumerate}
\end{lemma}

\begin{proof}
(i) The function $y\mapsto e^y-1-y$ has derivative $e^y-1$, which is
negative for $y<0$ and positive for $y>0$; its minimum value, at $y=0$, is
$0$.

(ii) Apply (i) with $y=m\log B-1$:
\[
        B^m=e^{m\log B}=e\cdot e^{\,m\log B-1}\ge e\,(1+m\log B-1)=e\,m\log B,
\]
and divide by $m>0$.

(iii) Apply (i) with $y=u-1$: $e^{u-1}\ge u$, and take logarithms (both
sides are positive).

(iv) Both bounds are equalities at $t=1$, so it suffices to compare
derivatives for $t\ge1$.  For the lower bound,
\[
        \frac{d}{dt}L(t)
        =\frac{6t(t^2+4t+1)-3(t^2-1)(2t+4)}{(t^2+4t+1)^2}
        =\frac{12(t^2+t+1)}{(t^2+4t+1)^2},
\]
and therefore
\[
        \frac{1}{t}-\frac{d}{dt}L(t)
        =\frac{(t^2+4t+1)^2-12t(t^2+t+1)}{t\,(t^2+4t+1)^2}
        =\frac{t^4-4t^3+6t^2-4t+1}{t\,(t^2+4t+1)^2}
        =\frac{(t-1)^4}{t\,(t^2+4t+1)^2}\ \ge\ 0 .
\]
For the upper bound, writing $U(t)=(t^2+4t-5)/(4t+2)$,
\[
        \frac{d}{dt}U(t)
        =\frac{(2t+4)(4t+2)-4(t^2+4t-5)}{(4t+2)^2}
        =\frac{4t^2+4t+28}{(4t+2)^2},
\]
and therefore
\[
        \frac{d}{dt}U(t)-\frac1t
        =\frac{4t^3+4t^2+28t-(4t+2)^2}{t\,(4t+2)^2}
        =\frac{4(t-1)^3}{t\,(4t+2)^2}\ \ge\ 0
        \qquad(t\ge1).
\]
Integrating the two derivative inequalities from $1$ to $t$ gives (iv).

(v) $e=\sum_{k\ge0}1/k!\ge1+1+\frac12+\frac16+\frac1{24}=\frac{65}{24}$,
since all terms of the series are positive.
\end{proof}

We record the six rational values of $L$ and $U$ that will be used; each is
obtained from Lemma~\ref{lem:app-tools}(iv) by clearing denominators.
\begin{equation}\label{eq:app-pade-values}
\begin{aligned}
        \log\frac{64}{63}&\ \ge\ L\!\left(\frac{64}{63}\right)
        =\frac{3\,(64^2-63^2)}{64^2+4\cdot64\cdot63+63^2}
        =\frac{3\cdot127}{24193},
        &
        \log\frac73&\ \ge\ L\!\left(\frac73\right)
        =\frac{3\,(49-9)}{49+84+9}
        =\frac{60}{71},
        \\[1mm]
        \log\frac{34}{29}&\ \ge\ L\!\left(\frac{34}{29}\right)
        =\frac{3\,(34^2-29^2)}{34^2+4\cdot34\cdot29+29^2}
        =\frac{945}{5941},
        &
        \log\frac76&\ \ge\ L\!\left(\frac76\right)
        =\frac{3\,(49-36)}{49+168+36}
        =\frac{39}{253},
        \\[1mm]
        \log 2&\ \le\ U(2)
        =\frac{1\cdot7}{2\cdot5}
        =\frac{7}{10},
        &
        \log\frac76&\ \le\ U\!\left(\frac76\right)
        =\frac{\frac16\cdot\frac{37}6}{2\cdot\frac{10}3}
        =\frac{37}{240}.
\end{aligned}
\end{equation}

\subsection*{Case A: $0<\theta\le1/6$}

Throughout this subsection set
\[
        \lambda_*=\frac{127}{24193},
        \qquad
        s=\frac{89}{100}.
\]

\begin{lemma}[Linear minorant of $\log B_0$]\label{lem:app-B0-linear}
For every $0\le\theta\le\frac16$,
\[
        \log B_0(\theta)\ \ge\ \lambda_*+s\left(\frac16-\theta\right).
\]
\end{lemma}

\begin{proof}
\emph{Value at the right endpoint.}
Since $4^{-1/6}=2^{-1/3}$,
\[
        B_0(1/6)
        =\left(\frac76\right)^{2/3}4^{-1/6}\,\frac{2/3}{7/12}
        =7^{2/3}6^{-2/3}\cdot2^{-1/3}\cdot\frac{8}{7}
        =8\,(7\cdot6^2\cdot2)^{-1/3}
        =\frac{8}{504^{1/3}}
        =\frac{4}{63^{1/3}},
\]
so that $B_0(1/6)^3=64/63$ and, by \eqref{eq:app-pade-values},
\begin{equation}\label{eq:app-lam-star}
        \log B_0(1/6)=\frac13\log\frac{64}{63}
        \ \ge\ \frac13\cdot\frac{3\cdot127}{24193}
        =\frac{127}{24193}=\lambda_* .
\end{equation}

\emph{Slope.}  From
\[
        \log B_0(\theta)
        =(1-2\theta)\log\frac76-\theta\log(8-24\theta)
        +\log\Bigl(\frac12+\theta\Bigr)-\log\Bigl(\frac5{12}+\theta\Bigr),
\]
the negated derivative $c(\theta):=-\dfrac{d}{d\theta}\log B_0(\theta)$ is
\begin{equation}\label{eq:app-c-def}
        c(\theta)
        =2\log\frac76+\log(8-24\theta)-\frac{24\theta}{8-24\theta}
        -\frac{1}{1/2+\theta}+\frac{1}{5/12+\theta},
\end{equation}
and differentiating once more,
\begin{equation}\label{eq:app-c-prime}
        c'(\theta)
        =-\frac{24}{8-24\theta}-\frac{192}{(8-24\theta)^2}
        +\frac{1}{(1/2+\theta)^2}-\frac{1}{(5/12+\theta)^2}.
\end{equation}
On $[0,1/6]$ we have $8-24\theta\ge4>0$, so the first two terms of
\eqref{eq:app-c-prime} are negative; and $0<\frac5{12}+\theta<\frac12+\theta$
gives $(1/2+\theta)^{-2}<(5/12+\theta)^{-2}$, so the sum of the last two
terms is negative as well.  Hence $c'<0$ on $[0,1/6]$, i.e.\ $c$ is strictly
decreasing there, and for every $\theta\in[0,1/6]$,
\begin{equation}\label{eq:app-c-lb}
        c(\theta)\ \ge\ c(1/6)
        =2\log\frac76+\log4-1-\frac32+\frac{12}7
        =2\log\frac73-\frac{11}{14}
        \ \ge\ \frac{120}{71}-\frac{11}{14}
        =\frac{899}{994}
        \ \ge\ \frac{89}{100}=s,
\end{equation}
where the value of $c(1/6)$ is read off from \eqref{eq:app-c-def} (using
$2\log\frac76+\log4=2\log\frac73$ and $-1-\frac32+\frac{12}7=-\frac{11}{14}$),
the penultimate inequality is \eqref{eq:app-pade-values} at $t=7/3$, and the
last comparison is $899\cdot100=89900\ge88466=89\cdot994$.

\emph{Conclusion.}  By the fundamental theorem of calculus, for
$0\le\theta\le1/6$,
\[
        \log B_0(\theta)
        =\log B_0(1/6)+\int_\theta^{1/6}c(u)\,du
        \ \ge\ \lambda_*+s\left(\frac16-\theta\right). \qedhere
\]
\end{proof}

\begin{lemma}[Weighted polynomial inequality]\label{lem:app-cubic}
For every $0<\theta\le\frac16$,
\[
        P(\theta)\left[\lambda_*+s\left(\frac16-\theta\right)\right]
        \ \ge\ 72\,\lambda_* .
\]
\end{lemma}

\begin{proof}
For $0<\theta<1/4$ one has $\theta(1/2-2\theta)^2>0$, so the claim is
equivalent to $G(\theta)\ge0$ on $(0,1/6]$, where
\[
        G(\theta)
        :=\left(\frac23-2\theta\right)
        \left[\lambda_*+s\left(\frac16-\theta\right)\right]
        -72\lambda_*\,\theta\left(\frac12-2\theta\right)^2 .
\]
$G$ is a cubic polynomial with rational coefficients.  At $\theta=1/6$ the
first term equals $\frac13\lambda_*$ and the second equals
$72\lambda_*\cdot\frac16\cdot\frac1{36}=\frac13\lambda_*$, so $G(1/6)=0$ and
$\left(\frac16-\theta\right)$ divides $G$.  Polynomial division gives the
exact factorisation
\begin{equation}\label{eq:app-G-factor}
        G(\theta)=\left(\frac16-\theta\right)Q(\theta),
        \qquad
        Q(\theta)
        =288\lambda_*\theta^2-(2s+96\lambda_*)\theta
        +\left(\frac{2s}3+4\lambda_*\right),
\end{equation}
which is verified by expanding both sides to
\[
        \left(\frac{2\lambda_*}3+\frac s9\right)
        -\left(s+20\lambda_*\right)\theta
        +\left(2s+144\lambda_*\right)\theta^2
        -288\lambda_*\theta^3 .
\]
On $[0,1/6]$,
\[
        Q'(\theta)=576\lambda_*\theta-(2s+96\lambda_*)
        \ \le\ 96\lambda_*-(2s+96\lambda_*)=-2s\ <\ 0,
\]
so $Q$ is strictly decreasing on $[0,1/6]$, and
\begin{equation}\label{eq:app-Q-endpoint}
        Q(\theta)\ \ge\ Q(1/6)
        =8\lambda_*-\frac{2s+96\lambda_*}{6}+\frac{2s}3+4\lambda_*
        =\frac s3-4\lambda_*
        =\frac{89}{300}-\frac{508}{24193}
        =\frac{2000777}{7257900}\ >\ 0 .
\end{equation}
Hence $Q>0$ on $[0,1/6]$ and $G\ge0$ on $(0,1/6]$ (with equality only at
$\theta=1/6$), which proves the lemma.
\end{proof}

\begin{proposition}[Case A, all real parameters]\label{prop:app-caseA-full}
For every real $0<\theta\le\frac16$ and every real $m>0$,
\[
        \frac{99}{100m}P(\theta)B_0(\theta)^m\ \ge\
        \frac{99}{100}\cdot72\cdot\frac{127}{24193}\cdot\frac{65}{24}
        =\frac{58841640}{58063200}\ >\ 1 .
\]
In particular \eqref{eq:app-caseA} holds for every integer pair of the
residual strip.
\end{proposition}

\begin{proof}
By Lemma~\ref{lem:app-B0-linear},
$\log B_0(\theta)\ge\lambda_*+s(\frac16-\theta)>0$.  By
Lemma~\ref{lem:app-tools}(ii) applied to $B=B_0(\theta)$,
\[
        \frac{B_0(\theta)^m}{m}\ \ge\ e\log B_0(\theta)
        \ \ge\ e\left[\lambda_*+s\left(\frac16-\theta\right)\right].
\]
Since $0<\theta\le1/6$ gives $\frac23-2\theta\ge\frac13>0$, we have
$P(\theta)>0$, and multiplying by $\frac{99}{100}P(\theta)$ preserves the
inequality:
\[
        \frac{99}{100m}P(\theta)B_0(\theta)^m
        \ \ge\ \frac{99e}{100}\,P(\theta)
        \left[\lambda_*+s\left(\frac16-\theta\right)\right]
        \ \ge\ \frac{99e}{100}\cdot72\,\lambda_*,
\]
the last step by Lemma~\ref{lem:app-cubic}.  Finally, by
Lemma~\ref{lem:app-tools}(v),
\[
        \frac{99e}{100}\cdot72\,\lambda_*
        \ \ge\ \frac{99}{100}\cdot72\cdot\frac{127}{24193}\cdot\frac{65}{24}
        =\frac{99\cdot72\cdot127\cdot65}{100\cdot24193\cdot24}
        =\frac{58841640}{58063200}\ >\ 1 . \qedhere
\]
\end{proof}

\subsection*{Case B: $1/6\le\theta<1/4$}

\begin{lemma}[Prefactor bound]\label{lem:app-P72}
For every $\frac16\le\theta<\frac14$, $\;P(\theta)\ge72$.
\end{lemma}

\begin{proof}
One checks by expansion the polynomial identity
\[
        \frac23-2\theta-72\,\theta\left(\frac12-2\theta\right)^2
        =-\frac23\,(6\theta-1)(72\theta^2-24\theta+1).
\]
On $[1/6,1/4]$ the factor $6\theta-1$ is nonnegative.  The quadratic
$72\theta^2-24\theta+1$ is convex, and its values at the endpoints are
\[
        72\cdot\frac1{36}-24\cdot\frac16+1=-1,
        \qquad
        72\cdot\frac1{16}-24\cdot\frac14+1=-\frac12,
\]
so by convexity it is at most $\max\{-1,-\frac12\}<0$ on all of $[1/6,1/4]$.
Hence the identity's right-hand side is nonnegative on $[1/6,1/4]$, i.e.
\[
        \frac23-2\theta\ \ge\ 72\,\theta\left(\frac12-2\theta\right)^2
        \qquad\left(\frac16\le\theta\le\frac14\right),
\]
and dividing by $\theta(1/2-2\theta)^2>0$ (legitimate for $\theta<1/4$)
gives $P(\theta)\ge72$.
\end{proof}

\begin{lemma}[Uniform lower bound for $\log B_1$]\label{lem:app-B1-quad}
For every $\frac16\le\theta\le\frac14$,
\[
        \log B_1(\theta)\ \ge\
        \frac{945}{5941}+\frac{39}{253}-\frac{4225}{13824}
        =\frac{157634591}{20778481152}
        \ \ge\ \frac{7}{1000}.
\]
\end{lemma}

\begin{proof}
Write $8-24\theta=4(2-6\theta)$, where $2-6\theta\in[\frac12,1]$ on the
stated range.  By Lemma~\ref{lem:app-tools}(iii) with $u=2-6\theta$,
\[
        \log(2-6\theta)\ \le\ 1-6\theta,
\]
hence, using $\theta>0$,
\[
        \log B_1(\theta)
        =\log\frac{34}{29}+(1-2\theta)\log\frac76
        -\theta\log4-\theta\log(2-6\theta)
        \ \ge\ C-K\theta+6\theta^2,
\]
where
\[
        C=\log\frac{34}{29}+\log\frac76,
        \qquad
        K=1+2\log2+2\log\frac76 .
\]
By \eqref{eq:app-pade-values},
\[
        C\ \ge\ \bar C:=\frac{945}{5941}+\frac{39}{253},
        \qquad
        K\ \le\ \bar K:=1+2\cdot\frac{7}{10}+2\cdot\frac{37}{240}
        =1+\frac75+\frac{37}{120}=\frac{325}{120}=\frac{65}{24}.
\]
Since $\theta>0$, replacing $C$ by $\bar C$ and $K$ by $\bar K$ can only
decrease $C-K\theta+6\theta^2$, and completing the square,
\[
        \log B_1(\theta)
        \ \ge\ \bar C-\bar K\theta+6\theta^2
        =\bar C-\frac{\bar K^2}{24}+6\left(\theta-\frac{\bar K}{12}\right)^2
        \ \ge\ \bar C-\frac{\bar K^2}{24}
        =\frac{945}{5941}+\frac{39}{253}-\frac{4225}{13824}.
\]
The exact value of the right-hand side is
$\frac{157634591}{20778481152}$ (the three denominators
$5941=13\cdot457$, $253=11\cdot23$ and $13824=2^9\cdot3^3$ are pairwise
coprime), and the final comparison
\[
        \frac{157634591}{20778481152}\ \ge\ \frac{7}{1000}
\]
is the integer inequality
$1000\cdot157634591=157634591000\ \ge\ 145449368064=7\cdot20778481152$.
\end{proof}

\begin{proposition}[Case B, all real parameters]\label{prop:app-caseB-full}
For every real $\frac16\le\theta<\frac14$ and every real $m>0$,
\[
        \frac{99}{100m}P(\theta)B_1(\theta)^m\ \ge\
        \frac{99}{100}\cdot72\cdot\frac{7}{1000}\cdot\frac{65}{24}
        =\frac{3243240}{2400000}\ >\ 1 .
\]
In particular \eqref{eq:app-caseB} holds for every integer pair of the
residual strip.
\end{proposition}

\begin{proof}
By Lemma~\ref{lem:app-B1-quad}, $\log B_1(\theta)\ge\frac7{1000}>0$.  By
Lemma~\ref{lem:app-tools}(ii) applied to $B=B_1(\theta)$, and then
Lemma~\ref{lem:app-P72} together with $\log B_1(\theta)>0$,
\[
        \frac{99}{100m}P(\theta)B_1(\theta)^m
        \ \ge\ \frac{99e}{100}\,P(\theta)\log B_1(\theta)
        \ \ge\ \frac{99e}{100}\cdot72\cdot\frac{7}{1000}
        \ \ge\ \frac{99}{100}\cdot72\cdot\frac{7}{1000}\cdot\frac{65}{24},
\]
the last step by Lemma~\ref{lem:app-tools}(v).  Finally
\[
        \frac{99\cdot72\cdot7\cdot65}{100\cdot1000\cdot24}
        =\frac{3243240}{2400000}\ >\ 1 . \qedhere
\]
\end{proof}

\begin{remark}[No computer assistance]\label{rmk:app-no-computer}
Propositions~\ref{prop:app-caseA-full} and~\ref{prop:app-caseB-full} prove
\eqref{eq:app-caseA} and \eqref{eq:app-caseB} for all real $\theta$ in the
stated ranges and all real $m>0$; the hypotheses $m\ge63$, $r\ge2$, $n>2r$
and $\theta=r/m$ of the residual strip are not needed here.  Consequently
the scripted verification of the finite range odd $63\le m\le499$ and the
twelve scripted integer comparisons of earlier versions of this appendix are
now superfluous, and the only computer-assisted ingredient remaining in the
proof of Theorem \emph{regionI-full} of [R1] is Proposition \emph{M61} of [R1]
(the residual strip for $m\le61$).  The margins of the two propositions are
$\frac{58841640}{58063200}-1\approx1.3\%$ (Case A; this is intrinsically
tight, since $\frac{99}{100m}P(\theta)B_0(\theta)^m\to\frac{99e}{100}\cdot
72\log B_0(1/6)\approx1.017$ along $\theta=1/6$, $m\log B_0(1/6)\to1$) and
$\frac{3243240}{2400000}-1\approx35\%$ (Case B).
\end{remark}
\section{The consolidated proof map and what remains}\label{sec:map}

\subsection*{The full theorem, post-cleanup}

For every graphon $W$ with edge density $p$ and every odd $m\ge3$,
$t(C_m,W)\ge p^m-p(1-p)^{m-1}$, by the following chain.
\begin{enumerate}[label=\textup{(\arabic*)},itemsep=2pt]
\item $p\le\tfrac12$: the right side is nonpositive.
\item If $\lambda_{\max}(A)\le q$: Lemma~\ref{lem:no-frontier}, at every
density, every odd $m\ge3$.
\item $\tfrac12<p<\tfrac23$ (Region II, frontier case): the one-frontier
reduction and Huber elimination of [R2] (valid for all odd $m\ge3$,
Remark~\emph{muniform} in \S\ref{sec:smallm}), then the scalar inequality
$R_m\le C_m\psi(\xi,\rho)$:
  \begin{itemize}[leftmargin=1.5em]
  \item odd $m\ge15$: Zones A (analytic, [R2]), B
  (\S\ref{sec:zoneB}, now analytic), C (\S\ref{sec:zoneC}, now analytic);
  \item odd $3\le m\le13$: \S\ref{sec:smallm} (uniform mechanism; two
  residual certificate steps, see below).  Alternatively the per-cycle
  theorems $C_5$--$C_{13}$ of [R1] remain valid as an independent route.
  \end{itemize}
\item $p\ge\tfrac23$ (Region I): the two-sided identity, edge deletion,
kernel expansion and mixture-of-diagonals of [R1]; then diagonal
positivity: pointwise regimes and integration-by-parts ([R1], analytic),
and the residual strip (including $r=1$) by \S\ref{sec:strip}, whose
constant inequalities are the analytic \S\ref{sec:appconst}.
\end{enumerate}

\subsection*{Remaining computer-assisted content}

\begin{center}
\small
\begin{tabular}{>{\raggedright\arraybackslash}p{0.30\linewidth}
                >{\raggedright\arraybackslash}p{0.30\linewidth}
                >{\raggedright\arraybackslash}p{0.32\linewidth}}
\toprule
Step & Status before this note & Status after this note \\
\midrule
(C1) $\operatorname{StripCert}(61)$, residual strip $m\le61$
& exact interval checker, $\sim$2400 pair-domains
& \emph{eliminated}: unified strip theorem for $r\ge1$, all odd $m$,
$q\le\tfrac13$; tail $m\ge31$ analytic, $49$ finite pairs with a
$21$-row displayed table (\S\ref{sec:strip}) \\
(C2) Appendix-B constants of [R1]
& scripted sweep odd $63\le m\le499$ + 12 integer comparisons
& \emph{eliminated}: analytic for all real $\theta$, $m>0$
(\S\ref{sec:appconst}) \\
(C3) Zone B battle
& 23-box certificate
& \emph{eliminated}: analytic; 15 displayed hand-checkable rows
(\S\ref{sec:zoneB}) \\
(C4) Zone C moderate-$e$
& 3001-box certificate, $m$ up to 1716
& \emph{eliminated}: analytic; 21 displayed hand-checkable rows
(\S\ref{sec:zoneC}) \\
(C5) $C_5$--$C_{13}$ per-cycle theorems (for $m\le13$ in Region II)
& long bespoke chapters with own certificates
& removed from the main chain by \S\ref{sec:smallm}; the small-$m$ route
retains two certificate steps: the Bernstein tables of
Prop.~\emph{band} and the moderate-range subdivision of
Prop.~\emph{moderate} (both exact rational, uniform protocol) \\
\bottomrule
\end{tabular}
\end{center}

Two audit conventions are inherited from [R1]/[R2] and kept here: a
displayed exact rational comparison (a finite table row a human can check
by hand, possibly with patience) is counted as analytic; an
interval-arithmetic subdivision run by a script is counted as
computer-assisted.

\subsection*{How to audit this note}

Each of \S\ref{sec:smallm}--\S\ref{sec:appconst} was produced with the
following pipeline: numeric cartography of the target inequality (dense
float grids over the exact domains, high-precision and exact rational
evaluation at extremes); an analytic proof draft; a machine validation of
every displayed inequality (companion scripts, see below); an
\emph{independent} verifier that re-derived every displayed claim from
scratch with its own script; and an adversarial reviewer instructed to
refute the logic (band disjointness, domain of monotonicity claims,
rounding directions, degenerate cases).  Defects found by this pipeline
(one unsafe rounding chain in an earlier version of
\S\ref{sec:zoneC}, five display-level constant errors in an earlier
version of \S\ref{sec:smallm}) were repaired and re-verified.  The
companion validation and verification scripts are collected in
\texttt{claude\_jul17/} in this directory, with a \texttt{README}
mapping each script to the section it certifies.

\subsection*{What genuinely remains open}

\begin{enumerate}[label=\textup{(\roman*)},itemsep=2pt]
\item A fully prose-analytic treatment of the two residual certificate
steps inside \S\ref{sec:smallm} (small-$m$ moderate range and Tur\'an
band).  The cartography shows the adapted zone arguments cover
$m\in\{7,\dots,13\}$ analytically; the honest obstruction is $m=3$ (and
partially $m=5$), where the payment-to-defect ratio approaches
$c'_3=0.9479$ at the corner --- any prose proof must be sharp to within
$5\%$ there.
\item The $21$ table rows of \S\ref{sec:strip} are one-line exact
rational inequalities in the sanctioned displayed style, but their
cross-multiplications involve large integers; a reader without a
computer can check them only with patience.  A prose tail argument
reaching $m\ge9$ (rather than $m\ge31$) would remove them.
\item The aesthetic endpoint: a proof with no density split at all.  By
Section~\ref{sec:dichotomy} this requires a payment mechanism exact on
the entire multipartite equality family while covering one-frontier
configurations --- a genuinely new idea, not an optimisation of the
present ones.
\end{enumerate}

\end{document}
